\documentclass[11pt]{article}
\usepackage{amsmath,amssymb,amsthm}
\usepackage[margin=1in]{geometry}
\newcommand{\R}{\mathbb R}
\newcommand{\E}{\mathbb E}
\newcommand{\GR}{G_{\R}}
\newcommand{\ip}[2]{\langle #1,#2\rangle}
\newcommand{\norm}[1]{\lVert #1\rVert}
\newcommand{\N}{\mathcal N}
\newcommand{\tr}{\operatorname{tr}}
\newtheorem{theorem}{Theorem}
\newtheorem{lemma}[theorem]{Lemma}
\newtheorem{proposition}[theorem]{Proposition}
\newtheorem{corollary}[theorem]{Corollary}
\allowdisplaybreaks
\setlength{\emergencystretch}{2em}
\title{Reduced matrix sequences and Gaussian cubic estimates\\for the real Grothendieck constant}
\author{}
\date{October 2, 2026}
\begin{document}
\maketitle
\begin{abstract}
We construct an explicit sequence of rational matrices whose Grothendieck
ratios converge to the reciprocal of a dimension-uniform Gaussian operator
norm. Independently checked whole-line inequalities and a complete continuous
source partition prove
\[
 \sup_d\|P_1-(711/1000)P_3\|_{L^\infty(\gamma_d)\to L^1(\gamma_d)}
 \le\frac{40}{71}.
\]
The resulting reduced matrix sequence proves $\GR\ge71/40=1.775$.
A self-contained signed-tensor rounding argument gives
$\GR<1760/987<1.785$, determining $\GR$ to three significant digits
as $1.78$. Exact ratio-preserving maps impose symmetry, zero diagonal,
zero row and column sums, and doubly stochastic entrywise absolute values.
The analytic cell entries require $O(d)$ arithmetic operations after
one-dimensional Gaussian moments are known. Fully rational finite-matrix
lower bounds converge to $1.775$ at rate $O(d^{-1})$; the actual matrix
ratio limit is the reciprocal of the true operator norm. An exhaustive
reduced sequence has monotone ratios converging to the full constant.
The norm proof uses coupled cubic-chaos moment inequalities, scalar Stein
certificates, and exact source geometry, with every numerical allowance
bound to its complete local proof. The final lower bound does not use the
unreconstructed historical coefficient-$0.699$ premise.
\end{abstract}

\section{Introduction and the numerical target}
For a nonzero finite real matrix $A$, define the bilinear norms
\[
 C(A)=\max_{s_i,t_j\in\{-1,1\}}\sum_{ij}A_{ij}s_it_j,
 \qquad
 V(A)=\max_{\norm{u_i}=\norm{v_j}=1}
                 \sum_{ij}A_{ij}\ip{u_i}{v_j}.
\]
The vector dimension is arbitrary; finite dimension $m+n$ suffices for an
$m$-by-$n$ matrix. Write $G(A)=V(A)/C(A)$ and
$\GR=\sup_A G(A)$.

The complete norm certificate in
Appendix~\ref{sec:new-final-goal} proves
\begin{equation}\label{eq:new-final-headline}
 \N(711/1000):=
 \sup_{d\geq1}\|P_1-(711/1000)P_3\|_{L^\infty(\gamma_d)\to L^1(\gamma_d)}
 \leq\frac{40}{71},\qquad
 \GR\geq\frac{71}{40}=1.775.
\end{equation}
Peng~\cite{peng} proves $\GR\leq17813/10000=1.7813$.
Thus the resulting enclosure is
\[
 1.775\leq\GR\leq1.7813,
\]
which determines the constant to three significant digits as $1.78$.
Proposition~\ref{prop:new-self-contained-upper} also proves the
self-contained upper bound $\GR<1760/987<1.785$, sufficient for
this rounding conclusion. The result is a three-significant-digit
approximation, rather than a determination to three decimal places.

The explicit Gaussian--Hermite matrix family has actual ratio limit
$1/\N(711/1000)\geq1.775$. Its scalar norm is not identified with the
certified cap, and this family is not assumed to attain the unrestricted
constant. Proposition~\ref{prop:universal-monotone} separately gives an
exhaustive reduced sequence whose ratios increase to $\GR$.
The logarithmic-cutoff construction and
Corollary~\ref{cor:new-rational-log-cutoff-bound} give finite rational
lower bounds and justify the limiting lower endpoint without numerical
extrapolation. Proposition~\ref{prop:new-compact-cubic-cell-entry}
expresses every analytic seed entry through one-dimensional truncated
Gaussian moments.

The norm proof combines the coupled cumulant theorem in Section~5,
the uniform cubic absolute-moment theorem in Section~6, the
support--source comparisons in Section~7, and the actual cubic projection
in Section~\ref{sec:r3}. All local remainder allowances, target-coefficient
spectral gates, and continuous source domains are checked independently.
The final closed moment refinement keeps every earlier valid cap and
exhausts all of its parent domains. The exact reduced-class maps are
proved in Section~2. At $b=1/2$, the published sharp operator inequality
of Saha et al.~\cite{saha} gives the exactly evaluable benchmark
$6\pi/11$ in Corollary~\ref{cor:half-benchmark}.

For historical comparison, the September 30 project estimate is retained
solely as the separate all-dimension premise
\begin{equation}\label{eq:baseline}
 \N(699/1000)\leq\frac{1000}{1773}.
\end{equation}
The earlier report attributes this estimate to a 55-interval amplitude
cover and inherited integrations~\cite{baseline}. That complete cover
has not been reconstructed here. Assuming \eqref{eq:baseline} gives
$\GR\geq1.773$; the fresh coefficient-$0.711$ theorem and its finite
matrix bounds have their own complete proof and do not use this premise.

\section{The exact $PAQ$ reduction}
The right transformation is denoted $R$ here to distinguish it from a
Gaussian cubic $Q$. Suppose distinct columns of $P$ have disjoint supports,
and distinct rows of $R$ have disjoint supports. Put
\[
 p_i=\sum_a|P_{ai}|,\qquad r_j=\sum_b|R_{jb}|.
\]
Then
\begin{equation}\label{eq:copies}
 C(PAR)=C(D_pAD_r),\qquad V(PAR)=V(D_pAD_r).
\end{equation}
In particular, if $p_i=\alpha>0$ for every original row and
$r_j=\beta>0$ for every original column, both norms are multiplied by
$\alpha\beta$, and $G(PAR)=G(A)$.

For scalar signs, the aggregate in row group $i$ lies in
$[-p_i,p_i]$. The maximum of the resulting bilinear form on the product
of these intervals occurs at a vertex, and every vertex is attainable
by aligned copied signs. For vectors the aggregate lies in the ball of
radius $p_i$; aligned copies attain its boundary. Ball vectors may be
completed to unit vectors using separate left and right orthogonal
coordinates without changing cross inner products. This proves
\eqref{eq:copies} for both norms.

Signed permutations and common scaling are special cases. Ordinary
diagonal balancing and arbitrary orthogonal changes of coordinates are
not consequences of this theorem. The requirements concern total
absolute copying weights, not the number of copies.

\subsection{A universal absolute doubly stochastic class}
Let $Z$ be a nonzero integer matrix with zero rows and columns removed.
Set
\[
 r_i=\sum_j|Z_{ij}|,\qquad c_j=\sum_i|Z_{ij}|,
 \qquad L=\sum_i r_i=\sum_jc_j.
\]
Make $r_i$ copies of row $i$, and $c_j$ copies of column $j$, and define
the square matrix of order $L$ by
\begin{equation}\label{eq:clone}
 \boxed{\quad B_{(i,a),(j,b)}=\frac{Z_{ij}}{r_ic_j},
 \quad 1\le a\le r_i,\quad1\le b\le c_j.\quad}
\end{equation}
Each original row or column has total copying weight one. Thus
\begin{equation}\label{eq:clone-norm}
 C(B)=C(Z),\qquad V(B)=V(Z),\qquad G(B)=G(Z).
\end{equation}
Every row and column of $|B|$ sums to one. For example, a cloned row
has sum $\sum_jc_j|Z_{ij}|/(r_ic_j)=1$.

For rational input first clear its denominator. For real input choose
rational approximations; both norms satisfy
\begin{equation}\label{eq:lipschitz}
 |C(A)-C(A')|,\ |V(A)-V(A')|
                 \le\sum_{ij}|A_{ij}-A'_{ij}|.
\end{equation}
Consequently the supremum defining $\GR$ is unchanged by restricting
to rational matrices whose entrywise absolute values are doubly stochastic,
over all orders. The clone order can be enormous, but both objectives
can be evaluated on the smaller seed $Z$.

\subsection{Symmetry, zero diagonal and zero sums}
For $B$ in \eqref{eq:clone}, set
\begin{equation}\label{eq:balanced}
 K(B)=\begin{pmatrix}0&B\\B^T&0\end{pmatrix},\qquad
 J_0=\frac12\begin{pmatrix}1&-1\\-1&1\end{pmatrix},\qquad
 H(B)=J_0\otimes K(B).
\end{equation}
Then $H(B)$ is rational, symmetric, zero-diagonal, has zero row and
column sums, and $|H(B)|$ is doubly stochastic. Its order is $4L$ and
\begin{equation}\label{eq:balanced-norm}
 C(H(B))=4C(B),\qquad V(H(B))=4V(B).
\end{equation}
The two off-diagonal terms in $K(B)$ use independent left/right variable
lists and can optimize simultaneously, giving factor two. The signed
two-copy map $J_0$ gives a further factor two by \eqref{eq:copies}.
Hence all these restrictions can be imposed simultaneously without
changing the supremum or any individual rational seed ratio.

\subsection{A monotone sequence converging to the full constant}
The analytic cubic sequence below has limit $1/\N(b)$.
The reduction also permits an unconditional sequence converging to the
full Grothendieck constant, without assuming that a cubic operator is
optimal.

\begin{proposition}\label{prop:universal-monotone}
Let $\mathcal A_n$ consist of all nonzero real matrices of dimensions
at most $n$ whose entries lie in
$2^{-n}\mathbb Z\cap[-1,1]$. Choose a matrix $S_n\in\mathcal A_n$
maximizing $G$, with a fixed lexicographic rule for ties. Clear its
dyadic denominator and apply \eqref{eq:clone} and \eqref{eq:balanced},
obtaining $\widetilde H_n$. Then every $\widetilde H_n$ belongs to the simultaneous reduced
class and
\[
 G(\widetilde H_n)\le G(\widetilde H_{n+1}),\qquad
 \lim_{n\to\infty}G(\widetilde H_n)=G_{\mathbb R}.
\]
\end{proposition}
\begin{proof}
Each candidate set is finite and the sets are nested. Both norms are
positive on a nonzero matrix, and the exact reductions preserve their
ratio, so the maxima exist and increase monotonically.
For any fixed nonzero real matrix $A$, homogeneity permits the
normalization $\max_{ij}|A_{ij}|\le1$. Its dyadic approximants $A^{(n)}$
belong to $\mathcal A_n$ for sufficiently large $n$ and have total
entry error tending to zero. The two norms satisfy
\[
 |C(A)-C(B)|\le\sum_{ij}|A_{ij}-B_{ij}|,
 \qquad |V(A)-V(B)|\le\sum_{ij}|A_{ij}-B_{ij}|.
\]
Consequently $G(A^{(n)})\to G(A)$ and
$\liminf_nG(\widetilde H_n)\ge G(A)$. Taking the supremum over $A$ proves the
asserted limit; the reverse inequality follows from the definition of
$G_{\mathbb R}$.
\end{proof}

This is a finite exact specification at every stage. The scalar norm is
computed by sign enumeration. The vector norm is the maximum of the
rational linear functional
$\sum_{ij}A_{ij}X_{i,m+j}$ on the compact set of positive semidefinite
$(m+n)$-by-$(m+n)$ matrices with diagonal one. Its polynomial
description makes exact algebraic optimization and comparison possible.
Enumeration is computationally expensive. The analytic subfamily gives
controlled candidates and a quantitative limit, while retaining a best
previous candidate prevents a weaker trial from replacing it.

\section{An explicit analytic seed and its limit}
Let $b\in[0,1)$ be computable (a rational coefficient suffices), and put
$T_b=P_1-bP_3$, where $P_j$ projects onto total
Hermite degree $j$ in Gaussian space. Start at $n\ge2$ and let $d_n=n^2$. In each coordinate
partition the line by the grid of mesh $n^{-2}$ from $-n$ to $n$, together
with its two exterior intervals. The grid has zero as an edge and is
antipodally symmetric. Form product cells and retain the cells whose
first coordinate lies on the negative half-line. These represent all
antipodal pairs.

For a normalized product Hermite $H_\xi$, define
\[
 F_{j,n}[C,\xi]=2\int_C H_\xi\,d\gamma_{d_n},\qquad |\xi|=j,
 \quad j\in\{1,3\}.
\]
The finite seed is
\begin{equation}\label{eq:seed}
 A_n(b)=F_{1,n}F_{1,n}^T-bF_{3,n}F_{3,n}^T.
\end{equation}
All entries are analytic and computable. With normalized one-variable
Hermites $h_k=\mathrm{He}_k/\sqrt{k!}$,
\[
 \int_a^b h_0\,d\gamma=\Phi(b)-\Phi(a),\qquad
 \int_a^b h_k\,d\gamma
 =\frac{\varphi(a)\mathrm{He}_{k-1}(a)
       -\varphi(b)\mathrm{He}_{k-1}(b)}{\sqrt{k!}}\quad(k\ge1),
\]
with boundary terms zero at infinite endpoints. Products give the
multivariate entries, so there is no unevaluated high-dimensional integral.

Extending cell signs oddly shows $C(A_n(b))\le\N(b)$, where
$\N(b)$ is the supremum of the scalar operator norm over all Gaussian
dimensions. Conversely any bounded odd functions in a fixed dimension
can be approximated by these partitions in their first coordinates as
$n\to\infty$. Conditional cell averaging converges in $L_2$; the
remaining coordinates integrate out. The partitions need not be nested:
continuous compactly supported functions are dense in $L_2$, and their
cell-average error tends to zero as the mesh shrinks and the cutoff grows.
Conditional averaging is an $L_2$ contraction, so density gives the same
conclusion for arbitrary $L_2$ functions. Therefore
\begin{equation}\label{eq:classical-limit}
 C(A_n(b))\longrightarrow\N(b).
\end{equation}
The supremum can be restricted to odd functions because $T_b$ annihilates
the even part. Maximization over bounded step functions reduces to signs
at the vertices of the finite cubes.

For $X\sim\gamma_d$, the radial vector field $X/\norm X$ gives
\begin{equation}\label{eq:radial}
 q_d(b)=a_d\left(1-\frac{b}{2(d+2)}\right),\qquad
 a_d=\frac{2\Gamma((d+1)/2)^2}{d\Gamma(d/2)^2}
 \longrightarrow1.
\end{equation}
Its total degree-one energy is $a_d$; its degree-three energy is
$a_d/[2(d+2)]$, by Gaussian radial integration. The ordinary $L_2$
norm of $T_b$ is one. Consequently $V(A_n(b))\le1$.

For an explicit approximation bound, let $r_n(X)$ be the vector of
cell midpoints, with exterior coordinates clipped to $\pm n$, and set
$v_n=r_n/\|r_n\|$. All its coordinates are nonzero because zero is an
edge of the grid. Thus this is an odd unit vector field, constant on cells.
With $d=n^2$, $h=n^{-2}$, and
\[
 p_n=\min\{1,2n\varphi(n)\},\qquad
 E_n=\sqrt{\frac1{n^2(n^2-2)}+8n\varphi(n)},
\]
the elementary normalization inequality and Gaussian Mills bound give
\begin{equation}\label{eq:direction-error}
 \|v_n-X/\|X\|\|_{L_2}^2
 \le dh^2\E\|X\|^{-2}+4\Pr(\max_i|X_i|>n)
 \le E_n^2.
\end{equation}
Here $\E\|X\|^{-2}=1/(d-2)$, and the union bound gives
$\Pr(\max_i|X_i|>n)\le2d\varphi(n)/n$.
In particular $E_n=n^{-2}+O(n^{-4})$.

The vector objective admits a sharper, quadratic error estimate.
Write $R=\|X\|$, $\mu=\E R$, $w=X/R$, and $e=v_n-w$.
The degree-one and degree-three projections above give
\[
 T_bw=t(R)w,\qquad
 t(R)=\frac{\mu R}{d}
 \left(1-\frac b2+\frac{bR^2}{2(d+2)}\right)\ge0.
\]
Pointwise unit length implies $e\cdot w=-\|e\|^2/2$.
Since the lowest odd spectral coefficient of $T_b$ is $-b$,
\[
 \langle w,T_bw\rangle-\langle v_n,T_bv_n\rangle
 =\E[t(R)\|e\|^2]-\langle e,T_be\rangle
 \le\E[(t(R)+b)\|e\|^2].
\]
Off the exterior event, $\|e\|^2\le dh^2/R^2$; on it,
$\|e\|^2\le4$. Using
$\mu\E R^{-1}=d a_d/(d-1)$, $a_d\le1$, and
\[
 \E[R\mathbf1_E]\le\sqrt{d\Pr(E)},\qquad
 \E[R^3\mathbf1_E]\le\sqrt{d(d+2)(d+4)\Pr(E)},
\]
proves the entirely explicit bound
\begin{align}\label{eq:quadratic-error}
 L_n(b)={}&\frac1{n^2}\left[
 \frac{1-b/2}{n^2-1}+\frac b{2(n^2+2)}+\frac b{n^2-2}\right]\notag\\
 &+4\left[1-\frac b2+\frac b2\sqrt{\frac{n^2+4}{n^2+2}}\right]
 \sqrt{p_n}+4bp_n.
\end{align}
Thus
\begin{equation}\label{eq:vector-bound}
 V(A_n(b))\ge q_{n^2}(b)-L_n(b)\longrightarrow1,
 \qquad L_n(b)=(1+b)n^{-4}+O(n^{-6}).
\end{equation}
The exponentially small terms in \eqref{eq:quadratic-error} are retained
in every finite bound. The first-order direction error cancels because
the radial field is pointwise aligned with its operator image.

\subsection{The sequence inside the reduced class}
Choose computable symmetric dyadic matrices $\widehat A_n(b)$ with
total entry error $\eta_n\le n^{-4}$. Directed entry intervals select
the first dyadic approximation whose certified error meets that budget;
this avoids deciding an exact rounding tie. Clear the denominator to
$Z_n(b)$, apply \eqref{eq:clone}, and then \eqref{eq:balanced}:
\begin{equation}\label{eq:final-sequence}
 \boxed{\quad H_n(b)=H(B(Z_n(b))).\quad}
\end{equation}
Every $H_n(b)$ lies in the simultaneous reduced class, and
\[
 G(H_n(b))=G(\widehat A_n(b))
 \ge\frac{q_{n^2}(b)-L_n(b)-\eta_n}{\N(b)+\eta_n}.
\]
Indeed $C(A_n(b))$ is bounded away from zero: a first-coordinate
hemisphere gives $(2/\pi)(1-b/6)>5/(3\pi)>1/4$.
Thus $n\ge2$ also ensures $\widehat A_n(b)\ne0$. Equations
\eqref{eq:lipschitz}--\eqref{eq:vector-bound} prove
\begin{equation}\label{eq:limit}
 \lim_{n\to\infty}G(H_n(b))=\frac1{\N(b)}.
\end{equation}
In particular \eqref{eq:baseline} proves
$\lim_nG(H_n(.699))\ge1.773$ assuming that premise.
The exact normalization preserves this endpoint.
The gamma recurrence gives
$a_{d+2}/a_d=(d+1)^2/[d(d+2)]$, and consequently
$a_d>1-1/(2d)$: along each parity the quotient
$a_d/(1-1/(2d))$ decreases to one.
Therefore any proved $\N(b)\le B$ gives the finite, computable lower bound
\begin{equation}\label{eq:finite-ratio}
 G(H_n(b))\ge
 \frac{(1-1/(2n^2))(1-b/[2(n^2+2)])-L_n(b)-n^{-4}}
 {B+n^{-4}}
 =\frac1B-\frac{1+b}{2Bn^2}+O(n^{-4}).
\end{equation}

% Primary input: Saha et al., New Lower and Upper Bounds for the
% Grothendieck Constant, arXiv:2608.11158v2, Sections 8--9.
% The external computer-assisted proof is invoked, not replayed here.
\begin{corollary}[An exactly evaluable Gaussian benchmark]
\label{cor:half-benchmark}
Let $T=P_1-\frac12P_3$ on standard Gaussian space and let $A_n$
be its antipodally folded product-cell compression, with dimension
$d=n^2$, mesh $n^{-2}$, cutoff $\pm n$, and $0$ a coordinate-cell
boundary.  Then, for every $n\geq2$,
\[
 C(A_n)=\frac{11}{6\pi},\qquad V(A_n)\leq1,
 \qquad \lim_{n\to\infty}G(A_n)=\frac{6\pi}{11}.
\]
Choose a symmetric rational approximation $\widehat A_n$ with total
entrywise error at most $\eta_n=n^{-4}$, and apply the exact cloning,
symmetric dilation and balancing reductions to obtain $H_n$.
Thus $H_n$ is symmetric, has zero diagonal and zero row sums, and
$|H_n|$ is doubly stochastic.  The same actual ratio limit holds:
\[
 \lim_{n\to\infty}G(H_n)=\frac{6\pi}{11}
 =1.713595992867159948\ldots.
\]
The rationalized ratios need not lie below their limiting value.
\end{corollary}
\begin{proof}
The all-odd-sign inequality proved by Saha et al.~\cite{saha},
Sections 8--9, is
\[
 2\langle P_1f,P_1g\rangle-\langle P_3f,P_3g\rangle
 \leq\frac{11}{6}\frac2\pi.
\]
Replacing $g$ by $-g$ gives the corresponding absolute bound.
It also holds for odd functions of absolute value at most one:
successively maximize the bilinear pairing in each function, choosing
the odd sign of the other function's image under $T$.
For $s(X)=\operatorname{sgn}(X_1)$ and
$\nu=\sqrt{2/\pi}$, the raw Hermite polynomial
$\operatorname{He}_3(z)=z^3-3z$ gives
\[
 P_1s=\nu X_1,\qquad P_3s=-\frac\nu6
 \operatorname{He}_3(X_1),\qquad
 \langle s,Ts\rangle=\nu^2-\frac{\nu^2}{12}
 =\frac{11}{6\pi}.
\]
The hemisphere sign is exactly constant on each product cell.
Consequently its folded matrix assignment attains the external upper
bound, establishing the exact scalar norm for every finite $n$.
For any pair of unit-vector assignments, their odd cell fields $v,w$
have $\int\|v\|^2\,d\gamma_d=\int\|w\|^2\,d\gamma_d=1$;
the $L^2$ operator norm $\|T\|_{2\to2}=1$ and Cauchy--Schwarz give
$V(A_n)\leq1$.
The radial unit field and the established midpoint estimate give
$V(A_n)\to1$.
Finally, total entrywise error bounds both norm errors by $\eta_n$,
and every subsequent reduction preserves $G$ exactly.  This proves
both asserted limits.
\end{proof}

For completeness, the following are explicit finite lower bounds.
Put $d=n^2$, $p_n=2n\varphi(n)$, and
\begin{align*}
 e_n={}&\frac1d\left(\frac{3}{4(d-1)}+
             \frac1{4(d+2)}+\frac1{2(d-2)}\right)\\
 &+4\left(\frac34+\frac14\sqrt{\frac{d+4}{d+2}}\right)
       \sqrt{p_n}+2p_n,\\
 q_n={}&\left(1-\frac1{2d}\right)
             \left(1-\frac1{4(d+2)}\right)-e_n.
\end{align*}
The gamma-ratio lower bound $a_d>1-1/(2d)$ and the quadratic
midpoint loss imply
\[
 G(A_n)\geq\frac{q_n}{11/(6\pi)},\qquad
 G(H_n)\geq\frac{q_n-n^{-4}}{11/(6\pi)+n^{-4}}.
\]
Directed MPFR256 evaluation, rounded down to twelve decimal places,
gives
\[
\begin{array}{c|cc}
 n & G(A_n)\text{ lower bound} & G(H_n)\text{ lower bound}\\\hline
 32&1.712339483749&1.712335051217\\
 128&1.713517545081&1.713517527758\\
 1024&1.713594767206&1.713594767202
\end{array}
\]
These evaluations are recorded in
\texttt{saha\_benchmark\_certificate.json}; the external scalar
theorem is the sole new norm input and has not been independently
replayed here.  This benchmark does not invoke the retained, unverified
global norm premise used for the earlier $1.773$ candidate.

The same external theorem determines the complete analytic range
$0\le b\le1/2$. Indeed
$T_b=2bT_{1/2}+(1-2b)P_1$, and $\N(0)=2/\pi$.
Convexity and the same hemisphere equality give
\[
 \N(b)=\frac2\pi(1-b/6),\qquad
 \lim_nG(H_n(b))=\frac{\pi}{2(1-b/6)}
 \quad(0\le b\le1/2).
\]
This exactly soluble range ends at $6\pi/11$; improving it requires
$b>1/2$ or a different operator family.

\section{Coefficient choices and construction ceilings}
The final certified sequence uses $b=711/1000$, with the complete norm
cap $40/71$ in Theorem~\ref{thm:new-final-goal-global-norm}, proving the
lower endpoint $71/40=1.775$. The following construction ceilings describe
earlier coefficient choices and their reflected sign witnesses.

The initial candidate uses $b=18/25=.72$ in
\eqref{eq:seed}--\eqref{eq:final-sequence}. For this coefficient,
the target requires the all-dimension inequality
\begin{equation}\label{eq:target}
 \boxed{\quad\N(18/25)\le\frac{40}{71}.
 \quad}\qquad \frac1{40/71}=1.775.
\end{equation}
The previous $.699$ operator cannot reach this target: its retained
directed diffuse-cubic sign witness limits its reciprocal-norm construction
to $1.77440101857\ldots$ \cite{history}. Changing the coefficient is
therefore necessary within the pure cubic family.

The $.72$ coefficient has a fresh directed lower witness. Let
$h_3(x)=(x^3-3x)/\sqrt6$, let all Gaussians below be independent, and put
\[
 \begin{gathered}
 Y_N=N^{-1/2}\sum_{j=1}^Nh_3(X_j),\qquad a=511/200,\quad c=393/1000,\\
 f_N=\operatorname{sgn}(Y_N+aZ+cZ^3),\qquad
 g_N=\operatorname{sgn}(-Y_N+aZ+cZ^3).
 \end{gathered}
\]
The spectral central limit calculation assigns ordinary $Z$-degree $j$
and diffuse-cubic degree $k$ to total original Hermite degree $j+3k$.
Its full degree accounting is proved in Appendix~\ref{app:witness}. Writing
$q(z)=az+cz^3$, the limiting moments are
\[
 m_Y=2\int_{\R}\varphi(z)\varphi(q(z))\,dz,\quad
 m_1=2\int_{\R}\varphi(z)q'(z)\varphi(q(z))\,dz,
\]
\[
 m_3=\frac{2}{\sqrt6}\int_{\R}
              (z^2-1)\varphi(z)q'(z)\varphi(q(z))\,dz.
\]
The scalar pairing is $m_1^2-b(m_3^2-m_Y^2)$. Directed Taylor
integration and explicit tails give
\begin{equation}\label{eq:new-witness}
 \N(.72)>.563241882518393914882,
 \qquad
 \frac1{\N(.72)}<1.7754361511.
\end{equation}
This is a ceiling for the construction, not an achieved lower bound.
The witness lies below $40/71$ by more than $.00013839917$, and hence
does not rule out \eqref{eq:target}. Old whole-frame profile transfers
fail to certify \eqref{eq:target}; their failure is a limitation of
those sufficient estimates, not an obstruction from the true norm.

\begin{proposition}[Exact reflected norm witness at $b=7171/10000$]
Let $\N(b)=\sup_d\|P_1-bP_3\|_{L^\infty(\gamma_d)\to L^1(\gamma_d)}$.
Then
\[
 \N\left(\frac{7171}{10000}\right)
 >\frac{56323352479769}{10^{14}}.
\]
Consequently the asymptotically unit-numerator reciprocal-norm construction
for this fixed coefficient has ceiling strictly below
$1775462496412/10^{12}=1.775462496412$.
This is a construction ceiling, not a Grothendieck lower endpoint.
\end{proposition}
\begin{proof}
Use the normalized Hermite $h_3(x)=(x^3-3x)/\sqrt6$. Set $a=26412631/10^7$, $c=4084047/10^7$, $q(z)=az+cz^3$ and
$Y_n=n^{-1/2}\sum_{j=1}^nh_3(X_j)$ for independent standard Gaussian variables.
Use the actual original-space signs
$f_n=\operatorname{sgn}(Y_n+q(Z))$ and
$g_n=\operatorname{sgn}(-Y_n+q(Z))$.
The original-degree projection accounting for the reflected sign witness
applies because $q'(z)\ge a>0$; in particular the diffuse channel retains
original degree three.  Thus the limiting first- and third-chaos pairings are
$m_1^2$ and $m_3^2-m_Y^2$, respectively, where
\[
 m_Y=\frac2\pi\int_0^\infty e^{-(z^2+q(z)^2)/2}\,dz,
\quad
 m_1=\frac2\pi\int_0^\infty q'(z)e^{-(z^2+q(z)^2)/2}\,dz,
\]
\[
 m_3=\frac2{\pi\sqrt6}\int_0^\infty(z^2-1)q'(z)
                       e^{-(z^2+q(z)^2)/2}\,dz.
\]
The exact-rational directed certificate
\texttt{certify\_\allowbreak cubic7171\_\allowbreak norm\_\allowbreak witness.py}
integrates Taylor terms through degree eight on 320 rational cells of $[0,5]$,
bounds the degree-ten remainder over each entire cell, and bounds the omitted
Gaussian tails explicitly.  At 384-bit MPFR precision it encloses
\[
 m_1^2-\frac{7171}{10000}(m_3^2-m_Y^2)
 \in[0.5632335247976901337793,\,0.5632335247976901344356].
\]
Every finite-dimensional sign pairing is bounded above by $\N(7171/10000)$;
taking the limit proves the assertion.  The reciprocal enclosure has an upper
endpoint below $1.77546249641158$, proving the stated
construction ceiling.  The witness remains below $40/71$ and therefore
does not exclude a future certified lower endpoint of $1.775$.
A separate Arb replay at 512 bits, Taylor order twelve and 400 rational cells
recomputed all three integrals and Gaussian tails without importing the
producer, independently accepting the stated inequalities.
\end{proof}

A fresh reflected-sign witness also excludes $b=141/200=.705$ from
reaching the target within this reciprocal-norm family. Use the exact head
$q(z)=(31091537/10^7)z+(49105809/10^8)z^3$ in the same construction.
Directed integration gives
\[
 \N(.705)>\frac{563383848679}{10^{12}}>\frac{40}{71},
 \qquad \frac1{\N(.705)}<1.774988761828<1.775.
\]
An independent 512-bit Arb replay, with Taylor order twelve and 400
rational cells, checks all three moments and the omitted Gaussian tails.
The same replay at $b=.713$ gives a construction ceiling below
$1.775407850864$, leaving the target possible at the base coefficient $b=.713$. These are ceilings for the specified matrix
limits, and impose no upper bound on $G_{\mathbb R}$.

\section{A coupled cumulant inequality for real cubic chaos}\label{sec:cumulants}
Let $Q=I_3(f)$ be a real third Gaussian chaos, normalized by
$\E Q^2=6\|f\|_F^2=1$, where $f$ is a real fully symmetric
third-order tensor. The flattening Gram matrices below are symmetric
matrices; they need not be fully symmetric as four-index tensors.
The norm $\|\cdot\|_F$ is the Euclidean norm of all entries. Define
\[
 \kappa=\E Q^4-3,\qquad \kappa_6=\E Q^6-15\E Q^4+30,
 \qquad U=\tfrac13\|DQ\|^2-1,
 \qquad K=194/\sqrt{10}.
\]
\begin{theorem}\label{thm:633}
In every finite Gaussian dimension,
\begin{equation}\label{eq:deficit}
 \kappa_6\le K\kappa^{3/2}
 -633\sqrt\kappa\left(\E U^2-\frac{7\kappa}{45}\right).
\end{equation}
Also $0<\kappa\le90$ and
$7\kappa/45\le\E U^2\le2\kappa/9$. Equality holds in
\eqref{eq:deficit} for $Q=(Z^3-3Z)/\sqrt6$.
\end{theorem}

The scalar endpoint $\kappa=0$ is included by continuity in the later
price calculations. A nonzero normalized cubic tensor in finite dimension
has $\kappa>0$.

Write $F_{ab,c}=f_{abc}$, $A=FF^T$, $G=F^TF$, and let $S$ be the full
symmetrization of the four-index tensor $A$. Put
\[
 s=\|A\|_F^2=\|G\|_F^2,\quad h=\|S\|_F^2,
 \quad V=\tr A^3,\quad V_4=\tr A^4.
\]
Symmetrization is an orthogonal projection, so
$\langle A,S\rangle=h$, $0\le h\le s$, and
\begin{equation}\label{eq:snorm}
 \|S+tA\|_F^2=t^2s+(1+2t)h \qquad(t\in\R).
\end{equation}
The Gaussian product formula gives
\begin{equation}\label{eq:products}
 U=3I_4(S)+12I_2(G),\quad
 \E U^2=216h+288s,\quad \kappa=1296s+1944h.
\end{equation}
For example, the first identity follows by applying
$D_iI_3(f)=3I_2(f_i)$ and
\[
 I_2(g)^2=I_4(\operatorname{Sym}(g\otimes g))
 +4I_2(g\otimes_1 g)+2\|g\|_F^2
\]
to each tensor slice. Orthogonality of Hermite degrees gives the variance.
Since $\tr A=1/6$, $s\le1/36$, which also proves $\kappa\le90$.

The sixth-cumulant identity is
\begin{equation}\label{eq:wick6}
 \frac{\kappa_6}{46656}
 =\frac{135}{2}\tr S^3+45\tr(AS^2)+\frac{165}{2}\tr(A^2S)
 -20\tr A^3+\frac{135}{2}\langle S,G\otimes G\rangle.
\end{equation}
It agrees with the real cubic identity in \cite[Eq.~(8.12)]{HJM}.
For completeness, the accompanying exact program derives it from connected
Wick diagrams. For $r$ copies of $f$, a loopless multigraph with
multiplicities $m_{ij}$ and degree three at each vertex has weight
$(3!)^r/\prod_{i<j}m_{ij}!$. Cumulants retain precisely the connected
graphs, by the exponential formula for the logarithm of the Wick generating
function. There are seven labelled connected graphs in the fourth cumulant
and 640 in the sixth. Expanding each term of \eqref{eq:wick6}, including
the three pairings in $S$, gives the same rational coefficient in each of
the six sixth-order graph classes. This checks a formal tensor identity;
finite examples of tensors are not substituted for that identity.

Set
\[
 y=\sqrt{h/s},\qquad a=(V/s^{3/2})^{1/3},\qquad d=1+\tfrac32y^2.
\]
Then $0\le y,a\le1$ and
\begin{equation}\label{eq:variance-shape}
 \kappa=1296sd,\qquad \frac{\E U^2}{\kappa}
 =\frac19\left(1+\frac1d\right).
\end{equation}
The stated variance interval follows immediately from $1\le d\le5/2$.

For any of the four pairs
\[
 (p,q)=(17/33,-1),\ (62/105,-3/4),\ (28/39,-1/2),\ (3/4,-1),
\]
put
\[
 c_p=\frac{55/2-15q}{p(p-q)},\quad
 c_q=\frac{55/2-15p}{q(q-p)},\quad
 \zeta=135/2-c_p-c_q,\quad \beta=20+c_pp^3+c_qq^3.
\]
These pairs have $c_p,c_q>0$ and $\zeta\ge0$. Direct expansion completes
the first four terms on the right of \eqref{eq:wick6} as
\[
 c_p\tr(S+pA)^3+c_q\tr(S+qA)^3+\zeta\tr S^3-\beta V.
\]
For a real symmetric matrix $B$, $\tr B^3\le\|B\|_F^3$.
Also
\[
 \|\operatorname{Sym}(G\otimes G)\|_F^2=(s^2+2V_4)/3,
 \qquad V_4\le V^{4/3}.
\]
Together with \eqref{eq:snorm}, these inequalities prove
$\kappa_6/\kappa^{3/2}\le\mathcal J_{p,q}(y,a)/d^{3/2}$, where
\begin{align}\label{eq:J}
 \mathcal J_{p,q}={}&c_p[p^2+(1+2p)y^2]^{3/2}
 +c_q[q^2+(1+2q)y^2]^{3/2}\notag\\
 &+\zeta y^3-\beta a^3+\frac{135}{2}y\sqrt{(1+2a^4)/3}.
\end{align}
An additional valid bound is $\mathcal M(y,a)/d^{3/2}$, with
\begin{align}\label{eq:M}
 \mathcal M={}&\frac{135}{2}y^3+\frac{135}{2}y\sqrt{(1+2a^4)/3}
 +45ay^2-20a^3\notag\\
 &+\frac{165}{2}\min\left\{a^{3/2}y,
 y^2a^3+ya^2\sqrt{(1-y^2)(1-a^2)}\right\}.
\end{align}
For the first mixed bound use
$\tr(A^2S)\le\sqrt{hV_4}$ and $V_4\le\sqrt{s}V$.
For the second, write $R=S-(h/s)A$ and apply Cauchy--Schwarz to
$\langle A^2-(V/s)A,R\rangle$, then use $V_4\le V^{4/3}$.
Finally $\tr(AS^2)\le V^{1/3}h$ since
$\|A\|_{\mathrm{op}}\le V^{1/3}$. No positivity of $S$ is assumed.

The exact scalar comparison, supplied in Appendix~\ref{app:tensor}, is
\begin{equation}\label{eq:633scalar}
 \frac{\min\{\mathcal M,\mathcal J_{17/33,-1},
 \mathcal J_{62/105,-3/4},\mathcal J_{28/39,-1/2},
 \mathcal J_{3/4,-1}\}}{d^{3/2}}+\frac{211}{3d}
 \le\frac{194}{\sqrt{10}}+\frac{422}{15}
 \quad(0\le y,a\le1).
\end{equation}
Its certificate consists of 84 rational rectangles and 2100 exact
Bernstein coefficient checks in $\mathbb Q(\sqrt{10})$.
Combining \eqref{eq:variance-shape}--\eqref{eq:633scalar} with
\[
 633\left\{\frac19(1+1/d)-\frac7{45}\right\}
 =\frac{211}{3d}-\frac{422}{15}
\]
proves Theorem~\ref{thm:633}. At $Q=(Z^3-3Z)/\sqrt6$,
$y=a=1$, $\kappa=90$, and $\kappa_6=52380=K90^{3/2}$.

\section{A uniform cubic absolute-moment bound}\label{sec:stein}
\begin{theorem}\label{thm:l1}
For every normalized real third Gaussian chaos $Q$ in every finite
dimension,
\begin{equation}\label{eq:rho}
 \E|Q|<\rho:=849/1000,\qquad
 \|P_3\|_{L_\infty\to L_2}^2\le720801/1000000.
\end{equation}
\end{theorem}

Put $c=\sqrt{2/\pi}$ and
$R(x)=e^{x^2/2}\int_x^\infty e^{-t^2/2}\,dt$ for $x\ge0$.
The odd extension of $g(x)=-1+cR(x)$ solves
$g'(x)-xg(x)=|x|-c$, with
\[
 g'(x)=-c(1-|x|R(|x|)).
\]
The cubic Gaussian Stein identity is
\begin{equation}\label{eq:stein}
 \E[Qv(Q)]=\E[(1+U)v'(Q)],
 \qquad U=\tfrac13\|DQ\|^2-1.
\end{equation}
It follows directly from Gaussian integration by parts and the
Ornstein--Uhlenbeck eigenvalue of $Q$. Consequently, for every odd
continuously differentiable regression $\psi$ of polynomial growth and
every $\eta>0$,
\begin{equation}\label{eq:stein-young}
 \E|Q|\le c+\E\left[\psi'(Q)-Q\psi(Q)
 +\frac{(g'(Q)-\psi'(Q))^2}{2\eta}\right]
 +\frac\eta2\E U^2.
\end{equation}
Indeed $\E|Q|-c=-\E[Ug'(Q)]$ and
$\E[\psi'(Q)-Q\psi(Q)]=-\E[U\psi'(Q)]$; Young's inequality
applied to $-U(g'-\psi')$ gives \eqref{eq:stein-young}.

The eight rational rows in Appendix~\ref{app:stein} supply
$\eta>0$, coefficients $q_0,q_2,q_4,q_6$ with $q_6>0$, and a
regression $\psi$, such that on the entire real line
\begin{equation}\label{eq:stein-majorant}
 \psi'(x)-x\psi(x)+\frac{(g'(x)-\psi'(x))^2}{2\eta}
 \le q_0+q_2x^2+q_4x^4+q_6x^6.
\end{equation}
For $x\ge0$, each regression has the explicit rational form
\begin{equation}\label{eq:regression}
 \psi(x)=v_0x+\frac{v_1x^2}{2}+\frac{v_2x^3}{3}
 +\frac{v_3x^4}{4}
 +\frac{x}{x+2}\sum_{j=0}^{12}v_{4+j}
 T_j\left(\frac{x-2}{x+2}\right),
\end{equation}
extended oddly. It is continuously differentiable at zero.
All $v_j$, all $q_j$, and $\eta$ are rational and are supplied in the
source archive as exact inputs, separately from floating discovery output.

Combining \eqref{eq:stein-young}--\eqref{eq:stein-majorant} with
Theorem~\ref{thm:633} gives a price affine in $\E U^2$.
Its maximum therefore occurs at one of the two permitted variance
endpoints. The resulting functions of $\kappa$ are
\begin{align}\label{eq:cornerprices}
 p_0(\kappa)={}&c+q_0+q_2+(3+\kappa)q_4
 +[15+15\kappa+K\kappa^{3/2}]q_6+7\eta\kappa/90,\notag\\
 p_1(\kappa)={}&c+q_0+q_2+(3+\kappa)q_4
 +[15+15\kappa+(K-633/15)\kappa^{3/2}]q_6+\eta\kappa/9.
\end{align}
Both are convex for $\kappa\ge0$ because $q_6>0$ and $K-633/15>0$.
On a closed cumulant interval, four corner prices thus suffice.
The eight intervals in Appendix~\ref{app:stein} cover $[0,90]$,
and every directed corner enclosure is below $.848521<.849$.
This proves the first assertion of Theorem~\ref{thm:l1}.

For the second assertion, normalize $P_3f$ to a variance-one cubic $Q$
when that projection is nonzero. For $|f|\le1$,
$\|P_3f\|_2=\E fQ\le\E|Q|<\rho$. Taking the supremum and squaring
proves the projection norm bound.

% Fragment for the main paper; requires its theorem/proof environments.
% Self-contained proof fragment for insertion after the sharp633/Stein lemmas.
\begin{theorem}[A refitted uniform cubic absolute-moment bound]\label{thm:l1-repriced}
For every normalized real third Gaussian chaos $Q$ in every finite dimension,
\[
 \mathbb E|Q|<\frac{847}{1000},\qquad
 \|P_3\|_{L_\infty\to L_2}^2\le\frac{717409}{10^6}.
\]
\end{theorem}
\begin{proof}
Parameterize the coupled shape by
$\mathbb E U^2=\kappa(7/45+s/15)$, where
$\kappa=\mathbb E Q^4-3\in[0,90]$ and $s\in[0,1]$.
For any complete-halfline scalar row with frozen rational parameters
$(q_0,q_2,q_4,q_6,\eta)$, the Stein identity and sharp633 theorem imply
\[
\mathbb E|Q|\le c+q_0+q_2+(3+\kappa)q_4
+[15+15\kappa+(K-633s/15)\kappa^{3/2}]q_6
+\eta\kappa(7/90+s/30),
\]
where $c=\sqrt{2/\pi}$ and $K=194/\sqrt{10}$.
Eight accepted scalar rows are supplemented by two fresh rational Chebyshev rows.
Their entire-halfline inequalities are verified by directed compact Taylor
and exact polynomial-tail comparisons; a separate direct-Chebyshev
implementation replays the new rows with 768-bit sixth-order Taylor bounds
on $[0,80]$ and independently reconstructed Fraction tails.
Each row has $q_6>0$ and $\eta>0$, so its price is convex in $\kappa$
and affine in $s$ on every allowed rectangle. Its validity is global in
cumulant, independent of the interval where it was discovered.
The exact certificate partitions $[0,90]$ into 2780 rational intervals,
selects one of the ten rows on each, and verifies all four corner prices.
The largest outward rational cap is
$423485406457/500000000000=.846970812914<.847$.
An independent MPFR implementation reconstructs all 11120 selected
prices and 111200 ten-row candidate checks, verifies source hashes and
exact exhaustive coverage, and agrees with the bound. This proves the
first assertion using only the freshly included sharp633 and Stein inputs.
Normalizing a nonzero cubic projection proves the second by duality.
\end{proof}

\begin{corollary}\label{cor:low}
Let $f,g$ be signs, $h=(f+g)/2$, $k=(f-g)/2$, and
$u=\|P_1h\|_2$. For $T_{.72}=P_1-(18/25)P_3$, if $u\le21/100$, then
\begin{equation}\label{eq:lowcap}
 \langle f,T_{.72}g\rangle\le.56307672<40/71.
\end{equation}
\end{corollary}
\begin{proof}
Self-adjointness and orthogonality give
\[
 \langle f,T_{.72}g\rangle
 =\|P_1h\|_2^2-.72\|P_3h\|_2^2
 -\|P_1k\|_2^2+.72\|P_3k\|_2^2
 \le u^2+.72\rho^2.
\]
Since $|k|\le1$, Theorem~\ref{thm:l1} applies. At $u=.21$ the exact
right side is $441/10000+6487209/12500000=.56307672$.
\end{proof}
The amplitude is the actual first-projection norm, without division by
$\sqrt{2/\pi}$. The corollary covers a complete interval of amplitudes,
but does not cover the middle range needed for \eqref{eq:target}.

With the refitted cap $847/1000$, the same elementary energy proof gives
$\langle f,T_{.72}g\rangle\le u^2+(18/25)(847/1000)^2$.
It therefore also closes $0\le u\le541/2500=.2164$, where its
rational cap is $7042043/12500000=.56336344<40/71$.
The frozen regional certificates below retain their original, valid
choice $\rho=.849$.

\section{A support--source comparison}\label{sec:source}
Let $Z$ be standard Gaussian, independent of a transverse Gaussian
vector $W$. Suppose $|h|+|k|\le1$, set
\[
 u=\E[Zh]\ge0,\qquad p=\E|k|,\qquad
 r=\|P_3^W\E_Z k\|_2.
\]
The source here is the transverse cubic projection. Its independence
of $Z$ is essential to the following comparison.
For $a>0$ and $y\ge0$ write
\[
 J_a(y)=\E_Z\max\{a|Z|,y\}
 =ac e^{-y^2/(2a^2)}+y\operatorname{erf}(y/(\sqrt2a)).
\]
\begin{proposition}\label{prop:source}
For every $a>0$ and $\lambda\ge0$,
\begin{equation}\label{eq:sourcegate}
 au+r+\lambda p\le
 (1-\rho^2)J_a(\lambda)+\rho^2J_a(\lambda+1/\rho),
 \qquad \rho=849/1000.
\end{equation}
\end{proposition}
\begin{proof}
Set $F(y)=J_a(y+\lambda)$ and $x=1/\rho$.
Its third derivative on $[0,\infty)$ is nonpositive:
\[
 F'''(y)=-\frac{c(y+\lambda)}{a^3}e^{-(y+\lambda)^2/(2a^2)}\le0.
\]
Let $D=F(x)-F(0)$ and interpolate with the quadratic
\[
 H(y)=F(0)+\beta y+\gamma y^2,\quad
 \beta=2D/x-F'(x),\quad \gamma=[xF'(x)-D]/x^2.
\]
The Hermite interpolation remainder is
$F'''(\xi)y(y-x)^2/6$, so $F(y)\le H(y)$ for every $y\ge0$.
Since $F'$ is increasing and concave, the trapezoid comparison gives
$D\ge x(F'(0)+F'(x))/2$, hence $\beta\ge0$; also
$D\le xF'(x)$ implies $\gamma\ge0$.
For any $Y\ge0$ with $\E Y^2=1$ and $\E Y\le\rho$, therefore,
\[
 \E F(Y)\le F(0)+\beta\rho+\gamma
 =(1-\rho^2)F(0)+\rho^2F(1/\rho).
\]

For $r>0$ normalize $P_3^W\E_Z k$ to a variance-one cubic $Q(W)$.
Then $\E[Qk]=r$, and pointwise
\[
 aZh+Qk+\lambda|k|\le\max\{a|Z|,|Q|+\lambda\}.
\]
Independence and Theorem~\ref{thm:l1} allow the preceding comparison
with $Y=|Q|$, proving \eqref{eq:sourcegate}. If $r=0$, choose any
normalized transverse cubic: its pairing with $k$ is zero by orthogonality.
An unused Gaussian coordinate can be adjoined if necessary.
\end{proof}

At $\lambda=0$, minimizing the right side of \eqref{eq:sourcegate}
minus $au$ gives the explicit whole-amplitude source cap
\begin{equation}\label{eq:Rrho}
 r\le R_\rho(u)=
 \begin{cases}
 \rho,&0\le u\le c(1-\rho^2),\\
 \rho\operatorname{erf}\!\sqrt{-\log v},
 &c(1-\rho^2)<u<c,\\
 0,&u=c,
 \end{cases}
 \quad v=\frac{u/c-(1-\rho^2)}{\rho^2}.
\end{equation}
Indeed the derivative of
$(1-\rho^2)ac+\rho^2J_a(1/\rho)$ is
$c[1-\rho^2+\rho^2e^{-1/(2a^2\rho^2)}]$.
In the middle range the minimizing value is
$a=1/(\rho\sqrt{-2\log v})$. The two endpoints follow by limits.
Retaining $\lambda>0$ preserves the actual disagreement mass in the
comparison. Additional source information is still needed to close the
full operator norm bound.

\begin{lemma}[Physical clipping of source cells]
\label{lem:physical-source-clipping}
Let $c=\sqrt{2/\pi}$, $0<\rho<1$, and use the support bound
$r\le R_\rho(u)$ of \eqref{eq:Rrho}.  Define
\[
 U_\rho(r)=c\{1-\rho^2+\rho^2
  e^{-[\operatorname{erf}^{-1}(r/\rho)]^2}\},
 \qquad 0<r<\rho,
\]
with $U_\rho(0)=c$ and $U_\rho(\rho)=c(1-\rho^2)$.
Every physical point in a cell
$[u_-,u_+]\times[r_-,r_+]$ also satisfies
\[
 u\le\min\{u_+,U_\rho(r_-)\}.
\]
Thus any outward rational upper bound on $U_\rho(r_-)$ may replace
the cell's amplitude endpoint in all native and R3 price maxima.
If this upper bound is below $u_-$, the cell is physically empty.
\end{lemma}
\begin{proof}
The source envelope is decreasing.  For $0<r<\rho$, solving
$R_\rho(u)=r$ gives the displayed inverse; its endpoint values
follow by limits.  Hence $r\ge r_-$ and $r\le R_\rho(u)$ imply
the asserted restriction.  To certify the inverse without assuming
a special-function inverse, bracket
$s=\operatorname{erf}^{-1}(r_-/\rho)$ by rational numbers
$s_-\le s\le s_+$ using directed monotone evaluations of
$\operatorname{erf}$.  Since $e^{-s^2}$ decreases for $s\ge0$,
an outward evaluation using $s_-$ gives the required upper bound.
\end{proof}

\subsection{Retaining the longitudinal cubic moment}
The preceding comparison can also retain
$v=\E[H_3(Z)h]$, with $H_3(z)=(z^3-3z)/\sqrt6$.
Take $A>0$, $B\ge0$, and write
\[
 L(z)=Az+Bz^3=az+dH_3(z),\qquad a=A+3B,\quad d=\sqrt6B,
 \qquad J_{A,B}(x)=\E_Z\max\{|L(Z)|,|x|\}.
\]
\begin{proposition}\label{prop:cubic-source}
For every real $x$,
\begin{equation}\label{eq:cubic-majorant}
 J_{A,B}(x)\le c\left[
 A+2B+\frac{x^2}{2A}-\frac{A+6B}{24A^4}x^4
 +\frac{A^2+20AB+120B^2}{240A^7}x^6\right].
\end{equation}
If $r>0$ and $Q=P_3^W\E_Zk/r$ has cumulant $\kappa$ and shape
$V=\E[(\|DQ\|^2/3-1)^2]$, then
\begin{align}\label{eq:cubic-source-price}
 au+dv+r\le c\bigg[&q_0+q_2+(3+\kappa)q_4\notag\\
 &+q_6\left\{15+15\kappa+K\kappa^{3/2}
 -633\sqrt\kappa(V-7\kappa/45)\right\}\bigg],
\end{align}
where $(q_0,q_2,q_4,q_6)$ are the four coefficients in the brackets
of \eqref{eq:cubic-majorant}.
\end{proposition}
\begin{proof}
For $x\ge0$ let $z\ge0$ solve $x=Az+Bz^3$.
Differentiating the distribution function gives
\[
 J'_{A,B}(x)=\Pr(|Z|\le z),\qquad
 J''_{A,B}(x)=\frac{ce^{-z^2/2}}{A+3Bz^2}.
\]
Moreover $J_{A,B}(0)=c(A+2B)$ and $J'_{A,B}(0)=0$.
Set $w=z^2$, $q=x^2=w(A+Bw)^2$, and
$F(q)=J''_{A,B}(\sqrt q)$. With $D_1=A+Bw$ and $D_3=A+3Bw$,
\[
 F(q)=ce^{-w/2}/D_3,\qquad
 \frac d{dq}=\frac1{D_1D_3}\frac d{dw}.
\]
The recurrence
\[
 P_0=1,\qquad
 P_{j+1}=(P_j/2-\partial_wP_j)D_1D_3
 +2jBP_jD_3+3(2j+1)BP_jD_1
\]
gives
\[
 (-1)^jF^{(j)}(q)=
 \frac{ce^{-w/2}P_j(A,B,w)}{D_1^{2j}D_3^{2j+1}}.
\]
For $j=1,2,3$, the expanded numerator has respectively 5, 12 and
22 nonnegative coefficients. The explicit third numerator is given in
Appendix~\ref{app:cubic-source}; hence $F'''(q)\le0$ on the whole
half-line. Taylor's formula with its integral remainder gives
\[
 F(q)\le\frac cA-\frac{c(A+6B)}{2A^4}q
 +\frac{c(A^2+20AB+120B^2)}{8A^7}q^2.
\]
Integrating twice in $x$, using the two initial values, proves
\eqref{eq:cubic-majorant}. Evenness covers negative $x$.

Pointwise $L(Z)h+Qk\le\max\{|L(Z)|,|Q|\}$, so independence and
projection orthogonality give $au+dv+r\le\E J_{A,B}(Q)$.
Apply \eqref{eq:cubic-majorant} and Theorem~\ref{thm:633}.
The fourth moment is exact, so its negative coefficient causes no
inequality reversal; the sixth coefficient is positive.
\end{proof}
This comparison retains the actual longitudinal cubic moment.
A global norm proof must quantify that moment jointly with the same
source and profile; assigning it a favorable numerical value would not
prove an operator bound.

\begin{proposition}[Cubic-head contact source gate]\label{prop:contact-head-source}
Let $A>0$, $B\ge0$, $L(z)=Az+Bz^3=a z+dH_3(z)$, where
$H_3(z)=(z^3-3z)/\sqrt6$, $a=A+3B$, $d=\sqrt6B$, and let
$J(x)=\mathbb E_Z\max\{|L(Z)|,|x|\}$.
Set $G(s)=J(\sqrt s)$ on $s\ge0$.  Then
$G''\le0$, $G'''\ge0$ and $G''''\le0$.
For every $s_0>0$, the cubic Hermite interpolant $P$ matching
$G(0),G'(0),G(s_0),G'(s_0)$ satisfies
\[
 J(x)\le P(x^2)=p_0+p_2x^2+p_4x^4+p_6x^6
 \quad(x\in\mathbb R),\qquad p_6\ge0.
\]
Here $p_0=\nu(A+2B)$, $p_2=\nu/(2A)$, $\nu=\sqrt{2/\pi}$, and, with
$D=G(s_0)-p_0-p_2s_0$ and $E=G'(s_0)-p_2$,
\[
 p_4=(3D-s_0E)/s_0^2,\qquad p_6=(s_0E-2D)/s_0^3.
\]
Any componentwise upward rational rounding $q_j\ge p_j$ gives a global
majorant with $q_6\ge0$.  Its coefficients already contain $\nu$.
\end{proposition}
\begin{proof}
Write $s=w(A+Bw)^2$, $D_1=A+Bw$, $D_3=A+3Bw$.
Since $ds/dw=D_1D_3>0$, Gaussian tail differentiation gives
\[
 G'(s)=\frac\nu2\int_0^1
       \frac{e^{-t^2w(s)/2}}{A+Bw(s)}\,dt.
\]
For $K_\tau(s)=e^{-\tau w(s)}/(A+Bw(s))$, $\tau\ge0$, differentiation
shows
\[
 K_\tau^{(j)}(s)=(-1)^j
  \frac{e^{-\tau w}P_j(A,B,w,\tau)}{D_1^{2j+1}D_3^{2j}},\qquad j=1,2,3.
\]
Starting from $P_0=1$, the numerator recurrence is
\[
 P_{j+1}=(\tau P_j-\partial_wP_j)D_1D_3
       +(2j+1)BP_jD_3+6jBP_jD_1.
\]
For $j=1,2,3$, exact expansion gives respectively $5$, $12$ and $22$
nonnegative coefficients.  Integrating the derivative signs at
$\tau=t^2/2$ proves the claimed signs of $G$.
The Hermite interpolation error is
\[
 G(s)-P(s)=G''''(\xi)s^2(s-s_0)^2/4!\le0
 \quad(s\ge0),
\]
including outside the interpolation interval.  The leading coefficient
$p_6=G[0,0,s_0,s_0]\ge0$ by the sign of $G'''$.
\end{proof}

\begin{corollary}[Moment price and retained longitudinal completion]
Let $h=(f+g)/2$, $k=(f-g)/2$, $P_1h=uZ$, and
$v=\mathbb E[H_3(Z)h]$.  Put
$r=\|P_3^W\mathbb E_Zk\|_2$ and
$Q=P_3^W\mathbb E_Zk/r$ when $r>0$.
Write $\kappa=\E Q^4-3$ and
$\E(\|DQ\|^2/3-1)^2=\kappa(7/45+\zeta/15)$, with $0\le\zeta\le1$.
The endpoint $\kappa=0$ is interpreted by continuity.
For a majorant from Proposition~\ref{prop:contact-head-source}, the physical
diamond and the coefficient-$633$ moment theorem give
\[
 au+dv+r\le A_q(\kappa,\zeta)
 =q_0+q_2+(3+\kappa)q_4+
 \left[15+15\kappa+
  \left(\frac{194}{\sqrt{10}}-\frac{633\zeta}{15}\right)\kappa^{3/2}\right]q_6.
\]
No extra $\nu$ multiplies this price.  For a fixed row the price is convex
in $\kappa$ and affine in $\zeta$.
If a certified frame completed its longitudinal penalty by
$-bv^2\le2\beta v+\beta^2/b$, with $\beta$ multiplying normalized $H_3$,
then for any fixed $\lambda\ge0$ its saved completed bound may be augmented by
\[
 \lambda\left(A_q-au-r+\frac{\beta d}{b}\right)
       +\frac{\lambda^2d^2}{4b}.
\]
With $\lambda$ fixed throughout each price rectangle, the existing cumulant,
shape, amplitude and source-chord certification can be reused.
If a frame instead stores $\beta_{\rm raw}$ multiplying $z^3-3z$, use
$\beta=\sqrt6\,\beta_{\rm raw}$.  The same correction is then
\[
 \lambda\left(A_q-(A+3B)u-r+\frac{6\beta_{\rm raw}B}{b}\right)
       +\frac{3\lambda^2B^2}{2b}.
\]
Both the contact row and $\lambda$ are chosen once for each whole cell;
they cannot be optimized independently at its corner points and then treated
as a single convex price.
\end{corollary}
\begin{proof}
Pointwise $L(Z)h+Qk\le\max(|L(Z)|,|Q|)$, and $\mathbb EQk=r$.
Take expectations, use the global majorant, and price its second, fourth and
sixth moments.  The sixth coefficient is nonnegative and
$194/\sqrt{10}-633/15>0$, proving the price properties.
Before completing the longitudinal penalty, add the nonnegative quantity
$\lambda(A_q-au-dv-r)$ and complete with
$(\beta+\lambda d/2)^2/b$ in place of $\beta^2/b$.
Subtracting the old completion gives the displayed correction.
\end{proof}

For the exact head $A=26412631/10^7$, $B=4084047/10^7$ and contact
$w_0=189/50$, so $s_0=w_0(A+Bw_0)^2$, one certified rational row is
\[
\begin{split}
q_0&=137957132894348330419/(5\cdot10^{19}),\\
q_2&=1510422344526876849/10^{19},\\
q_4&=-163849393411225061/10^{20},\\
q_6&=23608461930109/(25\cdot10^{17}).
\end{split}
\]
Its uniform four-corner price on $\kappa\in[2,4]$, $\zeta\in[0,1]$ is at most
$29040583649083683/10^{16}$.
The exact coefficient-sign proof, rational rounding and directed corner
prices are reproduced by
\texttt{certify\_\allowbreak direct\_\allowbreak contact\_\allowbreak majorant.py}.
This source proposition alone does not assert a global operator norm bound.
The scalar comparison was independently accepted by an exact rational
recurrence check and a separate 768-bit MPFR replay of all $28$ upward-rounded
coefficients and $98$ directed moment prices.  A global norm conclusion still
requires complete fixed-row, fixed-multiplier cell coverage.

\begin{corollary}[A physical source exclusion at the remaining native bottleneck]
In the notation of Proposition~\ref{prop:contact-head-source}, every physical
source with $u\ge151/200$, $\kappa\in[2,4]$ and $\zeta\in[0,1]$ satisfies
\[
 r\le\frac{3218564576843929081}{12500000000000000000}
 <\frac{257485167}{10^9}<\frac{129}{500}.
\]
Thus the entire source domain $u\ge.755$, $r\ge.258$, $\kappa\in[2,4]$
is physically impossible.  This exclusion is independent of the operator
coefficient $b$ and uses no longitudinal-completion multiplier.
\end{corollary}
\begin{proof}
Use the globally valid contact family at $A=147/50$, $B=0$, $w_0=11/2$,
so that $s_0=237699/5000$.  One upward-rounded rational row is
\[
\begin{split}
q_0&=234578060876042414629/10^{20},\\
q_2&=13569465319776621699/10^{20},\\
q_4&=-57073674634011611/(5\cdot10^{19}),\\
q_6&=648580884546013/10^{20}>0.
\end{split}
\]
These coefficients already contain $\nu$.  The price for this single fixed
row is convex in $\kappa$ and affine in $\zeta$.  Directed evaluation of the
four corners of $[2,4]\times[0,1]$ gives
\[
 A_q(\kappa,\zeta)
 \le\frac{30964814576843929081}{12500000000000000000}.
\]
The physical source gate is simply $(147/50)u+r\le A_q$ since $d=0$.
Subtract $(147/50)(151/200)$ to obtain the asserted bound.
All four rational coefficient roundings and four corner prices passed a
separate 768-bit MPFR replay; the globally valid contact-family argument
also applies at $B=0$.
This is a source-domain exclusion, not a complete operator norm theorem.
\end{proof}

\section{Conditional Gaussian profiles and a mixed cubic penalty}\label{sec:mixed}
% Self-contained fragment.  Requires amsmath, amssymb, and theorem environments
% named proposition and proof.  No numerical-profile certificate is asserted.
\begin{proposition}[Sharp conditional Gaussian profile and a quartic budget]
\label{prop:conditional-gaussian-budget}
Let $Z$ be standard normal and let measurable real functions $h,k$ satisfy
$|h|+|k|\leq1$.  Put
\[
 \nu=\sqrt{2/\pi},\qquad A=\mathbb E[Zh],\qquad
 p=|\mathbb E k|,\qquad C=1-\nu^2.
\]
If $p=2\Phi(t)-1$, $t\in[0,\infty]$, then
\begin{equation}\label{eq:sharp-conditional-profile}
 |A|\leq\nu e^{-t^2/2}.
\end{equation}
This bound is sharp for every $p$.  In particular,
\begin{equation}\label{eq:quartic-conditional-budget}
 A^2+p^2-Cp^4\leq\nu^2.
\end{equation}
The quartic budget improves the ellipse
$A^2+\nu^2p^2\leq\nu^2$ by $Cp^2(1-p^2)$.
\end{proposition}

\begin{proof}
Write $\alpha=|k|$ and $\alpha_0=\mathbf1_{\{|Z|\leq t\}}$.
Since $\mathbb E\alpha\geq p=\mathbb E\alpha_0$ and
$(|Z|-t)(\alpha-\alpha_0)\geq0$ pointwise,
\[
 \mathbb E[|Z|\alpha]\geq\mathbb E[|Z|\alpha_0].
\]
Consequently $|A|\leq\mathbb E[|Z|(1-\alpha)]\leq2\varphi(t)
=\nu e^{-t^2/2}$.  Equality is attained by
$h=\operatorname{sgn}(Z)\mathbf1_{\{|Z|>t\}}$ and
$k=\mathbf1_{\{|Z|\leq t\}}$.

For the polynomial budget set $p(t)=2\Phi(t)-1$, so
$p'=\nu e^{-t^2/2}$ and $p''=-tp'$.  Define
\[
 D(t)=\nu^2e^{-t^2}-\nu^2+p(t)^2-Cp(t)^4,
 \qquad g(t)=p(t)-2Cp(t)^3-tp'(t).
\]
Then $D'=2p'g$ and $g'=p'(t^2-6Cp^2)$.
The ratio $t/p(t)$ strictly increases from $1/\nu$ to infinity:
indeed $p(t)>tp'(t)$ for $t>0$ because $p'$ strictly decreases.
Furthermore $1/\nu^2<6C$ and $1-2C>0$; for example these follow
already from $3<\pi<22/7$.
Thus $g'$ changes sign exactly once, from negative to positive.
As $g(0)=0$ and $\lim_{t\to\infty}g(t)=1-2C>0$, the function $g$
has exactly one further zero.  It follows that $D$ first decreases
and then increases.  Its two endpoint limits are zero, so $D\leq0$.
Combining this with \eqref{eq:sharp-conditional-profile} proves
\eqref{eq:quartic-conditional-budget}.
\end{proof}

\begin{proposition}[A mixed cubic completion with one scalar source]
\label{prop:mixed-cubic-completion}
Let $Z$ and $W$ be independent standard Gaussian coordinates.  Let $f,g$
be odd signs and put $h=(f+g)/2$, $k=(f-g)/2$,
\[
 A(W)=\mathbb E_Z[Zh],\qquad k_0(W)=\mathbb E_Z k.
\]
Suppose $r=\|P_3^Wk_0\|_2>0$ and write
$Q=P_3^Wk_0/r=I_3(F)$, with $6\|F\|^2=1$.  Here Wiener integrals are
normalized by $\mathbb E I_j(H)^2=j!\|H\|^2$.
Set
\[
 G=F\otimes_2F,\quad s=\|G\|^2,\quad
 D=1+\frac{3\|\operatorname{Sym}(F\otimes_1F)\|^2}{2s},\quad
 R_2=\frac{I_2(G)}{\sqrt{2s}},\quad \kappa=\mathbb E Q^4-3.
\]
Then $\mathbb ER_2=0$, $\mathbb ER_2^2=1$, and
\begin{equation}\label{eq:mixed-source-normalization}
 \theta:=\mathbb E[Q^2R_2]=18\sqrt{2s},\qquad
 \theta^2=\frac{\kappa}{2D}.
\end{equation}
For $t=\mathbb E[AR_2]$, any $b,w>0$ and $a,c\in\mathbb R$ satisfy
\begin{align}
 -bt^2\leq{}&w\{\nu^2+a^2+2ac+c^2(3+\kappa)\}
 -w\mathbb E k_0^2+wC\mathbb E k_0^4
 \notag\\[-1mm]
 &-2w\mathbb E[(a+cQ^2)A]
 -\frac{bw c^2\kappa}{2D(b+w)},
 \label{eq:quartic-mixed-completion}
\end{align}
where $\nu^2=2/\pi$ and $C=1-\nu^2$.
The direction $ZR_2(W)$ is a unit cubic Hermite function orthogonal to
$H_3(Z)=(Z^3-3Z)/\sqrt6$.  Hence its penalty $-bt^2$ can be retained
in addition to a completion using only the longitudinal cubic direction.
\end{proposition}

\begin{proof}
The Gaussian product formula yields
\[
 P_2(Q^2)=18I_2(G),\qquad
 \kappa=1296s+1944\|\operatorname{Sym}(F\otimes_1F)\|^2=1296sD.
\]
These identities give \eqref{eq:mixed-source-normalization}.
Also $G$ is positive semidefinite, since
$G_{ij}=\sum_{a,b}F_{iab}F_{jab}$.  Thus $R_2$ is the normalized,
centered quadratic form determined by this actual cubic source.

Put $\eta=a+cQ^2$.  Orthogonal projection onto the unit vector $R_2$
gives $\mathbb E(A-\eta)^2\geq(t-c\theta)^2$.  Completing one scalar
square gives
\[
 bt^2+w(t-c\theta)^2\geq\frac{bw}{b+w}c^2\theta^2,
\]
and therefore
\[
 -bt^2\leq w\mathbb E(A-\eta)^2-
 \frac{bw}{b+w}c^2\theta^2.
\]
Conditionally on $W$, $h,k$ satisfy $|h|+|k|=1$, so
Proposition~\ref{prop:conditional-gaussian-budget} gives
$A^2\leq\nu^2-k_0^2+Ck_0^4$.  Expanding the square and using
$\mathbb E\eta^2=a^2+2ac+c^2(3+\kappa)$ proves
\eqref{eq:quartic-mixed-completion}.  The final orthogonality claim
follows from independence of $Z,W$ and $\mathbb ER_2=0$.
\end{proof}

\paragraph{Insertion into a local scalar criterion.}
If the original conditional potential contains $2\ell A$ and local
moments $\mathbf k=(k_0,k_1,k_2,k_3)$, formula
\eqref{eq:quartic-mixed-completion} replaces
\[
 \ell\quad\hbox{by}\quad\ell-w(a+cq^2)
\]
at scalar source value $q$, adds the local polynomial
$-wk_0^2+wCk_0^4$, and adds only the explicit outer price
\[
 w\{\nu^2+a^2+2ac+c^2(3+\kappa)\}
 -\frac{bw c^2\kappa}{2D(b+w)}.
\]
It introduces no additional source or local moment coordinate.
This insertion is an exact inequality, not a global norm certificate:
the modified local potential and its source tail must still be bounded
on their complete domains.  In particular the quadratic dependence
on $q$ requires a new chord allowance.  If the unmodified local
potential is a supremum of affine functions of $q$, then its modified
physical supremum $\Psi$ obeys, for $q\in[q_0,q_1]$,
\[
 \Psi(q)\leq\operatorname{Chord}_{q_0,q_1}\Psi(q)
       +2w|c|\nu(q-q_0)(q_1-q).
\]
Indeed adding $2w|c|\nu q^2$ makes every candidate convex, because
$|A|\leq\nu$.  Any pre-existing curvature allowance must also be retained.

\section{Continuous regional norm certificates}\label{sec:regional}
We next convert local inequalities into bounds on complete intervals of
the agreement amplitude. All source and cumulant parameters are covered
continuously; finite sampling is used only to choose the rational inputs.

Throughout this section the longitudinal Hermites are orthonormal:
\[
 H_0=1,\quad H_1=Z,\quad H_2=(Z^2-1)/\sqrt2,\quad
 H_3=(Z^3-3Z)/\sqrt6.
\]
Section~\ref{sec:r3} explicitly switches to the raw Hermite convention.

\begin{proposition}[A transverse Gram comparison]\label{prop:gram}
Let $b>0$, let $f,g$ be odd signs, put $h=(f+g)/2$, $k=(f-g)/2$,
and rotate so that $P_1h=uZ$. In the transverse coordinates $W$, set
$k_j(W)=\E_Z[H_j(Z)k]$, $h_j(W)=\E_Z[H_j(Z)h]$, and
$r=\|P_3^Wk_0\|_2$. Suppose the even block $E$ on coordinates $(0,2)$
and odd block $R$ on coordinates $(1,3)$ satisfy
\[
 E\succeq0,\quad E-\operatorname{diag}(-1,b)\succeq0,
 \quad R\succeq0,\quad R-\operatorname{diag}(-1,b)\succeq0,
 \quad R-\operatorname{diag}(b,0)\succeq0,
 \qquad E_{22}>0.
\]
Let $M$ be the corresponding four-coordinate matrix. For arbitrary
real $\ell,\beta,\delta,\omega$, define
\[
 s=(\delta,0,-\omega,0),\qquad
 \xi=\frac{(\omega-E_{02})^2}{E_{22}}-E_{00}+b-2\delta.
\]
Let $\mathcal P(t)$ be the supremum, over conditional functions satisfying
$|h|+|k|\le1$, of
\[
 \mathbf k^TM\mathbf k+2t\,s^T\mathbf k
       +2\ell h_1+2\beta h_3+\beta^2/b,
 \qquad \mathbf k=(k_0,k_1,k_2,k_3)^T.
\]
For $r>0$, with $Q=P_3^Wk_0/r$, one has
\begin{equation}\label{eq:gram-regional}
 \langle f,T_bg\rangle\le
 u^2-2\ell u+\E\mathcal P(rQ)+\xi r^2.
\end{equation}
For $r=0$, replace the expectation by $\mathcal P(0)$.
\end{proposition}
\begin{proof}
Retain $-b(\E h_3)^2$ from the agreement penalty and use
$-bx^2\le2\beta x+\beta^2/b$. Expand the disagreement function by
longitudinal and transverse Hermite degree. In transverse degree one,
the even block must dominate $\operatorname{diag}(-1,b)$; in transverse
degrees zero and two, the odd block must dominate the two displayed
diagonal matrices. In the source direction $Q$, write
$a=\E[Qk_2]$. The even quadratic and source terms are
\[
 E_{22}a^2+2(E_{02}-\omega)ra+(E_{00}+2\delta+\xi)r^2
 \ge br^2,
\]
by their exact Schur completion. The component of $P_3^Wk_0$
orthogonal to $Q$ is zero by definition. All remaining transverse
components are dominated by the positive blocks. Finally insert
$2\ell\E Zh-2\ell u=0$. These comparisons give
\eqref{eq:gram-regional} after conditional maximization.
\end{proof}

For completeness, the finite local comparison is particularly simple.
If the actual conditional moment vector lies in a box centered at $c$
with half widths $w$, exact quadratic completion and the diamond
constraint give the entire box bound
\begin{align}\label{eq:physical-box}
 \mathcal L(t;h,k)\le{}&
 2\E_Z\max\{|(Mc+ts)\cdot H|,|\ell H_1+\beta H_3|\}
       -c^TMc+\beta^2/b+w^T|M|w,
\end{align}
where $\mathcal L$ denotes the objective defining $\mathcal P$ and
$|M|$ is entrywise absolute value. This is a bound for actual moments
in that box, with no stationary-point assumption. Bessel's inequality
and $|k_j|\le\E|H_j|$ restrict the moments to
\[
 \sum_{j=0}^3k_j^2\le1,\qquad
 (|k_0|,|k_1|,|k_2|,|k_3|)\le(1,4/5,7/10,5/8).
\]
The expectation in \eqref{eq:physical-box} is the sum of two Gaussian
absolute-cubic integrals, since
$\max(|A|,|B|)=(|A+B|+|A-B|)/2$.
Certified real-root isolation and Gaussian antiderivatives evaluate
these integrals with outward rounding.

The potential $\mathcal P$ is even and convex in $t$. Its certified
nodal heights therefore give upper chords on every source interval;
a tail slope at least $2\|s\|$ gives an affine extension to infinity.
Denote this even convex upper function by $F$. To price $\E F(rQ)$,
solve $g_F'(x)-xg_F(x)=F(rx)-\E F(rN)$ for a standard normal $N$.
For any $\lambda\ge0$, any odd regression $\psi$, and $\eta>0$, the
same Stein calculation as in Section~\ref{sec:stein} gives
\begin{align}\label{eq:regional-scalar-price}
 \E F(rQ)\le{}&\E F(rN)+\lambda(\rho-c)
 +\E\left[\psi'(Q)-Q\psi(Q)
 +J_\eta(g_F'(Q)-\lambda g'(Q)-\psi'(Q))\right]
 +\frac\eta2\E U^2,\\
 J_\eta(x)={}&
 \begin{cases}x^2/(2\eta),&x\le\eta,\\x-\eta/2,&x\ge\eta.
 \end{cases}\notag
\end{align}
Here $c=\sqrt{2/\pi}$, $g$ is the absolute-value Stein solution above,
and $\rho=849/1000$. The clipped quadratic follows by maximizing
$-vx-\eta v^2/2$ over $v\ge-1$, since $U\ge-1$.
Smooth approximation justifies the finitely many target knots.

A whole-line sextic majorant for the bracket in
\eqref{eq:regional-scalar-price}, with nonnegative sextic coefficient,
has a price convex in $\kappa$ and affine in the shared variance shape.
Consequently its four endpoint comparisons cover each complete
cumulant--shape rectangle. Source interpolation adds only
\[
 (-\xi)_+\max_i(r_{i+1}-r_i)^2/4
\]
to the endpoint prices, because $\E F(rQ)$ is convex in $r$.
Checking the two endpoints of $u^2-2\ell u$ covers an entire amplitude
interval. These facts explain the continuous scope of the following
certificates.

\begin{theorem}[Complete regional amplitude bounds]\label{thm:regional-coverage}
For odd Gaussian signs $f,g$, put $u=\|P_1((f+g)/2)\|_2$.
In every finite Gaussian dimension,
\[
 \langle f,T_{.72}g\rangle<
 \begin{cases}
 563367/10^6,&0\le u\le341/500,\\[2pt]
 563334/10^6,&77/100\le u\le\sqrt{2/\pi}.
 \end{cases}
 \qquad <\frac{40}{71}.
\]
\end{theorem}
\begin{proof}
The seven frozen profiles in Appendix~\ref{app:regional} cover the
intervals $[.14,.42]$, $[.30,.485]$, $[.40,.54]$, $[.50,.638]$,
$[.60,.662]$, $[.662,.682]$, and $[.77,.797885]$.
Each satisfies Proposition~\ref{prop:gram}, complete local moment-box
comparisons, and complete scalar inequalities of
\eqref{eq:regional-scalar-price}. Directed prices and exact convex
endpoint comparisons cover every source amplitude, cumulant and
shared shape in each region. Their first six rational caps are all
at most $563367/10^6$; the seventh cap is $563334/10^6$.
For $u\le.14$, the energy comparison gives
$u^2+.72\rho^2\le.53857672$. The intervals therefore cover every
$u\le.682$. Finally $u\le\E|Z|=\sqrt{2/\pi}<.797885$, so the last
interval contains every feasible $u\ge.77$.
\end{proof}

The accepted broad profile uses a source grid ending exactly at
$\rho$, which is valid at every agreement amplitude by projection
norm duality. Hence its original discovery interval imposes no
restriction on its larger accepted interval. The other extensions
retain their verified outward source endpoint. This accounts for
the continuous scope of the theorem. The remaining entire amplitude
range is $.682<u<.77$; a local source-parameter subregion inside that
range does not complete its full norm cover.

\section{Retaining the actual cubic projection}\label{sec:r3}
% Self-contained mathematical fragment. Requires amsmath, amssymb, amsthm,
% and proposition/proof environments. All Hermites here are unnormalized.
Let $T_b=P_1-bP_3$, $b>0$, on real standard Gaussian space, where $P_j$
is the orthogonal projection onto total Hermite degree $j$.
For odd signs $f,g$ put $h=(f+g)/2$, $k=(f-g)/2$.
Then $(h,k)$ takes the four values $(\pm1,0),(0,\pm1)$.
Rotate the Gaussian coordinates as $(Z,W)$ so that $P_1h=uZ$,
with $Z$ independent of $W$.  Use the raw longitudinal Hermites
\[
 H_0=1,\quad H_1=Z,\quad H_2=Z^2-1,\quad H_3=Z^3-3Z,
 \qquad m_j(W)=\mathbb E_Z[H_jk].
\]
If $r=\|P_3^Wm_0\|_2>0$, define the normalized cubic source
$Q=P_3^Wm_0/r$.  Thus $\mathbb EQ^2=1$ and
$\mathbb E[Qm_0]=r$.  All remaining projections act on $W$.

\begin{proposition}[An exact residual of the physical cubic source]
\label{prop:physical-r3-source}
Put
\[
 \kappa=\mathbb EQ^4-3,\quad c_6=\mathbb EQ^6-15\mathbb EQ^4+30,
 \quad S(q)=q^3-3q,
\]
\[
 \mathcal R=Q^3-(3+\kappa)Q-P_1^W(Q^3),\quad
 \Delta=\mathcal R-S(Q),\quad
 D=\kappa^2+\|P_1^W(Q^3)\|_2^2.
\]
The residual $\mathcal R$ is odd, orthogonal to first chaos and to $Q$, and
\begin{equation}\label{eq:physical-r3-identities}
 \mathbb E\mathcal R^2=6+9\kappa+c_6-D,\quad
 \mathbb E\Delta^2=D,\quad
 \mathbb E[Q\Delta]=-\kappa,\quad
 \mathbb E[S(Q)\Delta]=-D.
\end{equation}
\end{proposition}
\begin{proof}
Since $Q$ has pure degree three, $P_1^W(Q^3)$ is orthogonal to $Q$.
The coefficient removed from $Q^3$ along the unit vector $Q$ is
$\mathbb EQ^4=3+\kappa$.  Also
$\Delta=-\kappa Q-P_1^W(Q^3)$.  Orthogonal expansion of these two
expressions proves all four identities.
\end{proof}

Write $Q=I_3(F)$ and put $s_F=\|F\otimes_2F\|_F^2$ and
$z=1296s_F/\kappa\in[2/5,1]$, as in the tensor shape comparison
of Section~\ref{sec:cumulants}.
The residual parameter has the dimension-free bounds
\begin{equation}\label{eq:r3-D-price}
 \kappa^2\le D\le\kappa^2+\frac94(\kappa z)^{3/2}.
\end{equation}
Indeed Gaussian differentiation and the product formula give
$\E[X_iQ^3]=324\langle F_i,F\otimes_2F\rangle$.
The flattening map of $F$ has squared operator norm at most $\sqrt{s_F}$,
so the sum of the squared inner products is at most $s_F^{3/2}$.
Multiplication by $324^2$ proves \eqref{eq:r3-D-price}.

\begin{proposition}[An inverse-free physical dual criterion]
\label{prop:physical-r3-dual}
Fix $p=(p_0,p_2)$ and a symmetric matrix
$R=\left(\begin{smallmatrix}R_{11}&R_{13}\\R_{13}&R_{33}\end{smallmatrix}\right)$
such that
\begin{equation}\label{eq:physical-r3-gates}
 p_2^2\geq\frac b2(1+p_0^2),\qquad
 R\succeq0,\quad R-\operatorname{diag}(-1,b/6)\succeq0,
 \quad R-\operatorname{diag}(b,0)\succeq0.
\end{equation}
Choose real $\ell,\beta,\gamma,y,\zeta,c,d$ and $G>0$.  Write
\[
 \sigma^2=p_0^2+2p_2^2,\quad \rho=1+\zeta^2/G,\quad
 \mathfrak g(q)=cq+dS(q),\quad a(q)=\zeta[-S(q)+\mathfrak g(q)/(2G)].
\]
For $v=(s,v_1,v_3)$ define
\begin{align*}
 L(z)&=\ell z+\beta H_3(z),\\
 B(q,z;v)&=q\gamma+q(y-p_0p_2r)H_2(z)
       +[a(q)+\rho s][p_0+p_2H_2(z)]\\
 &\quad +(R_{11}v_1+R_{13}v_3)z
       +(R_{13}v_1+R_{33}v_3)H_3(z),
\end{align*}
and the physical dual ellipsoid
\[
 \mathcal E_p=\{(s,v_1,v_3):s^2/\sigma^2+v_1^2+v_3^2/6\leq1\}.
\]
Put
\begin{equation}\label{eq:physical-r3-local}
 \Phi(q)=\sup_{v\in\mathcal E_p}\left\{
 2\mathbb E_Z\max(|L(Z)|,|B(q,Z;v)|)
 -\rho s^2-R_{11}v_1^2-2R_{13}v_1v_3-R_{33}v_3^2
 +\frac{\mathfrak g(q)^2}{4G}\right\}.
\end{equation}
Then
\begin{align}\label{eq:physical-r3-criterion}
 \langle T_bf,g\rangle\leq{}&u^2-2\ell u+\frac{6\beta^2}{b}
 +(b-p_0^2)r^2-2\gamma r+\frac{y^2}{p_2^2}
 +\mathbb E\Phi(Q)\notag\\
 &+\zeta^2(6+9\kappa+c_6)-c\kappa+(G-d-\zeta^2)D.
\end{align}
The ellipsoid in \eqref{eq:physical-r3-local} may be restricted to
$v_1\geq0$ without changing the supremum, and $\Phi(-q)=\Phi(q)$.
\end{proposition}
\begin{proof}
Let $V=p_0m_0+p_2m_2$ and $t=\mathbb E[Qm_2]$.
Let $\Pi$ project onto odd transverse degrees at least three,
orthogonally to $Q$.  Hermite orthogonality and
\eqref{eq:physical-r3-gates} give
\begin{align*}
 -\|P_1k\|_2^2+b\|P_3k\|_2^2
 \leq{}&\mathbb EV^2+
 \mathbb E(R_{11}m_1^2+2R_{13}m_1m_3+R_{33}m_3^2)\\
 &+(b-p_0^2)r^2-2p_0p_2rt-p_2^2t^2-\|\Pi V\|_2^2.
\end{align*}
Indeed the relevant raw-Hermite blocks are
$\operatorname{diag}(-1,b/2)$ on $(m_0,m_2)$ in transverse degree one,
$\operatorname{diag}(-1,b/6)$ on $(m_1,m_3)$ in degree zero, and
$\operatorname{diag}(b,0)$ on that latter pair in degree two.
The first gate is exactly the first block inequality for $pp^T$.
In even longitudinal/transverse degree three,
$P_3^Wm_0=rQ$; subtracting the $Q$ component of $V$ leaves precisely
the indicated unused loss.  Higher odd transverse degrees supply the
remaining part of $\|\Pi V\|_2^2$.

Now $\mathcal R$ lies in the range of $\Pi$, so
\[
 -\|\Pi V\|_2^2\leq-2\zeta\mathbb E[\mathcal RV]
                          +\zeta^2\mathbb E\mathcal R^2.
\]
The pointwise square inequality
\[
 -2\zeta\Delta V\leq \mathfrak g(Q)\Delta+G\Delta^2+
                    (\zeta V+\mathfrak g(Q)/2)^2/G
\]
and \eqref{eq:physical-r3-identities} yield the last line of
\eqref{eq:physical-r3-criterion}; the remaining local quadratic is
$\rho V^2$ with linear coefficient $2a(Q)V$ and constant $\mathfrak g(Q)^2/(4G)$.
Complete $-p_2^2t^2\leq2yt+y^2/p_2^2$, and insert the exact zeros
$2\gamma\mathbb E[Qm_0]-2\gamma r$ and
$2\ell\mathbb E[Zh]-2\ell u$.
Finally
$-b\|P_3h\|_2^2\leq-(b/6)(\mathbb E[H_3h])^2
\leq2\beta\mathbb E[H_3h]+6\beta^2/b$.

For each fixed $W$, Bessel's inequality gives
$m_0^2+m_1^2+m_2^2/2+m_3^2/6\leq\mathbb E_Zk^2\leq1$.
Hence $(V,m_1,m_3)\in\mathcal E_p$ by Cauchy--Schwarz in the even block.
For a positive quadratic form $A$, the identity
\[
 x^TAx=\sup_v\{2v^TAx-v^TAv\}
\]
holds also when $A$ is singular, with equality at $v=x$.
Apply it to $\operatorname{diag}(\rho,R)$ and interchange the two
suprema.  Maximizing the resulting linear functional over the four
pointwise possibilities for $(h,k)$ gives $2\max(|L|,|B|)$.
Restricting the dual supremum to $\mathcal E_p$ loses no value because
it contains every physical moment vector.  Conversely every dual vector
gives at most the physical supremum by the same quadratic identity.
This proves \eqref{eq:physical-r3-local} without a matrix inverse.

Under $z\mapsto-z$, the simultaneous reflection
$(v_1,v_3)\mapsto(-v_1,-v_3)$ preserves the value and the ellipsoid.
Thus $v_1\geq0$ suffices.  Under $(q,s)\mapsto(-q,-s)$ the even part
of $B$ changes sign; reflecting $z$ and changing the overall sign
of $B$ restores the same absolute value.  This proves the evenness in $q$.
\end{proof}

% Extends the previously proved physical R3 criterion. Raw Hermites.
\begin{proposition}[A full even-block criterion]
\label{prop:full-even-r3}
Retain the physical residual-source notation, including
$m_j=\mathbb E_Z[H_jk]$.
Let $E=(E_{ij})_{i,j\in\{0,2\}}$, $R=(R_{ij})_{i,j\in\{1,3\}}$
be symmetric matrices and $p=(p_0,p_2)^T$ satisfying
\begin{align*}
 E&\succeq0,& E-\operatorname{diag}(-1,b/2)&\succeq0,&
 E-pp^T&\succeq0,\\
 R&\succeq0,& R-\operatorname{diag}(-1,b/6)&\succeq0,&
 R-\operatorname{diag}(b,0)&\succeq0.
\end{align*}
In particular $E_{22}\geq b/2>0$.  For arbitrary real
$\ell,\beta,\gamma,y,\zeta,c,d$ and $G>0$, put
\begin{gather*}
 S(q)=q^3-3q,\quad \mathfrak g(q)=cq+dS(q),\\
 a(q)=\zeta[-S(q)+\mathfrak g(q)/(2G)],\quad L(z)=\ell z+\beta H_3(z).
\end{gather*}
With the coordinate order $(0,1,2,3)$, let $N$ have even block
$E+(\zeta^2/G)pp^T$ and odd block $R$, and set
\[
 b_0(q)=(q\gamma+a(q)p_0,\,0,\,
           q(y-E_{02}r)+a(q)p_2,\,0)^T,
 \quad H(z)=(H_0,H_1,H_2,H_3)^T.
\]
Define the physical ellipsoid and local dual value by
\begin{align*}
 \mathcal E&=\{v:v_0^2+v_1^2+v_2^2/2+v_3^2/6\leq1\},\\
 \Phi_E(q)&=\sup_{v\in\mathcal E}\left\{
 2\mathbb E_Z\max\bigl(|L(Z)|,|(b_0(q)+Nv)^TH(Z)|\bigr)
 -v^TNv+\frac{\mathfrak g(q)^2}{4G}\right\}.
\end{align*}
Then
\begin{align*}
 \langle T_bf,g\rangle\leq{}&
 u^2-2\ell u+\frac{6\beta^2}{b}
 +(b-E_{00})r^2-2\gamma r+\frac{y^2}{E_{22}}
 +\mathbb E\Phi_E(Q)\\
 &+\zeta^2(6+9\kappa+c_6)-c\kappa+(G-d-\zeta^2)D.
\end{align*}
The supremum may be restricted to $v_1\geq0$, and $\Phi_E$ is even.
\end{proposition}
\begin{proof}
Put $t=\mathbb E[Qm_2]$, $V=p_0m_0+p_2m_2$, and let $\Pi$
project onto odd transverse degrees at least three, orthogonally to
$Q$.  Hermite orthogonality gives
\begin{align*}
 -\|P_1k\|_2^2+b\|P_3k\|_2^2\leq{}&
 \mathbb E(m_0,m_2)E(m_0,m_2)^T+
 \mathbb E(m_1,m_3)R(m_1,m_3)^T\\
 &+(b-E_{00})r^2-2E_{02}rt-E_{22}t^2
 -\|\Pi V\|_2^2.
\end{align*}
To verify this inequality, the transverse degree-one even block is
$\operatorname{diag}(-1,b/2)$, while the degree-zero and degree-two
odd blocks are $\operatorname{diag}(-1,b/6)$ and
$\operatorname{diag}(b,0)$.  The corresponding gates cover them.
In transverse degree three, $P_3^Wm_0=rQ$ exactly.  Subtracting the
$Q$-component of the $E$ quadratic leaves a nonnegative multiple
$E_{22}\|P_{3,Q^\perp}^Wm_2\|_2^2$.  This covers the same component
of $\|\Pi V\|_2^2$, since $E_{22}\geq p_2^2$.
In every higher odd transverse degree, $E-pp^T\succeq0$ covers
the unused loss.  The other odd-block degrees use $R\succeq0$.

Apply the residual square completion to $V$ exactly as before:
$-\|\Pi V\|_2^2\leq-2\zeta\mathbb E[\mathcal RV]
+\zeta^2\mathbb E\mathcal R^2$, followed by the completion of
$-2\zeta\Delta V$.  This adds only $(\zeta^2/G)pp^T$ to the
even block and gives the displayed $D,c_6,\kappa$ terms.
Complete $-E_{22}t^2\leq2yt+y^2/E_{22}$ and the longitudinal
agreement penalty, and insert the exact linear and source identities.
For each transverse point, Bessel's inequality places the physical
four-moment vector in $\mathcal E$.
The identity $x^TNx=\sup_v(2v^TNx-v^TNv)$, valid also for
singular positive semidefinite $N$, and the four pointwise values of
$(h,k)$ now give precisely $\Phi_E$.
The reflections $z\mapsto-z$ and
$(q,v_0,v_2)\mapsto(-q,-v_0,-v_2)$ give the final symmetry claims.
\end{proof}

\paragraph{Transport and a necessary concavity allowance.}
Set $y(r)=y_0+E_{02}(r-r_0)$ and $v_y=y_0-E_{02}r_0$.
This freezes $b_0$ and hence the whole local certificate, while
\[
 (b-E_{00})r^2-2\gamma r+\frac{y(r)^2}{E_{22}}
 =a_r r^2-2\left(\gamma-\frac{E_{02}v_y}{E_{22}}\right)r
                +\frac{v_y^2}{E_{22}},\qquad
 a_r=b-E_{00}+\frac{E_{02}^2}{E_{22}}.
\]
Unlike the rank-one case, $a_r$ can be negative.  On a source cell
$r\in[r_-,r_+]$, endpoint interpolation remains an upper bound after
adding $\max(0,-a_r)(r_+-r_-)^2/4$.
No compact local certificate, or fixed-source price, alone proves a
global operator norm bound.

% Analytic extensions only. No new numerical norm certificate is asserted.
\begin{proposition}[Jensen reductions with a rank local oracle]
\label{prop:jensen-rank-reductions}
Retain the raw-Hermite rank residual-source notation
$p=(p_0,p_2)$, $L$, $B$, $\mathcal E_p$, $\Phi$, $b>0$ and $R$.
Let $0\le\omega<1$, $\nu^2=2/\pi$, and put $C_0=1-\nu^2$.
Use the exact source transport
\[
 y(r)=y_0+p_0p_2(r-r_0),\qquad v_y=y_0-p_0p_2r_0,
 \qquad \eta=\gamma-p_0v_y/p_2.
\]
Write $\mathcal C$ for the usual residual price, including
$6\beta^2/b$ and $v_y^2/p_2^2$, but excluding the amplitude and
source quadratics and $\mathbb E\Phi(Q)$.
Assume the three usual odd-block spectral gates for $R$.

For the ellipse reduction choose any $v_{\rm e}\ge\nu^2$ and assume only
\begin{equation}\label{eq:jensen-ellipse-even-gate}
 p_2^2\ge\frac b2\left(1+\frac{p_0^2}{1+v_{\rm e}\omega}\right).
\end{equation}
Then the ordinary rank local oracle, with constant longitudinal head,
gives
\begin{equation}\label{eq:jensen-ellipse-rank-price}
 \langle T_bf,g\rangle\le
 (1-\omega)u^2-2\ell u+(b-v_{\rm e}\omega)r^2-2\eta r
 +\mathcal C+v_{\rm e}\omega+\mathbb E\Phi(Q).
\end{equation}
The stronger old gate $p_2^2\ge(b/2)(1+p_0^2)$ is not required.

For the quartic reduction replace the even gate by
\begin{equation}\label{eq:jensen-quartic-even-gate}
 p_2^2\ge\frac b2\left(1+\frac{p_0^2}{1+\omega}\right).
\end{equation}
Define the modified local oracle by adding one independent bounded
dual variable $\tau\in[-1,1]$ to the rank oracle: replace $B(q,Z;v)$
inside $2\mathbb E_Z\max(|L|,|B|)$ by
$B(q,Z;v)+2\omega C_0\tau^3$, subtract $3\omega C_0\tau^4$
outside this expectation, and maximize over both $v\in\mathcal E_p$
and $\tau\in[-1,1]$. Denote the result by $\Phi_\omega(q)$.
Then
\begin{equation}\label{eq:jensen-quartic-rank-price}
 \langle T_bf,g\rangle\le
 (1-\omega)u^2-2\ell u+(b-\omega)r^2-2\eta r
 +\mathcal C+\omega\nu^2+\mathbb E\Phi_\omega(Q),
\end{equation}
and $0\le\Phi_\omega(q)-\Phi(q)\le\omega C_0$.
Neither extension introduces a quadratic dependence of the longitudinal
head on $q$. Existing residual curvature allowances remain necessary;
the quartic extension adds no new such allowance and its coarse tail
increases by at most $\omega C_0$.
\end{proposition}
\begin{proof}
Set $A(W)=\mathbb E_Z[Zh]$ and $m_0(W)=\mathbb E_Zk$.
Alignment of $P_1h=uZ$ gives $\mathbb EA=u$, and Jensen gives
$\mathbb EA^2\ge u^2$. The conditional profile budget proves
\[
 A^2+\nu^2m_0^2\le\nu^2,\qquad
 A^2+m_0^2-C_0m_0^4\le\nu^2.
\]
Thus the integrated quantities
\[
 v_{\rm e}-u^2-v_{\rm e}\mathbb Em_0^2,\qquad
 \nu^2-u^2-\mathbb Em_0^2+C_0\mathbb Em_0^4
\]
are nonnegative, using $v_{\rm e}\ge\nu^2$ for the first.

Apply the full even-block criterion with
$E=pp^T+a e_0e_0^T$, taking $a=v_{\rm e}\omega$ or $a=\omega$.
The gates $E\succeq0$ and $E-pp^T\succeq0$ are automatic.
The remaining gate is equivalent to
\[
 \det(E-\operatorname{diag}(-1,b/2))
 =(1+a)p_2^2-\frac b2(1+a+p_0^2)\ge0,
\]
which is exactly the respective displayed even gate. In particular
$p_2\ne0$. Add $\omega$ times the corresponding nonnegative
integrated budget to the full criterion. Its negative $m_0^2$ term
cancels the extra local even-block entry exactly, so the effective
local quadratic remains the original rank quadratic. The extra outer
terms are $\omega(v_{\rm e}-u^2-r^2v_{\rm e})$ in the ellipse case and
$\omega(\nu^2-u^2-r^2)$ in the quartic case. Exact $y(r)$ transport
then gives the two stated prices.

In the quartic case the additional local term is
$\omega C_0m_0^4$. Convexity of $x^4$ gives, for $|x|\le1$,
\[
 x^4=\sup_{|\tau|\le1}(4\tau^3x-3\tau^4),
\]
with equality at $\tau=x$. Since $m_0=\mathbb E_Zk$, the linear
part shifts the support field by $2\omega C_0\tau^3$, proving
the stated local formula by the same PSD duality as for the rank
oracle. The actual physical moment vector and $\tau=m_0$ are
admissible, and the two dual suprema may be interchanged with the
physical supremum. The added physical term lies between zero and
$\omega C_0$, proving the local gap and tail statements.
The extra dual terms do not depend on $q$, and $1-\omega>0$ preserves
the positive amplitude curvature needed for exact polygon pricing.
\end{proof}

These are analytic sufficient criteria. Any newly fitted parameters,
including ones violating the stronger old rank gate, require validation
of the true even-block gates and a complete local scalar certificate
before yielding an operator norm bound.

\begin{proposition}[A frozen-centering scalar majorant]
\label{prop:physical-r3-stein-majorant}
Let $U=\|\nabla_WQ\|^2/3-1$ and $\mathcal V=\mathbb EU^2$.
Fix any $\theta\in\mathbb R$, $\eta>0$, real $a_0,a_2$, and a
continuously differentiable test function $\varphi$ such that
$\varphi$ and $\varphi'$ have at most polynomial growth.
Set
\[
 e(q)=\varphi'(q)-a_0-a_2(q^2-1),\qquad
 \lambda(q)=-1-\theta(q^2-1),
\]
\[
 P(q)=\begin{cases}
 e(q)^2/(2\eta),&e(q)\geq\eta\lambda(q),\\
 \lambda(q)e(q)-\eta\lambda(q)^2/2,&e(q)<\eta\lambda(q).
 \end{cases}
\]
If, for every real $q$,
\begin{equation}\label{eq:physical-r3-scalar-majorant}
 \Phi(q)\leq A_0+A_2q^2+A_4q^4+A_6q^6+
 q\varphi(q)-[1+\theta(q^2-1)]\varphi'(q)-P(q),
\end{equation}
then
\begin{align}\label{eq:physical-r3-stein-price}
 \mathbb E\Phi(Q)\leq{}&A_0+A_2+(3+\kappa)A_4
 +(15+15\kappa+c_6)A_6\notag\\
 &+\frac\eta2[\mathcal V-2\theta\kappa/3+\theta^2(\kappa+2)]
 +a_2[\kappa/3-\theta(\kappa+2)].
\end{align}
In particular one may freeze $\theta=\kappa_0/[3(\kappa_0+2)]$
throughout a whole cumulant interval; the covariance correction in
\eqref{eq:physical-r3-stein-price} keeps its displayed positive sign.
When $\varphi$ is odd, the proposed majorant is even, so it suffices
to verify \eqref{eq:physical-r3-scalar-majorant} for $q\geq0$. The same
proof applies to the full-even-block potential $\Phi_E$ after replacing
$\Phi$ by $\Phi_E$.
\end{proposition}
\begin{proof}
Since $Q$ has pure degree three, Gaussian integration by parts gives
$\mathbb E[Q\varphi(Q)]=\mathbb E[(1+U)\varphi'(Q)]$.
It also gives $\mathbb EU=0$ and
$\mathbb E[U(Q^2-1)]=\kappa/3$; the latter follows by using
$\varphi(q)=q^3$.  Put $U_\theta=U-\theta(Q^2-1)$.
Then $U_\theta\geq\lambda(Q)$, its variance is
$\mathcal V-2\theta\kappa/3+\theta^2(\kappa+2)$, and
$\mathbb E[U_\theta(Q^2-1)]=\kappa/3-\theta(\kappa+2)$.
By construction $P(q)=\sup_{v\geq\lambda(q)}[ve(q)-\eta v^2/2]$.
Therefore $U_\theta e(Q)-P(Q)\leq\eta U_\theta^2/2$.
Take expectations in \eqref{eq:physical-r3-scalar-majorant} and apply
these identities.  This proves \eqref{eq:physical-r3-stein-price}.
\end{proof}

\paragraph{Exact transport of the source amplitude.}
Freeze a reference $r_0$ and $y_0$, and set
\[
 y(r)=y_0+p_0p_2(r-r_0),\qquad v_y=y_0/p_2-p_0r_0.
\]
Then $y(r)-p_0p_2r=y_0-p_0p_2r_0$, so the entire function $\Phi$ is
unchanged for every $r$.  Its outer quadratic simplifies exactly to
\begin{equation}\label{eq:physical-r3-transport}
 (b-p_0^2)r^2-2\gamma r+y(r)^2/p_2^2
 =br^2-2(\gamma-p_0v_y)r+v_y^2.
\end{equation}
Consequently no local amplitude error is needed for this transport;
both the $u$ and $r$ outer quadratics are convex and are bounded on
rectangles by their vertex values.  The source moments and cumulants
must still be covered separately.  A proposed scalar majorant becomes
a norm certificate only after it bounds \eqref{eq:physical-r3-local}
for the whole real $q$ line and its outer expectation is bounded for
every admissible physical source.  No finite grid maximization establishes
those requirements.  The case $r=0$ can be handled separately by a
source-free criterion, without normalizing a zero projection.

% Numerical corollary; uses the notation of r3_physical_dual_criterion.tex.
\begin{proposition}[A certified cubic-source rectangle]
\label{prop:r3-certified-rectangle}
For $b=18/25$, every odd Gaussian sign pair whose physical parameters satisfy
\[
 u\in[7549/10000,7551/10000],\qquad
 r\in[2574/10000,2577/10000],\qquad
 \kappa\in[51/100,55/100]
\]
obeys
\[
 \langle(P_1-bP_3)f,g\rangle\leq\frac{28167}{50000}
 =0.56334<\frac{40}{71}.
\]
The bound covers every admissible cubic tensor shape and every Gaussian
dimension.
\end{proposition}
\begin{proof}[Certificate proof]
The exact rational profile is frozen in
\texttt{frozen\_R3\_profile\_exact.json}.  The scalar majorant uses reserves
$2\cdot10^{-5}$ in its constant coefficient and $10^{-6}$ in its
degree-six coefficient.  Exact rational Bernstein derivative bounds and
directed $256$-bit MPFR arithmetic certify the whole interval $q\in[0,8]$:
the completed partition has $6751$ nodes and $3376$ accepted leaves.
A fresh independent traversal rechecked all accepted leaves and the
physical-center quadratic radii, with stable source hashes.
The independently replayed $[8,64]$ cover and exact nonnegative-coefficient
tail proof for $q\geq64$ remain valid after increasing the constant reserve.
Evenness covers negative $q$.

Use the fresh sixth-cumulant comparison and the exact cubic projection bound
\[
 c_6\leq\left[\frac{194}{\sqrt{10}}
              -\frac{633}{9}(z-2/5)\right]\kappa^{3/2},\qquad
 D\leq\kappa^2+\frac94(\kappa z)^{3/2},\quad z\in[2/5,1].
\]
In the outer price their respective coefficients are nonnegative.
After the exact $y(r)$ transport, the price is separately convex in
$u,r,\kappa,z$; thus its $16$ rectangle vertices suffice.
Independent directed evaluation gives the upper bound
$0.563333433603497\ldots<28167/50000$.
\end{proof}
This certifies the displayed rectangle of physical parameters.  It is not
a global bound for $\|P_1-(18/25)P_3\|_{\infty\to1}$.

% Uses the physical parameters and exact profile of r3_physical_dual_criterion.tex.
\begin{proposition}[Certified moving bands of cubic-source parameters]
\label{prop:r3-certified-moving-polygons}
Let $b=18/25$ and use the physical parameters $u,r,\kappa$ of an odd
Gaussian sign pair.  Define
\[
 t=r-\frac{5151}{20000}+3\left(u-\frac{151}{200}\right).
\]
In each of the following three cases, the pairing is bounded by the
listed cap, in every Gaussian dimension and for every admissible
dimensionless cubic tensor shape $z\in[2/5,1]$:
\[
\begin{array}{c|c|c|c}
 u& |t|\text{ at most}&\kappa&
 \langle(P_1-bP_3)f,g\rangle\text{ at most}\\ \hline
 {[.7538,.7562]}&.00025&[.3,.8]&28168/50000\\
 {[.75,.76]}&.001&[.3,.8]&451/800\\
 {[.7505,.7595]}&.001&[.15,1.1]&451/800
\end{array}
\]
All finite decimals in the table are exact rational numbers.  The first
cap is $0.56336<40/71$; the other two are $0.56375$.
\end{proposition}
\begin{proof}[Certificate proof]
Use the accepted whole-line scalar majorant from the exact frozen profile,
with constant reserve $1/50000$ and sixth-degree reserve $1/1000000$.
The transport $y(r)=y_0+p_0p_2(r-r_0)$ leaves this entire local majorant
unchanged.  Its outer amplitude part is
\[
 u^2-2\ell u+br^2-2(\gamma-p_0v_y)r+v_y^2,
 \qquad v_y=y_0/p_2-p_0r_0.
\]
After $r=5151/20000-3(u-151/200)+t$, its Hessian divided by two is
\[
 \begin{pmatrix}187/25&-54/25\\-54/25&18/25\end{pmatrix},
 \qquad \det=18/25>0.
\]
Thus it is jointly convex in $(u,t)$.

In the source-moment price the coefficients of the sixth cumulant and
the projection allowance are respectively
\[
 A=\frac{22262234067601129}{2500000000000000000000}>0,
 \qquad
 B=\frac{172316040932398871}{2500000000000000000000}>0.
\]
The fresh source bounds therefore give a price consisting of affine
terms in each source coordinate, plus
\[
 \left[A\left\{\frac{194}{\sqrt{10}}
       -\frac{633}{9}(z-2/5)\right\}
       +\frac94Bz^{3/2}\right]\kappa^{3/2}+B\kappa^2.
\]
The coefficient in braces is positive for $2/5\leq z\leq1$.
Consequently the price is separately convex in $\kappa$ and $z$.
Each affine parameter rectangle $(u,t,\kappa,z)$ is therefore bounded
by its $16$ vertex prices.  Independent exact-profile directed evaluation
at $256$ bits gives, respectively,
\[
 0.563356280368,\qquad 0.563648547036,\qquad 0.563686934863
\]
as rational upper bounds for the three vertex maxima.  These are strictly
below the respective table caps.
\end{proof}
The restrictions on $t$ are restrictions on the actual projection norm $r$.
These moving bands do not cover every $r$ at their listed amplitudes and
do not by themselves bound the global operator norm.

\begin{proposition}[Further certified source regions]
\label{prop:r3-additional-source-regions}
For $b=18/25$, use the physical parameters $u,r,\kappa$ of an odd
Gaussian sign pair.  In every Gaussian dimension and for every
admissible cubic tensor shape,
\begin{gather*}
 u\in[341/500,1381/2000],\quad 0\leq\kappa\leq1,
 \quad r\geq1/5\\
 \Longrightarrow\quad
 \langle T_bf,g\rangle\leq\frac{563371}{10^6}<\frac{40}{71}.
\end{gather*}
There are also the following full-even-block bounds.  Put
\[
 t=r-\frac{389}{1000}+\frac95\left(u-\frac7{10}\right).
\]
Under the common restriction $-1/200\leq t\leq1/50$, the bounds are
\[
\begin{array}{c|c|c}
 u&\kappa&\langle T_bf,g\rangle\text{ at most}\\\hline
 {[.69,.70]}&[0,1]&551/1000\\
 {[.70,.72]}&[0,1]&68/125\\
 {[.70,.72]}&[0,2]&451/800
\end{array}
\]
The first two table caps are below $40/71$; the last is $0.56375$.
Every finite decimal here denotes its exact rational value.
\end{proposition}
\begin{proof}[Certificate proof]
On $0\leq\kappa\leq1$, the accepted Stein row prices give
\[
 \mathbb E|Q|\leq\frac{830643632499}{10^{12}}
 <\rho_1:=\frac{207661}{250000}.
\]
The support--source comparison with this cap, $a=9/5$, and
$\lambda=0$ has directed right side below
$1.649920196798<1649921/10^6$.  Thus every such physical
source satisfies
\[
 r\leq\frac{1649921}{10^6}-\frac95u.
\]
The two accepted rank-one R3 profiles cover the source intervals
\[
 [1/5,199/500],\qquad[199/500,1649921/10^6-(9/5)u].
\]  Their whole-line local proofs use
constant reserve $1/50000$ and sixth-degree reserve $1/5000000$.
Independent replay accepts all $21873$ compact scalar leaves and the
exact tails starting at64.  Exact $y(r)$ transport and convex
endpoint pricing on the two source polygons give respective maxima
below $0.563249470254$ and $0.563370758117$, proving the first cap.
The complete replay is recorded in
\texttt{COMPLETE\_\allowbreak RESTRICTED\_\allowbreak DOMAIN\_\allowbreak REPLAY.json}.
The outer extension is recorded in
\texttt{independent\_\allowbreak finer\_\allowbreak amplitude\_\allowbreak extension.json}.

For the table, the accepted full-even profile uses constant reserve
$1/1000$ and sixth-degree reserve $1/10^6$.
Its independent whole-line replay checks $56194$ partition events,
$27636$ accepted leaves and $462$ exact physical-ellipsoid exclusions,
then glues the exact tail beyond64.  The six matrix gates and
$y(r)$ transport are exact.  Let
$a_r=b-E_{00}+E_{02}^2/E_{22}<0$.  On each $(u,t)$ rectangle,
product-corner interpolation incurs at most
\[
 (-a_r)\{(9/5)^2(\Delta u)^2+(\Delta t)^2\}/4.
\]
This follows because the mixed term is interpolated exactly and the
variance of the affine source is $(9/5)^2\operatorname{Var}(u)+
\operatorname{Var}(t)$; the positive $u^2$ term only improves the
comparison.  Source prices are separately convex in $\kappa$ and
the shared shape.  All48 vertex prices, including these reserves,
have independent directed maxima below
$0.550507802111$, $0.543476058131$, and $0.563687366986$.
They prove the table caps; the record is
\texttt{general\_r3\_source\_bands\_audit.json}.
\end{proof}
The first conclusion retains both $\kappa\leq1$ and $r\geq1/5$.
The table retains its restrictions on $t$ and $\kappa$.
Their union is a collection of physical-source regions, and does not
establish complete amplitude coverage or a global operator norm bound.

\begin{proposition}[Exact transport to a smaller cubic coefficient]
\label{prop:smaller-b-transport}
Freeze a valid physical Gaussian profile at $b_0>0$, including its
whole-line local and scalar majorants.  For every $0<b\leq b_0$ all
its positive semidefinite gates remain valid.  At the same physical
parameters $(u,r,\kappa,\text{shape})$, the resulting pairing upper
price changes exactly by
\[
 (b-b_0)r^2+c_\beta\left(\frac1b-\frac1{b_0}\right).
\]
Here $c_\beta=\beta^2$ for orthonormal longitudinal Hermites and
$c_\beta=6\beta^2$ for raw longitudinal Hermites in the R3 criterion.
No new local Gaussian integration is required.  In the first case
the local potential changes only by its constant
$\beta^2(1/b-1/b_0)$; in the second it is unchanged.
\end{proposition}
\begin{proof}
Each gate at the smaller coefficient is the old gate plus a positive
semidefinite diagonal matrix.  For example, the raw even gate gains
$\operatorname{diag}(0,(b_0-b)/2)$, and the raw odd gates gain
$\operatorname{diag}(0,(b_0-b)/6)$ or
$\operatorname{diag}(b_0-b,0)$.  The normalized gates behave likewise;
the independent gates $E\succeq pp^T$, $E\succeq0$, $R\succeq0$
do not change.  The rank-one raw gate also weakens.

The only remaining appearances of $b$ in the completed criterion are
the source term $br^2$ and the longitudinal square-completion constant.
The norm of the raw third Hermite is squared norm six, accounting for
the displayed factor.  All source residual and scalar majorant terms
are independent of $b$.  In the normalized native criterion the added
local constant cancels out of the centered Stein target, so its Stein
solution and all derivative inequalities remain the same.  Its added
constant is simply charged in the outer expectation.  This proves the
exact price identity.
\end{proof}

\paragraph{Uniform cell pricing and a global coefficient.}
Suppose an old valid uniform cap $U_0$ covers a physical source cell
$r\in[r_-,r_+]$ with $0\le r_-\le r_+$. Since $b-b_0\leq0$, the transported uniform cap
\[
 U_0+(b-b_0)r_-^2+c_\beta(1/b-1/b_0)
\]
is immediately valid.  For tighter endpoint repricing, the native
source curvature changes to $\xi(b)=\xi(b_0)+b-b_0$, and the full-even
R3 curvature changes to $a_r(b)=a_r(b_0)+b-b_0$.  Their respective
chord reserves must be updated to
\[
 (-\xi(b))_+(r_+-r_-)^2/4,
 \qquad (-a_r(b))_+(r_+-r_-)^2/4.
\]
Every such fixed-profile cell cap is convex in $b>0$.  A fixed
selection of profiles covering all physical cells therefore gives a
convex maximum of caps.  A norm certificate requires one coefficient
$b$ valid for the entire cover.  Minimizing separately for each sign
pair, at $b=\min\{b_0,\sqrt{c_\beta}/r\}$ when $r,c_\beta>0$,
does not prove a bound for a single operator.

\begin{theorem}[Complete amplitude branches at one coefficient]
\label{thm:same-b7171-regions}
Set $b_*=7171/10000$.  For every pair of odd Gaussian signs $f,g$,
in every finite Gaussian dimension, put
$u=\|P_1((f+g)/2)\|_2$.  Then
\[
 \langle f,T_{b_*}g\rangle\le
 \begin{cases}
 56338/10^5,&0\le u\le1381/2000,\\[2pt]
 563315250991/10^{12},&31/40\le u\le\sqrt{2/\pi}.
 \end{cases}
 \qquad <\frac{40}{71}.
\]
Both conclusions include every physical source amplitude, cumulant
and cubic tensor shape.
\end{theorem}
\begin{proof}[Certificate proof]
All the profiles here use the single coefficient $b_*$.
The elementary energy estimate with $\rho=847/1000$ gives
\[
 \langle f,T_{b_*}g\rangle
 \le u^2+b_*\rho^2
 \le\frac{7042043}{12500000}<\frac{40}{71}
 \quad(0\le u\le541/2500),
\]
where the displayed cap conservatively uses $b=18/25$.
Four previously replayed complete native frames, transported by
Proposition~\ref{prop:smaller-b-transport}, cover
$[.14,.42]$, $[.30,.485]$, $[.40,.54]$ and $[.50,.638]$.
Their outward caps are respectively below
$.561498051269$, $.562840841523$, $.563045009361$ and
$.563207127166$.  The intervals overlap each other and the elementary
branch.

The exact native-frame union on $[.60,.682]$ consists of $30268$
accepted cells and has maximum rational cap
$563067382411/10^{12}$.  On $[.682,.6905]$, $9661$ accepted native
cells and $41$ cells covered by the transported R3 profile have
maximum cap $563379835873/10^{12}<56338/10^5$.
The R3 replacement uses the independently replayed $11470+88$
compact leaves and exact tail beyond64, with constant and sixth-degree
reserves $1/50000$ and $1/5000000$.

For completeness of these source unions, the support--source proof
of \eqref{eq:Rrho} applies with the freshly proved $\rho=.847$.
For each amplitude subinterval $[u_-,u_+]$, its source grid covers
\[
 0\le r\le R_{.847}(u_-),
\]
an outward bound on the entire physical source range because
$R_{.847}$ is decreasing.  The cumulant intervals partition $[0,90]$.
Native prices are affine in the variance shape $s\in[0,1]$;
the R3 coordinate is exactly $z=2/5+(3/5)s\in[2/5,1]$.
Every cell has either one complete native bound or one same-coefficient
R3 bound.  Exact positive source chords and quadratic source maxima,
including interior vertices, cover its full continuous interior.
Independent price replays check $2240$ and $1344$ native endpoint
prices in these two unions, all selected profile hashes, and the
sixteen R3 replacement corners.  A separate exact verifier checks
all $112+42$ amplitude--cumulant groups and every gap-free source
partition.  This proves the first branch.

The upper native union on $[.775,.797885]$ has $5096$ cells and
maximum cap $563315250991/10^{12}$; its independent replay checks
$1568$ endpoint prices and the same exhaustive physical-source
partitions.  Finally $u\le\sqrt{2/\pi}<.797885$.
The complete union records are
\texttt{native\_\allowbreak lower\_\allowbreak base\_\allowbreak b7171\_\allowbreak target\_\allowbreak union.json},
\texttt{native\_\allowbreak R3\_\allowbreak b7171\_\allowbreak u682\_\allowbreak 6905\_\allowbreak union.json}, and
\texttt{native\_\allowbreak upper\_\allowbreak b7171\_\allowbreak target\_\allowbreak union.json}; their frozen
input hashes are checked in
\texttt{same\_\allowbreak b\_\allowbreak native\_\allowbreak R3\_\allowbreak glue\_\allowbreak audit.json} and the independent
native price replay ledgers.
\end{proof}
The uncovered amplitude range for this theorem alone is
$1381/2000<u<31/40$.  The theorem is an upper bound for the indicated
pairings, and does not by itself bound the unrestricted operator norm.

% Standalone analytic fragment. No numerical coverage is asserted here.
\begin{lemma}[Exact pricing on physical source polygons]
\label{lem:physical-source-polygon-price}
Fix an amplitude interval $[u_-,u_+]$, a source interval $[r_-,r_+]$,
and a cumulant interval. Suppose that, throughout the cumulant interval
and for every allowed shape, proved scalar inequalities imply
\[
 a_j u+r\le C_j,\qquad a_j>0\quad(1\le j\le m).
\]
Let $K$ be the intersection of these half-planes with the amplitude--source
rectangle, after any separately proved coordinate restrictions.
For rational inputs its vertices and the maximum on $K$ of
\[
 F(u,r)=u^2+L_u u+g r^2+L_r r+C
\]
are computable using rational arithmetic. For a nonempty $K$, it suffices
to check its vertices and, on each edge, the stationary point in the edge
interior when the restricted quadratic is concave. This includes segments
and singletons. An empty intersection contains no physical source.

For a native certificate interpolated in $r$, first split $K$ at every
source knot. On each resulting piece apply this rule to each of the four
cumulant--shape corner bounds, and take the maximum of the resulting
bounds. A fixed nonnegative contact multiplier may be included by changing
only $L_u,L_r,C$.
\end{lemma}
\begin{proof}
Every defining half-plane contains every physical source in the nominal
rectangle, so intersection and coordinate clipping preserve all physical
points. Rational line intersections give all vertices, retaining points
on the boundaries. If a full-dimensional polygon had a maximizing point
in its interior, sufficiently small horizontal displacements would stay
in the polygon. Their two objective values have average
$F(u,r)+\delta^2$, contradicting maximality. Thus its maximum lies on its
boundary. On the edge $(u,r)=(u_0,r_0)+t(\Delta u,\Delta r)$,
$0\le t\le1$, the quadratic coefficient is
$(\Delta u)^2+g(\Delta r)^2$. Endpoints suffice when this coefficient is
nonnegative; otherwise also check the unique stationary point if it lies
in $(0,1)$. All these operations are rational. The same calculation covers
the lower-dimensional cases.

A native source chord is valid only between successive saved knots.
Splitting at every knot therefore retains its validity on each closed
piece; shared boundaries have both valid endpoint bounds. For each fixed
cumulant--shape corner the chord changes only the linear and constant
coefficients in $r$. The already proved corner inequalities and positive
interpolation uniformly bound the whole cumulant--shape cell. Taking the
maximum over the four corner prices and all source pieces gives a uniform
bound. The scalar row chosen at a source endpoint is fixed throughout its
cumulant--shape cell; different endpoints may use different globally
valid scalar rows.

For a contact with parameters $a,B$, put $d^2=6B^2$ and let $A$ be its
uniform scalar price. The nonnegative multiplier $\lambda$ adds
\[
 -\lambda a u-\lambda r+
 \lambda\left(A+\frac{\beta d}{b}\right)
 +\frac{\lambda^2d^2}{4b}.
\]
Here $\beta$ is the normalized native Hermite coefficient; for a raw
Hermite R3 profile the product $\beta d$ is $6\beta_{\rm raw}B$.
The coefficient of $u^2$ remains one, so the same boundary computation
applies. The multiplier and scalar contact are fixed on the whole
cell. Selecting any one of several proved upper bounds on each cell
preserves validity.
\end{proof}

\paragraph{Uniform source planes.}
For the accepted even sextic scalar rows, the moment price is affine in
shape and convex in cumulant: the sextic coefficient is nonnegative and
the coefficient multiplying $\kappa^{3/2}$ is positive throughout the
shape interval. Hence the largest of the four endpoint prices is a
uniform $C_j$. A single fixed scalar row is used for each plane over the
whole cumulant interval. Taking the smallest uniform price among several
such rows is valid; taking separate minima at the four corners would
require a separate justification.


\begin{proposition}[Base scalar certificate at $.713$]
\label{prop:base-final29-713}
Set $\widehat b=713/1000$. Then
\[
 \N(\widehat b)\le\widehat B:=\frac{35245449}{62500000}
 =.563927184.
\]
Consequently the explicit reduced rational matrix sequence has actual limit
\[
 \lim_{n\to\infty}G(H_n(\widehat b))
 =\frac1{\N(\widehat b)}
 \ge\frac{62500000}{35245449}
 =1.7732785869744488\ldots .
\]
The rational reciprocal is an unconditional lower bound for
$G_{\mathbb R}$. The scalar certificate does not assert
$\N(\widehat b)=\widehat B$.
\end{proposition}
\begin{proof}[Certificate proof]
All branches use the same coefficient $\widehat b$. The independently
verified lower-amplitude ledger covers $u\in[0,1381/2000]$ with bound
$563295214999/10^{12}<\widehat B$. It combines the low-energy branch,
four transported regional frames and two complete native--R3 source
unions. The exact transfer from $\bar b=7171/10000$ is
\[
 (\widehat b-\bar b)r^2+
 \beta_{\rm norm}^2(1/\widehat b-1/\bar b),
\]
where $\beta_{\rm norm}^2=6\beta_{\rm raw}^2$ for an R3 profile.
The selected source unions retain the nonpositive first term; the
spectral gates remain valid at the smaller coefficient. The lower
ledger checks all $39929$ native cells and $41$ R3 cells in these two
unions, along with the four earlier regional frames.

The original upper cover has $25$ amplitude bands on
$[1381/2000,159577/200000]$ and $56$ cumulant bands on $[0,90]$.
The endpoint $159577/200000$ exceeds the maximum physical amplitude
$\sqrt{2/\pi}$. Every one of its $1400$ groups covers every physical
source from zero to an outward $R_{.847}(u_-)$ endpoint, with all shapes
included. Native shape $s\in[0,1]$ corresponds exactly to
$z=2/5+(3/5)s$ for the R3 prices. Directed checks verify all $25$
original source-grid upper endpoints and $250$ inverse clips.
Nineteen initial R3 profiles give $1064$ cumulant-band constants,
checked at all $4256$ cumulant--shape corners. Seven native tables
have $11648$ independently checked baseline endpoint prices, followed
by the exact transfer above. Their source quadratics also gain
$\widehat b-\bar b$.

Each cumulant band has ten proved uniform positive-head source planes.
The $560$ planes and their moment corner prices are independently
verified. Of the original $255360$ source cells, $253282$ retain
upper prices and $2078$ are exactly empty. The producer-free exact
replay checks every original interval, selected polygon or containing
rectangle, native radial knot and concave edge stationary point, using
Lemma~\ref{lem:physical-source-polygon-price}. It reconstructs all
selected contact and ordinary prices and accounts for every original
cell exactly once.

Three additional independently accepted whole-line profiles strengthen
$1069$ prices on the identical cover. An exact amplitude refinement
then replaces $2814$ difficult parent cells by $27340$ child intervals:
$20749$ are nonempty and $6591$ are exactly empty. The other $250468$
price cells remain unchanged. Each child inherits its parent's proved
source restrictions; its clipped polygon and selected price are checked
independently. The old parent bound is retained as an option. Five
further accepted whole-line profiles then strengthen prices on this
same refined cover in three successive monotone stages, changing
$2397$, $27889$ and $2428$ prices respectively. Their $224$, $448$ and
$448$ new cumulant--shape corner prices are independently recomputed.
The resulting $27$-profile cover has $271217$ price cells and $8669$
exact empty cells.

Finally, introduce a $68$-point cumulant grid while retaining every
previous bound. Only $364$ difficult parents are split, giving an exact
partition into $1794$ cumulant children. Of these, $1726$ retain prices
and $68$ are exactly empty; the remaining $270853$ price cells are
unchanged. Seven native endpoint tables on this finer grid have $13936$
independently checked baseline corner prices, transported exactly to
$\widehat b$. Fourteen new cumulant bands use ten fixed positive-head
planes each, chosen from the already proved $60$-row analytic bank.
All $140$ planes are independently checked at their $560$ moment
corners. Their intersections contain every physical child point.
An exact replay verifies every child cumulant partition and complete
clipped polygon. It reconstructs $1110$ newly selected native prices,
splitting every radial knot and including every concave edge stationary
point; the other $616$ nonempty children retain their valid parent bounds.
No new R3 constant is needed in this last refinement.

Two final independently accepted whole-line profiles are added on
this identical cumulant-refined cover. They strengthen $691$ and $343$
prices respectively, while retaining every prior bound. Each new profile
has $280$ independently recomputed cumulant--shape corner prices on
the mixed collection of original and refined bands. An exact replay
reconstructs every changed polygonal or rectangular price and verifies
that all other cell records and all empty cells are unchanged.

The final complete upper cover thus contains $272579$ price cells and
$8737$ exact empty cells. Its maximum recorded price is
\[
 \frac{4511417459}{8\cdot10^9}<\widehat B.
\]
\begingroup\raggedright
The frozen source is
\texttt{global\_\allowbreak upper\_\allowbreak b713\_\allowbreak native\_\allowbreak kappa68\_\allowbreak twentyninebank\_\allowbreak envelope.json.gz},
with SHA-256
\texttt{ea763375a19e63a49174bcf18975c948ab188f993edfb471014cbb2bfec04d2d}.
The final independent reconstruction is
\texttt{fullupper\_\allowbreak b713\_\allowbreak native\_\allowbreak kappa68\_\allowbreak twentyninebank\_\allowbreak monotone\_\allowbreak audit.json};
it binds the complete chain of accepted parent covers and replay hashes.
\par\endgroup
These upper and lower branches therefore bound every odd sign pairing
in every finite Gaussian dimension. Negating one sign gives the absolute
norm bound. The established Gaussian-to-matrix convergence and exact
ratio-preserving reductions give the actual matrix limit and lower bound.
\end{proof}

\begin{theorem}[Complete transported scalar certificate]
\label{thm:final-preferred-polygon}
Set $\widehat b=711/1000$. Then
\[
 \N(\widehat b)\le\widehat B:=\frac{112782853}{2\cdot10^8}
 =.563914265.
\]
Consequently the explicit reduced rational matrix sequence has actual limit
\[
 \lim_{n\to\infty}G(H_n(\widehat b))
 =\frac1{\N(\widehat b)}
 \ge\frac{2\cdot10^8}{112782853}
 =1.7733192119195636\ldots .
\]
The rational reciprocal is an unconditional lower bound for
$G_{\mathbb R}$; the norm need not equal its certified cap.
\end{theorem}
\begin{proof}[Coefficient transport]
Retain the complete upper physical partition and every selected frozen
criterion from Proposition~\ref{prop:base-final29-713}, at
$b_0=713/1000$. Amplitude, source, cumulant, shape and all physical
source restrictions are independent of the operator coefficient.
All $8737$ empty cells therefore remain empty.

For an ordinary frozen criterion let $\beta_i$ be its normalized
longitudinal coefficient: it is the recorded native coefficient, or
$\sqrt6$ times the recorded raw R3 coefficient. Its pointwise price
changes by
\[
 (\widehat b-b_0)r^2+
 \beta_i^2(1/\widehat b-1/b_0).
\]
For a fixed head-contact criterion, with multiplier $\lambda_i\ge0$
and head $a_cZ+B_cH_3(Z)$, put $d_c=\sqrt6B_c$. Completing the same
longitudinal square gives instead
\[
 (\widehat b-b_0)r^2+
 (\beta_i+\lambda_i d_c/2)^2
       (1/\widehat b-1/b_0).
\]
Thus the ordinary case is included by taking $\lambda_i=0$.
These identities follow by grouping the longitudinal cost
$\beta_i^2/b$, its contact cross term
$\lambda_i\beta_i d_c/b$ and the contact square cost
$\lambda_i^2d_c^2/(4b)$. All other frozen prices are unchanged.
All spectral matrix gates are checked at the smaller coefficient.

Let $a_i$ be the accepted old upper price and $r_{i,-}$ its source
cell's lower endpoint. Since $\widehat b-b_0<0$ and every physical
point in that cell satisfies $r\ge r_{i,-}\ge0$, a valid new uniform
price is
\[
 a_i+(\widehat b-b_0)r_{i,-}^2+
 (\beta_i+\lambda_i d_c/2)^2(1/\widehat b-1/b_0).
\]
An inherited parent bound is treated in precisely the same way, using
its original frozen criterion and multiplier; it already bounds every
physical child point. The coefficient $\widehat b$ is fixed throughout
all cells and is not chosen separately for different sign pairs.

An independent producer-free replay reconstructs all $43$ exact
profile bindings, verifies the target spectral gates and evaluates
these shifts on every one of the $272579$ price cells. The final assigned mosaic has $272557$ ordinary cells and $22$ R3
contact cells. Every one of their transport shifts is rational and is
computed exactly. The maximum rounded shifted price is
\[
 \frac{563914263671}{10^{12}}<\widehat B.
\]
The source is the hash-bound
\texttt{assigned\_\allowbreak twenty-nine\_\allowbreak coefficient\_\allowbreak scan.json};
the exhaustive independent cell replay is
\texttt{assigned\_\allowbreak b711\_\allowbreak shifted\_\allowbreak cells\_\allowbreak audit.json}.
The separately verified complete lower-amplitude branch at the same
coefficient, $u\in[0,1381/2000]$, has upper bound
$563307766601/10^{12}<\widehat B$. It includes the elementary branch,
four earlier regional frames and both complete lower source unions.
Together the transported upper and lower covers contain every physical
odd sign pair in every finite Gaussian dimension. Negating one sign
gives the absolute norm bound. Gaussian-to-matrix convergence and the
exact ratio-preserving reductions give the stated actual matrix limit.
\end{proof}

\section{Consequences for search and reflection representations}
The sequence lies in the full absolute doubly stochastic class. Formula \eqref{eq:final-sequence}
now performs that conversion. Evaluate the copies on their smaller seed,
merge proportional rows and columns with their total absolute weights,
and quotient signed row/column permutations.

At fixed order, decompose $|B|$ into a convex combination of permutation
matrices and restore the fixed entry signs, giving signed permutations
with simplex weights and no cancellation. A signed left permutation
fixes the first term to the identity. Disconnected support ratios are
$C$-weighted component averages, so keep a best component. A balanced
cycle-sign support gauges to a nonnegative matrix and has ratio one.
For two signed permutations with weights $t,1-t$, each nontrivial
component is an alternating cycle on $k\ge2$ rows and $k$ columns.
A frustrated cycle has $C=k-2\min(t,1-t)$ and
\[
 V\le k\sqrt{1-4t(1-t)\sin^2(\pi/(2k))}.
\]
To see the latter bound, gauge its matrix to $tI+(1-t)S$, where the
signed cyclic permutation has $S^k=-I$, and use $V\le k\|\cdot\|_{\rm op}$.
For $0\le t\le1/2$ the squared ratio is increasing, since
$k\sin^2(\pi/(2k))\le1$. Its maximum is at most
$k\cos(\pi/(2k))/(k-1)$, which is $\sqrt2$ at $k=2$ and at most
$3/2$ for $k\ge3$. Thus a search near $1.775$ needs frustrated cycles
and at least three permutation terms; this does not assert that three suffice.

A weighted-reflection representative is also explicit. Since $|H_n|$
is doubly stochastic, $\norm{H_n}_{\mathrm{op}}\le1$. If its order is
$M_n$, put
\[
 U_n=\begin{pmatrix}H_n&(I-H_n^2)^{1/2}\\
                     (I-H_n^2)^{1/2}&-H_n\end{pmatrix},\quad
 P_n=(I-U_n)/2,\quad D_n=\operatorname{diag}((2/M_n)I_{M_n},0).
\]
Then $P_n=P_n^T=P_n^2$, its diagonal is $1/2$, and
\[
 D_n-2D_n^{1/2}P_nD_n^{1/2}=(2/M_n)H_n\oplus0.
\]
Its ratio is unchanged. Zero auxiliary weights are essential to this
universal representation. Strictly equal positive weights on every
coordinate would give an orthogonal matrix, whose ratio is at most
$\pi/2$, and would wrongly exclude the target. Indeed, for an orthogonal matrix $O$ of order $m$, $V(O)\le m$.
Gaussian signs of $Z_i$ and $(O^TZ)_j$ give
\[
 C(O)\ge\frac2\pi\sum_{ij}O_{ij}\arcsin(O_{ij})
 \ge\frac2\pi\sum_{ij}O_{ij}^2=\frac{2m}\pi,
\]
so $G(O)\le\pi/2$.

\begin{thebibliography}{99}
\bibitem{peng} B. Peng, \emph{A computer-assisted upper bound of 1.7813 for the real Grothendieck constant}, arXiv:2609.20074v1, September 17, 2026.

\bibitem{saha} R. Saha, A. Li, A. Xue, S. Chaudhuri, A. Klivans,
P. K. Kothari and R. Meka, \emph{New Lower and Upper Bounds for the
Grothendieck Constant}, arXiv:2608.11158v2, 2026.
\bibitem{baseline} Retained project report,
\emph{Real Grothendieck continuation}, September 30, 2026,
Theorem 3 and its inherited analytic premises.
\bibitem{maps} Retained project report,
\emph{Matrix invariance and stochastic normalization}, September 30,
2026; independently reconstructed reductions in this continuation.
\bibitem{history} Recovered matrix-mining sign certificates and
 \emph{Baseline and transfer review}, October 1, 2026, included in the
source package. The fixed-coefficient sign witnesses have fresh directed numerical replays;
the historical global norm cover remains a separate premise.
\bibitem{songnote} \emph{Grothendieck 177128}, preceding delivered
note and source package, October 1, 2026. Its endpoint depends on the
cited Song--Yue uniform norm theorem and the stated nonlinear-cover scope.
\bibitem{HJM} S. Heilman, C. Jones and G. Malavolta,
\emph{Sharper Bounds for the Complex Grothendieck Constant},
arXiv:2609.25462v1 (2026), equation (8.12). The tensor identities used
here are also freshly checked by connected Wick diagrams.
\end{thebibliography}

\appendix
\section{Exact tensor comparison}\label{app:tensor}
We give the finite acceptance rules for \eqref{eq:633scalar}.
Set $x=1-y$ and $z=1-a$. Each of 84 leaves specifies a closed rational
rectangle in $[0,1]^2$, one of the upper bounds \eqref{eq:J}--\eqref{eq:M},
and rational tangent parameters. For $t>0$ the inequalities
\[
 \sqrt X\le(X+t^2)/(2t),\qquad
 \sqrt{XY}\le(tX+Y/t)/2
\]
replace square roots by polynomial upper bounds. For a cube $X^{3/2}$
multiply the first inequality by $X\ge0$.
Choose rational $y_0\in[0,1]$ and $l_0>0$ with
$l_0^2\ge1+3y_0^2/2$. Cauchy--Schwarz gives
\[
 \sqrt d\ge L(y)=(1+3y_0y/2)/l_0.
\]
The bracket $[(K+422/15)d-211/3]$ is positive, since its minimum
is $K-211/5>0$. It is enough to show that the selected polynomial upper
bound for $\mathcal M$ or $\mathcal J$ is at most this bracket times
$L(y)$.

The dedicated equality-corner rectangle is
$x,z\in[0,2^{-19}]$. It uses $(p,q)=(3/4,-1)$ and
$L(y)=(2+3y)/\sqrt{10}$. Here
$(1-y^2)^{3/2}\le(1-y^2)/512$; tangent parameters $7/4$ and $1$
handle the remaining roots. This retains equality at the corner.

Each difference polynomial has degree at most $(4,4)$ over
$\mathbb Q(\sqrt{10})$. After substitution of the affine rectangle
coordinates, let $c_{kl}$ be its power coefficients. Its Bernstein
coefficients are
\[
 b_{ij}=\sum_{k\le i,\,l\le j}c_{kl}
 \frac{\binom ik}{\binom4k}\frac{\binom jl}{\binom4l},
 \qquad 0\le i,j\le4.
\]
Every coefficient is nonnegative. The sign of $r+s\sqrt{10}$ is
decided by rational sign cases and comparisons of $r^2$ with $10s^2$.
Since the Bernstein basis is nonnegative and sums to one, the 25 checks
prove the entire rectangle. The leaf counts are 21, 36 and 7 for the
first three $\mathcal J$ bounds, 1 and 18 for the two Gram branches,
and 1 for the dedicated corner: 84 rectangles and 2100 coefficients.

The rectangles lie within the square, have pairwise disjoint interiors,
and have exact total area one. Their finite closed union covers the square:
an omitted point would have a relatively open neighborhood of positive
area. This checks every boundary as well as the interior. Two independent
programs reconstruct the polynomials from the formulas. Neither imports
the producer's coefficient lists or the older 99-rectangle data.

\section{Whole-line scalar certificates}\label{app:stein}
All parameters of \eqref{eq:regression} and \eqref{eq:stein-majorant}
are frozen in exact rational row files. The rows and complete cumulant
cover are shown below; displayed prices are rounded upward, while
acceptance uses directed rational enclosures.
\begin{center}\small
\begin{tabular}{lll}
Row & Cumulant interval & Corner price upper\\\hline
$0$ & $[0,39/100]$ & $.825942224105$\\
$2$ & $[39/100,39/25]$ & $.839598147439$\\
$3$ & $[39/25,4]$ & $.847061071421$\\
$4$ & $[4,34/5]$ & $.847358899934$\\
$k8$ & $[34/5,48/5]$ & $.848520760642$\\
$5$ & $[48/5,76/5]$ & $.848507934475$\\
$6$ & $[76/5,28]$ & $.846999761374$\\
$7$ & $[28,90]$ & $.838404593017$\\\hline
\end{tabular}
\end{center}
For each row, 512-bit outward Arb arithmetic and fourth-order Taylor
bounds verify the entire compact interval $[0,64]$ on an adaptive rational
partition. Mills' ratio is enclosed at the endpoints, and $R'=xR-1$
generates the needed Taylor coefficients.
For $x\ge64$, $0\le1-xR(x)\le x^{-2}$ and convexity of the square
reduce the comparison to two rational polynomial tests. Clearing the
positive denominators and substituting $x=64+t$ gives nonnegative
coefficients and positive constants, proving the complete unbounded tail.
Evenness supplies the negative half-line.

The independent verifier uses direct Chebyshev recurrence, 768-bit
arithmetic and sixth-order Taylor bounds on $[0,80]$, followed by a
separately reconstructed exact tail at 80. It also independently assembles
all 32 prices in \eqref{eq:cornerprices}. A further checker reconstructs
the 16 tail comparisons and 32 corner prices without importing either
scalar implementation. Complete comparisons, rather than sampled maxima,
are the acceptance criterion.

\section{Polynomial certificate for the cubic source gate}\label{app:cubic-source}
The third numerator in Proposition~\ref{prop:cubic-source} is explicitly
\begin{align*}
 P_3={}&A^6/8+A^5B(3w/2+21/4)\\
 &+A^4B^2(57w^2/8+195w/4+84)\\
 &+A^3B^3(17w^3+339w^2/2+540w+504)\\
 &+A^2B^4(171w^4/8+549w^3/2+1194w^2+1728w)\\
 &+AB^5(27w^5/2+837w^4/4+1116w^3+1980w^2)\\
 &+B^6(27w^6/8+243w^5/4+378w^4+756w^3).
\end{align*}
Every coefficient is nonnegative for $A>0$, $B,w\ge0$.
Thus this finite identity establishes the derivative sign on the entire
half-line; no numerical Gaussian quadrature is used.
As an exact example, take $A=97/40$, $B=1/6$. Then
\[
 a=117/40,\quad d=\sqrt6/6,\quad
 (q_0,q_2,q_4,q_6)=
 \left(\frac{331}{120},\frac{20}{97},
 -\frac{1096000}{265587843},
 \frac{106274560000}{727184560303017}\right).
\]
At $\kappa=53/100$, $u=151/200$, $v=-1/4$, the two variance
endpoints give the directed source caps
\[
 r<252844293/10^9 \quad(V=7\kappa/45),\qquad
 r<250945617/10^9 \quad(V=2\kappa/9).
\]
The same affine shape interpolation applies between them.
These conditional examples illustrate the gate; $v=-1/4$ is not assumed
for arbitrary optimizing signs. Independent exact polynomial and directed
price implementations verify the recurrence, integration constants and
two rational caps.

\section{Complete continuous regional certificates}\label{app:regional}
For $T=P_1-(18/25)P_3$, odd signs $f,g$, and
$u=\|P_1((f+g)/2)\|_2$, the following complete certificates give
the displayed bounds on $\langle f,Tg\rangle$:
\begin{center}
\begin{tabular}{c|r|r|c}
amplitude interval & local boxes & scalar rows & rational upper cap\\ \hline
$[7/50,21/50]$ & $63681$ & $24$ & $561523/10^6$\\
$[3/10,97/200]$ & $61829$ & $24$ & $562857/10^6$\\
$[2/5,27/50]$ & $58841$ & $24$ & $563039/10^6$\\
$[1/2,319/500]$ & $151929$ & $24$ & $563266/10^6$\\
$[3/5,331/500]$ & $153593$ & $24$ & $563302/10^6$\\
$[331/500,341/500]$ & $627497$ & $24$ & $563367/10^6$\\
$[77/100,159577/200000]$ & $53541$ & $24$ & $563334/10^6$
\end{tabular}
\end{center}
Every cap is strictly below $40/71$.
The respective directed totals are below
$0.561522410$,
$0.562856131$, $0.563038808$, $0.563265432$,
$0.563301723$, $0.563366577$, and $0.563333793$.
The first profile uses the universal transverse bound $r\le849/1000$
and therefore remains valid throughout its extended amplitude interval.
The other extensions are checked against the outward source endpoint.
The outer polynomial $u^2-2\ell u+C$ is convex, so rational endpoint
checks cover each entire interval.

Each certificate has an optimizer-free, directed-arithmetic replay of
all physical moment boxes, scalar compact leaves, infinite scalar tails,
the five matrix gates, and the continuous cumulant, shape, and source
interpolations. Infinite tails explicitly begin beyond every target and
regression hinge. Local height and scalar coefficient repairs are fully
priced. These are regional norm bounds; amplitudes outside the stated
coverage still require certification before a new unrestricted norm cap
or Grothendieck lower endpoint follows.
The primary union is $[.14,.682]\cup[.77,\sqrt{2/\pi}]$.
The separately proved low-energy branch covers $[0,.14]$.
Thus these earlier regional certificates alone leave the target amplitude
interval $(.682,.77)$ unclosed; the later complete cover is treated separately.
The earlier, tighter bound $9/16$ on $[.79,\sqrt{2/\pi}]$ and its
original portable certificate are retained as supplementary results.

\section{Original-degree accounting for the sign witness}\label{app:witness}
In this appendix $H_j=\operatorname{He}_j/\sqrt{j!}$ denotes the
orthonormal longitudinal Hermite polynomial.
The central limit theorem for $Y_N$ alone does not justify assigning it
ordinary Hermite degree one. Each summand has original degree three.
Here we verify all projections used in \eqref{eq:new-witness}.

First replace the sign by a bounded smooth approximation with bounded
derivatives. Expand in one transverse summand $H_3(X_i)/\sqrt N$.
Orthogonality and parity give
\[
 \E[H_1H_3]=\E[H_1H_3^2]=\E[H_2H_3]=0.
\]
Taylor expansion therefore makes the $H_1(X_i)$ coefficient
$O(N^{-3/2})$ and the $H_2(X_i)$ coefficient $O(N^{-1})$.
The $H_3(X_i)$ coefficient is
$N^{-1/2}\E\partial_y f+O(N^{-3/2})$.
Expand separately in distinct transverse coordinates for mixed terms.
The complete list of original-degree-one and degree-three basis types is
\begin{center}\small
\begin{tabular}{lcc}
Basis type & Coefficient order & Summed squared energy\\\hline
$H_1(Z)$ & $\to m_1$ & retained\\
$H_1(X_i)$ & $O(N^{-3/2})$ & $O(N^{-2})$\\
$H_3(Z)$ & $\to m_3$ & retained\\
$H_2(Z)H_1(X_i)$ & $O(N^{-3/2})$ & $O(N^{-2})$\\
$H_1(Z)H_2(X_i)$ & $O(N^{-1})$ & $O(N^{-1})$\\
$H_1(Z)H_1(X_i)H_1(X_j)$ & $O(N^{-3})$ & $O(N^{-4})$\\
$H_3(X_i)$ & $N^{-1/2}m_Y+o(N^{-1/2})$ & $\to m_Y^2$\\
$H_2(X_i)H_1(X_j)$ & $O(N^{-5/2})$ & $O(N^{-3})$\\
$H_1(X_i)H_1(X_j)H_1(X_k)$ & $O(N^{-9/2})$ & $O(N^{-6})$\\\hline
\end{tabular}
\end{center}
The indices $i,j,k$ are distinct. For a fixed smooth approximation the
orders are uniform in these indices; the last column counts every term
of the displayed type. The ordinary central limit theorem and normal
integration by parts identify the three retained moments.
Reflection of $Y_N$ preserves $m_1,m_3$ and changes $m_Y$ to $-m_Y$.
Thus the degree-one and degree-three cross-pairings tend respectively
to $m_1^2$ and $m_3^2-m_Y^2$.

Finally $q'(z)\ge511/200>0$. Conditional on $Y_N$, the normal density
and this derivative bound give a small-band estimate
\[
 \Pr(|Y_N+q(Z)|\le\varepsilon)
 \le\frac{2\varepsilon}{(511/200)\sqrt{2\pi}},
\]
uniformly in $N$. The same holds for the reflected witness.
Consequently the smooth approximation can be removed uniformly in
$L_2$; orthogonal projection is an $L_2$ contraction. This completes
the limiting degree calculation. The exact witness parameters,
directed compact integration and infinite-tail bounds are supplied in
\texttt{reduced\_search/operator/cubic720/README.md} and the adjacent
\texttt{cubic720\_norm\_witness.json}.

\section{Reproduction, retained premises and proof status}\label{app:status}
The commands and numerical tables in this appendix describe the retained
October 1 checkpoints. The completed continuation is stated in
Appendix~\ref{sec:new-final-goal}, where the complete same-coefficient
certificate proves $\GR\geq1.775$ and the full proof gives the
three-significant-digit conclusion $1.78$.

The source archive contains the exact row inputs, tensor rectangles,
complete-halfline checkers, formal Wick enumeration, implicit matrix
construction, independent audits, and directed MPFR numerator and sign
witness calculations. The main checks can be run from the archive root:
\begin{verbatim}
python reduced_search/energy/633/certify_wick_identity.py
python reduced_search/energy/633/discover_cover.py --replay
python reduced_search/energy/633/independent_replay.py
python reduced_search/energy/certify_stein_cubic_L1.py
python reduced_search/energy/audit_stein_L1_independent.py
python tonight/moments/certify_cubic_head_polynomial.py
python tonight/theory/certify_numerator_rate.py
python tonight/theory/certify_retained_finite_ratio.py
python tonight/theory/certify_saha_benchmark.py
python tonight/theory/certify_published_upper_rounding.py
python tonight/operator/cubic_L1_refinement/refit847/verify_cover.py
python tonight/audit/refit847/replay_complete_refit_cover.py
python tonight/regional_certificates/u014_042/code/verify_region.py
python tonight/regional_certificates/u030_0485/code/verify_region.py
python tonight/regional_certificates/u040_054/code/verify_region.py
python tonight/regional_certificates/u050_0638/code/verify_region.py
python tonight/regional_certificates/u060_0662/code/verify_region.py
python tonight/regional_certificates/u0662_0682/code/verify_region.py
python tonight/regional_certificates/u077_c/code/verify_region.py
python tonight/operator/projection/replay_R3_cell/verify_all.py
python tonight/operator/projection/replay_accepted_R3/verify_all.py --quick
python tonight/operator/preferred_global/verify_bundle.py
python reduced_search/operator/cubic7171/certify_cubic7171_norm_witness.py
python tonight/audit/replay_cubic7171_witness.py
python tonight/audit/replay_selected_cubic_witnesses.py
\end{verbatim}
Python, python-flint and the system MPFR library suffice for these accepted
checks; discovery programs have additional numerical dependencies.
Files are bound by a SHA-256 manifest. Floating optimization identifies
candidate inputs and is kept separate from proof outputs.

For $b=.72$, directed evaluations of \eqref{eq:vector-bound} give
\begin{center}\small
\begin{tabular}{rr}
$n$ & Certified vector numerator lower bound\\\hline
32 & $.999159371292$\\
128 & $.999947506710$\\
1024 & $.999999179839$\\\hline
\end{tabular}
\end{center}
These bounds are unconditional numerator estimates. A denominator bound
valid for all sign pairs in every dimension is still required before
\eqref{eq:finite-ratio} gives the target lower endpoint.

The global cap in Theorem~\ref{thm:final-preferred-polygon}, inserted
into \eqref{eq:finite-ratio} at $b=711/1000$, gives
\begin{center}\small
\begin{tabular}{rr}
$n$ & Unconditional lower bound for $G(H_n(.711))$\\\hline
32 & $1.771831610393$\\
128 & $1.773226593281$\\
1024 & $1.773317765118$\\\hline
\end{tabular}
\end{center}
These displayed lower bounds converge to $200000000/112782853$.
The actual ratio converges to $1/\N(.711)$, which may be larger.
All finite rounding losses, including the total entry error $n^{-4}$,
are charged by the directed checker
\texttt{evaluate\_\allowbreak finite\_\allowbreak ratio.py}.

For comparison, the retained $.699$ denominator premise already gives
explicit finite members of the reduced sequence. Directed evaluation of
\eqref{eq:finite-ratio}, with that premise stated separately, yields
\begin{center}\small
\begin{tabular}{rr}
$n$ & Conditional lower bound for $G(H_n(.699))$\\\hline
32 & $1.771523050434$\\
128 & $1.772908047308$\\
1024 & $1.772998563604$\\\hline
\end{tabular}
\end{center}
These conditional finite lower bounds converge quantitatively to
$1.773$; the actual ratio limit remains $1/\N(.699)$. These numbers
use the earlier scalar norm premise and do not reprove it. The very large matrices are specified implicitly
by their analytic seed and exact cloning maps.

The inherited project estimate \eqref{eq:baseline} remains an
explicit premise. Its 55-interval cover and inherited integrations are
not reconstructed by the present archive. The exact reductions,
asymptotic sequence, coupled cumulant theorem, cubic absolute-moment
theorem, and support--source propositions above have their own fresh
proofs and do not depend on that missing reconstruction.

The conditional project lower endpoint is $1.773$, at distance $.002$ from
$1.775$. The later $1.7712846852\ldots$ family was weaker and does not
replace it. The fixed $.699$ construction has a directed ceiling
$1.77440101857\ldots$; the fresh $.72$ ceiling is
$1.77543615103\ldots$, and the fresh $.7171$ ceiling is
$1.775462496412$. These ceilings bound the reciprocal-norm methods,
not the unrestricted Grothendieck constant. The regional amplitude
caps and the complete same-coefficient global envelope prove the fresh
October 1 unconditional lower endpoint $1.7733192119\ldots$.
The refined complete certificate in Section~\ref{sec:new-complete-refined-certificate}
improves that checkpoint using independently accepted source prices and
a complete refined cover. The final certificate in
Appendix~\ref{sec:new-final-goal} reaches $1.775$ and completes the
three-significant-digit determination.

% New-round theorem fragment. No numerical all-sign norm bound is asserted.
% Requires amsmath, amssymb, amsthm and proposition/proof environments.
\section{Finite odd multipliers with an identity tail}
Let $P_j^{(d)}$ be the orthogonal projection onto total Gaussian Hermite
degree $j$, and let $P_{\rm odd}^{(d)}$ be antipodal odd projection.
Fix computable real coefficients $\tau,c_0,\ldots,c_K$ and put
\[
 T^{(d)}=\tau P_{\rm odd}^{(d)}+
                 \sum_{k=0}^Kc_kP_{2k+1}^{(d)},\qquad
 \lambda_{2k+1}=\tau+c_k\ (k\leq K).
\]
Its remaining odd eigenvalues are $\tau$.  Suppose
\begin{equation}\label{eq:new-odd-spectral-cube}
 \tau+c_0=1,\qquad |\tau|\leq1,\qquad
                    |\tau+c_k|\leq1\quad(1\leq k\leq K).
\end{equation}
Define
\[
 N(T)=\sup_{d\geq1}\sup_{|f|,|g|\leq1}
             \langle f,T^{(d)}g\rangle.
\]
The absolute value may equivalently be inserted in this definition.
Odd signs suffice: odd projection preserves the unit ball of $L_\infty$,
and a bilinear functional over bounded odd functions has the same
supremum as over odd signs.

\begin{proposition}[Exact radial Hermite energies]
\label{prop:new-all-radial-energies}
For $X\sim\gamma_d$, write $R=\|X\|$, $V=X/R$, and
\[
 a_d=\frac{(\mathbb ER)^2}{d}
 =\frac{2\Gamma((d+1)/2)^2}{d\Gamma(d/2)^2},\qquad \alpha=d/2.
\]
With the generalized Laguerre polynomial $L_k^\alpha$, one has
\begin{align}
 P_{2k+1}^{(d)}V_i
  &=\frac{\mathbb ER}{d}\frac{(1/2)_k}{(\alpha+1)_k}
             X_iL_k^\alpha(R^2/2),\label{eq:new-vector-projection}\\
 e_{2k+1}(d):=\sum_i\|P_{2k+1}^{(d)}V_i\|_2^2
  &=a_d\frac{(1/2)_k^2}{k!(\alpha+1)_k}.
             \label{eq:new-radial-energy}
\end{align}
In particular
\[
 e_1=a_d,\quad e_3=\frac{a_d}{2(d+2)},\quad
 e_5=\frac{9a_d}{8(d+2)(d+4)},\quad
 e_7=\frac{75a_d}{16(d+2)(d+4)(d+6)}.
\]
The exact natural-vector objective is
\begin{equation}\label{eq:new-exact-natural-value}
 q_d(T)=\tau+\sum_{k=0}^Kc_ke_{2k+1}(d)\longrightarrow1.
\end{equation}
\end{proposition}
\begin{proof}
Gaussian projection commutes with orthogonal rotations.  Thus the
projection of the equivariant vector field $V$ onto degree $2k+1$ is
a constant times $X_iL_k^\alpha(R^2/2)$.  Indeed an equivariant vector
polynomial is $X_i$ times a radial polynomial, and orthogonality to
the lower odd degrees is precisely Laguerre orthogonality with
weight $t^\alpha e^{-t}$.
Put $t=R^2/2$.  Radial integration and Laguerre orthogonality give
\[
 \mathbb E[X_i^2L_k^\alpha(t)^2]=\frac{(\alpha+1)_k}{k!},\qquad
 \mathbb E[V_iX_iL_k^\alpha(t)]
       =\frac{\mathbb ER}{d}\frac{(1/2)_k}{k!}.
\]
For the second identity, use
\[
 \int_0^\infty t^{\alpha-1/2}e^{-t}L_k^\alpha(t)\,dt
       =\Gamma(\alpha+1/2)\frac{(1/2)_k}{k!},
\]
which follows by integrating the generating function
$(1-z)^{-\alpha-1}\exp[-tz/(1-z)]$ and comparing coefficients.
Taking the ratio proves \eqref{eq:new-vector-projection}; its squared
norm proves \eqref{eq:new-radial-energy}.  Parseval gives
$\sum_{k\geq0}e_{2k+1}(d)=1$.  Finally $a_d\to1$ and every fixed
$e_{2k+1}(d)$ with $k\geq1$ tends to zero.  Since $\tau+c_0=1$,
\eqref{eq:new-exact-natural-value} follows.
\end{proof}

\begin{proposition}[An explicit reduced matrix sequence]
\label{prop:new-odd-tail-matrix-sequence}
For $n\geq2$, set $d=n^2$.  Partition each Gaussian coordinate into
intervals of mesh $n^{-2}$ from $-n$ to $n$, together with the two
exterior intervals.  Retain one cell $C$ from each antipodal pair,
using the cells with first coordinate negative.  Set
\[
 w_C=2\gamma_d(C),\qquad D_w=\operatorname{diag}(w_C),\qquad
 F_j[C,\xi]=2\int_C H_\xi\,d\gamma_d\quad(|\xi|=j),
\]
where $H_\xi$ is the normalized product Hermite polynomial.  Define
\begin{equation}\label{eq:new-odd-tail-seed}
 A_n(T)=\tau D_w+\sum_{k=0}^Kc_kF_{2k+1}F_{2k+1}^T.
\end{equation}
Then
\begin{equation}\label{eq:new-tail-matrix-limits}
 C(A_n(T))\to N(T),\qquad V(A_n(T))\to1,
 \qquad G(A_n(T))\to\frac1{N(T)}.
\end{equation}
Choose dyadic rational approximants with total entrywise error at most
$\eta_n=n^{-4}$.  Clear their denominators, delete zero rows and
columns, and apply the exact absolute row/column cloning reduction,
the symmetric off-diagonal dilation, and the signed balanced copy.
The resulting rational matrices $H_n(T)$ are symmetric, zero-diagonal,
have zero row and column sums, and have $|H_n(T)|$ doubly stochastic.
Their ratios have the same limit $1/N(T)$.
\end{proposition}
\begin{proof}
Odd extension of the cell values turns the finite bilinear form
\eqref{eq:new-odd-tail-seed} into $\langle f,T^{(d)}g\rangle$.
The diagonal term is exact because for odd cell functions
\[
 \langle f,P_{\rm odd}^{(d)}g\rangle
           =\sum_Cw_Cf_Cg_C.
\]
The remaining terms are the exact Hermite projection pairings.
Thus $C(A_n(T))\leq N(T)$.  Every pair of bounded odd functions in a
fixed dimension can be approximated in $L_2$ by conditional averages
on the first coordinates of these grids as $n\to\infty$.
Continuous compactly supported approximation proves this even though
the grids are not nested.  The other coordinates integrate out, and
$T^{(d)}$ is consistent under such cylindrical embeddings.  Since
its $L_2$ norm is one, the pairings converge.  Finite cube maxima
occur at sign vertices, proving $\liminf C(A_n(T))\geq N(T)$.

For vector-valued odd cell functions of pointwise norm one, the same
$L_2$ operator bound and Cauchy--Schwarz give $V(A_n(T))\leq1$.
Use the unit vector at the clipped cell midpoint as the trial field.
The quantitative estimate below proves that its objective tends to
$q_d(T)$ and hence to one.  This proves the vector limit.  Positivity
is explicit: for $s=\operatorname{sgn}X_1$, first-chaos energy is
$2/\pi$ and all remaining odd eigenvalues are at least $-1$, whence
$N(T)\geq\langle s,Ts\rangle\geq4/\pi-1>0$.

Both finite norms change by at most total entrywise approximation
error.  Consequently rational approximation preserves both limits.
For completeness, if $Z$ is the cleared integer seed, put
$r_i=\sum_j|Z_{ij}|$, $s_j=\sum_i|Z_{ij}|$, and replace its entry by
$Z_{ij}/(r_is_j)$ on each of the $r_i$ row and $s_j$ column copies.
Both norms are unchanged, because every original row and column has
total absolute copying weight one; the absolute matrix is doubly
stochastic.  Off-diagonal symmetric dilation multiplies both norms
by two.  Tensoring with
$\tfrac12\left(\begin{smallmatrix}1&-1\\-1&1\end{smallmatrix}\right)$
multiplies both by two again and imposes zero sums.  All ratios are
therefore preserved exactly.
\end{proof}

\begin{proposition}[The general asymptotic vector value]
\label{prop:new-general-vector-value}
Drop \eqref{eq:new-odd-spectral-cube} and assume only that $T$ is
nonzero.  Write
\[
 M=\max\bigl(|\tau|,|\tau+c_0|,\ldots,|\tau+c_K|\bigr).
\]
Choose nonzero rational approximants, or discard the finitely many initial
zero approximants before applying the reductions. This is possible since
$C(A_n(T))\to N(T)>0$. The seed \eqref{eq:new-odd-tail-seed}
and these rational reductions satisfy
\[
 C(A_n(T))\to N(T),\qquad V(A_n(T))\to M,
                    \qquad G(H_n(T))\to M/N(T).
\]
For every fixed odd Hermite degree $j$ and $d\geq j$, the square-free
tensor trial field used in the proof has norm-square variance
\[
 \epsilon_{d,j}^2=
 \sum_{r=1}^j
 \frac{\binom jr\binom{d-j}{j-r}}{\binom dj}(3^r-1)
                  \leq\frac{j(3^j-1)}d.
\]
Terms with an impossible intersection size have binomial coefficient
zero.  Its unit normalization attains $|\lambda_j|$
within $2M\epsilon_{d,j}$ in continuum vector objective.
\end{proposition}
\begin{proof}
Let $m=\binom dj$ and define the vector with components indexed by
$j$-subsets $S$ of $\{1,\ldots,d\}$,
\[
 Y_S(X)=m^{-1/2}\prod_{i\in S}X_i,\qquad U=Y/\|Y\|.
\]
The exceptional zero set has Gaussian measure zero.  Every component
has pure degree $j$, and $\mathbb E\|Y\|^2=1$.  For two subsets
$S,S'$, their norm-square moment is $3^{|S\cap S'|}$.
Counting intersections proves the variance formula.  If $R$ is their
intersection size, $\mathbb ER=j^2/d$ and
$3^r-1\leq r(3^j-1)/j$ for $0\leq r\leq j$, giving its upper bound.
The elementary inequality $(\sqrt{x}-1)^2\leq(x-1)^2$ yields
\[
 \|U-Y\|_{L_2}^2=\mathbb E(\|Y\|-1)^2
                   \leq\epsilon_{d,j}^2.
\]
Thus
$|\langle U,TU\rangle-\lambda_j|\leq2M\epsilon_{d,j}$.
If $\lambda_j<0$, use $-U$ on the right.  Any selected nonzero odd
eigenvalue is covered; for the tail eigenvalue choose an odd
$j>2K+1$.  This proves that the continuum supremum of the vector
objective over all dimensions is $M$.

For each fixed trial dimension, conditional cell averages of $U$
on the first coordinates converge in $L_2$ as $n\to\infty$.
They have pointwise norm at most one; completing left and right
ball vectors in separate orthogonal coordinates gives admissible
finite unit vectors with the same cross pairings.  Therefore the
seed vector norms have $\liminf\geq M$.  The $L_2$ norm of $T$ gives
the reverse bound.  The scalar convergence and rational reductions
are the same as in the preceding proof, with operator norm $M$
replacing one.  A nonzero multiplier has $N(T)>0$, since otherwise
all its pairings against bounded functions vanish, contradicting its
nonzero Hermite action.
\end{proof}

\paragraph{An explicit finite error through degree seven.}
Here $K\leq3$.  Set missing $c_k$ equal to zero and put
\[
 \lambda_*=\min(\tau,\tau+c_0,\tau+c_1,\tau+c_2,\tau+c_3),
 \qquad C_* = |\tau-\lambda_*|.
\]
Define the following nonnegative rational expressions in $d$ and
the absolute coefficient values:
\begin{align*}
 B_0&=|c_0|+|c_1|/2+3|c_2|/8+5|c_3|/16,\\
 B_1&=\frac{|c_1|/2+3|c_2|/4+15|c_3|/16}{d+2},\\
 B_2&=\frac{3|c_2|/8+15|c_3|/16}{(d+2)(d+4)},\\
 B_3&=\frac{5|c_3|/16}{(d+2)(d+4)(d+6)}.
\end{align*}
Let $J_0=1$, $J_j=\prod_{a=1}^j(d+2a-1)$ and
$M_j=\prod_{a=0}^{j-1}(d+2a)$, with $M_0=1$.  Put
\begin{align*}
 Q_d^*&=C_*^2+2C_*\sum_{j=0}^3B_jJ_j
          +\frac1d\sum_{i,j=0}^3B_iB_jM_{i+j+1},\\
 \mathcal L_n(T)&=
 \frac{d}{n^4}\left[
 \frac{C_*}{d-2}+\frac{B_0}{d-1}+B_1
 +(d+1)B_2+(d+1)(d+3)B_3\right]
       +4\sqrt{Q_d^*}\sqrt{2n\varphi(n)},\quad d=n^2.
\end{align*}
Then $\mathcal L_n(T)=O(n^{-4})$ and
\begin{equation}\label{eq:new-tail-finite-lower}
 V(A_n(T))\geq q_{n^2}(T)-\mathcal L_n(T).
\end{equation}
In particular, if a separately proved all-sign bound is $N(T)\leq B$,
the rational reduced sequence satisfies the completely explicit bound
\[
 G(H_n(T))\geq
 \frac{\max(0,q_{n^2}(T)-\mathcal L_n(T)-n^{-4})}{B+n^{-4}}.
\]
This is a lower bound; it does not identify $N(T)$ with $B$.

\begin{proof}[Proof of the error estimate]
The radial projection formula gives $TV=t_d(R)V$, where
\[
 t_d(R)=\tau+\frac{\mathbb ER}{d}R\sum_{k=0}^3
 c_k\frac{(1/2)_k}{(d/2+1)_k}L_k^{d/2}(R^2/2).
\]
Let $v$ be the clipped-midpoint unit field and $\delta=v-V$.
Both fields are odd and unit length, so
$2\delta\cdot V=-\|\delta\|^2$.  Since the minimum odd eigenvalue
is $\lambda_*$,
\[
 \langle v,Tv\rangle\geq q_d(T)
             -\mathbb E[(t_d(R)-\lambda_*)\|\delta\|^2].
\]
The coefficientwise absolute Laguerre expansion shows
\[
 (t_d(R)-\lambda_*)_+\leq
 B_d(R):=C_*+\frac{\mathbb ER}{d}
                         \sum_{j=0}^3 B_jR^{2j+1}.
\]
On the no-clipping event normalization gives
$\|\delta\|^2\leq dn^{-4}/R^2$; on its complement it is at most
four.  The latter event has probability at most $2n\varphi(n)$.
Gaussian radial moments give
\[
 \mathbb ER^{2j}=M_j,\qquad
 \mathbb ER^{2j+1}=(\mathbb ER)J_j,\qquad
 \mathbb ER^{-1}=\frac{\mathbb ER}{d-1},\quad
 \mathbb ER^{-2}=\frac1{d-2}.
\]
Use $(\mathbb ER)^2/d=a_d\leq1$ in
$\mathbb E[B_d(R)/R^2]$ and $\mathbb EB_d(R)^2$, followed by
Cauchy--Schwarz on the clipping event.  These are exactly the two
terms of $\mathcal L_n(T)$.  Finally $B_j=O(d^{-j})$, $J_j=O(d^j)$,
$M_j=O(d^j)$, which proves the asserted rate.
\end{proof}

\paragraph{Two useful subfamilies.}
For $T=P_1-b_3P_3-b_5P_5-b_7P_7$ the exact numerator is
$a_d-b_3e_3-b_5e_5-b_7e_7$.  For the normalized identity-tail family
\[
 T_{\tau,b}=(1-\tau)P_1-bP_3+\tau P_{\rm odd},
 \qquad q_d(T_{\tau,b})=\tau+(1-\tau)a_d-be_3.
\]
In the latter family the spectral-cube conditions are
$|\tau|\leq1$ and $|\tau-b|\leq1$.  The identity tail is represented
by the diagonal in \eqref{eq:new-odd-tail-seed}, and is not discarded
at finite $n$.  Its numerator has leading loss
$(1-\tau+b)/(2d)$.  For every such family, a useful numerical endpoint
still requires a proved bound for the complete scalar norm $N(T)$.

% Genuine finite-matrix sequence refinement. Denominator convergence is
% qualitative; the rate applies to the displayed certified lower bounds.
\begin{proposition}[A powered-dimension cubic sequence]
\label{prop:new-powered-cubic-sequence}
Fix $0<b<1$, and suppose a complete all-sign theorem proves
$N(P_1-bP_3)\leq B$.  In the analytic cell-compression construction,
take
\[
 d_n=n^4,\qquad h_n=n^{-4},\qquad \text{coordinate cutoff }n,
                   \qquad \eta_n=n^{-8},\quad n\geq2.
\]
Here $h_n$ is the interval mesh and $\eta_n$ is the total entrywise
dyadic approximation error.  Apply the same exact cloning, symmetric
dilation and balanced copy, producing reduced rational matrices
$\widehat H_n(b)$.  Put $d=d_n$ and
\[
 p_n=\min(1,2n^3\varphi(n)),\qquad
 q_d(b)=a_d\left(1-\frac b{2(d+2)}\right),
\]
\begin{align*}
 \widehat{\mathcal L}_n(b)={}&dh_n^2
 \left[\frac{1-b/2}{d-1}+\frac b{2(d+2)}+\frac b{d-2}\right]\\
 &+4\left[1-\frac b2+\frac b2\sqrt{\frac{d+4}{d+2}}\right]
                         \sqrt{p_n}+4bp_n.
\end{align*}
Then the explicit finite lower bound is
\[
 G(\widehat H_n(b))\geq
 \frac{\max(0,q_d(b)-\widehat{\mathcal L}_n(b)-n^{-8})}
                           {B+n^{-8}}.
\]
The matrices have actual ratio limit $1/N(P_1-bP_3)$.
The displayed certified lower bounds converge to $1/B$ at rate
$O(n^{-4})$.  No rate for the actual scalar-norm convergence is
asserted.
\end{proposition}
\begin{proof}
For $V=X/R$, the radial projection identity gives
\[
 (P_1-bP_3)V=t_d(R)V,\qquad
 t_d(R)+b=b+\frac{\mathbb ER}{d}R
           \left[1-\frac b2+\frac{bR^2}{2(d+2)}\right]\geq0.
\]
The unit-field difference identity and minimum odd eigenvalue $-b$
give a trial-objective loss at most
$\mathbb E[(t_d(R)+b)\|v-V\|^2]$.
Off the clipping event, $\|v-V\|^2\leq dh_n^2/R^2$.
The exact negative radial moments and $a_d\leq1$ give the first
line of $\widehat{\mathcal L}_n$.  The clipping event has probability
at most $p_n$ by the union bound and Gaussian Mills inequality.
On that event $\|v-V\|^2\leq4$.  Cauchy--Schwarz and
\[
 \mathbb ER^2=d,\qquad \mathbb ER^6=d(d+2)(d+4),\qquad a_d\leq1
\]
give its remaining two terms.  These steps use no independence of
the clipping event and $R$.

As before, analytic cell compression has scalar norm at most $B$
and vector norm at least $q_d-\widehat{\mathcal L}_n$.
Entrywise approximation changes each norm by at most $\eta_n$ and
the exact reductions preserve the ratio, proving the finite bound.
The qualitative scalar/vector limits hold for the new nonnested
grids because both the dimension and cutoff tend to infinity and
the mesh tends to zero.
Finally
\[
 q_{n^4}(b)=1-\frac{1+b}{2n^4}+O(n^{-8}),\qquad
 \widehat{\mathcal L}_n(b)=(1+b)n^{-8}+O(n^{-12})
                         +O(n^{3/2}e^{-n^2/4}).
\]
The latter exponential term is $o(n^{-m})$ for every fixed $m$.
Together with $\eta_n=n^{-8}$ this proves the stated lower-bound
rate.  Qualitative convergence of the scalar norm supplies the
actual limit, without supplying a quantitative denominator rate.
\end{proof}

\paragraph{Extension to a normalized odd tail.}
For $T_{\tau,b}=(1-\tau)P_1-bP_3+\tau P_{\rm odd}$, let
$\alpha=1-\tau\geq b/2$, and impose the spectral-cube conditions.
The same proof holds with $q_d=\tau+\alpha a_d-be_3$ and
\[
 \widehat{\mathcal L}_n(\tau,b)=dh_n^2
 \left[\frac{\alpha-b/2}{d-1}+\frac b{2(d+2)}+\frac b{d-2}\right]
 +4\left[\alpha-\frac b2+\frac b2\sqrt{\frac{d+4}{d+2}}\right]
                       \sqrt{p_n}+4bp_n.
\]
Its leading natural-vector loss is $(\alpha+b)/(2n^4)$, and its
midpoint bulk loss is $(\alpha+b)n^{-8}$.  A separately certified
complete norm bound remains necessary for every selected $(\tau,b)$.

% Exact mathematical envelope lemma. Every numerical input still needs its
% own complete all-sign criterion and directed constant certificate.
\begin{proposition}[Exact envelope of criteria with a common quadratic]
\label{prop:new-rank-dominance-envelope}
Let $P\subset\mathbb R^2$ be a bounded convex polygon, allowing a
segment, a point, or the empty set.  Fix $b>0$ and finitely many
criteria
\[
 F_i(u,r)=u^2+br^2+A_i u+B_i r+C_i,\qquad 1\leq i\leq m.
\]
For each $i$, form the closed dominance polygon
\[
 P_i=P\cap\bigcap_j
 \{(u,r):(A_i-A_j)u+(B_i-B_j)r\leq C_j-C_i\}.
\]
If $P$ is nonempty, then
\begin{equation}\label{eq:new-exact-rank-envelope}
 \max_{(u,r)\in P}\min_iF_i(u,r)
       =\max_{i:P_i\ne\varnothing}\;
                         \max_{v\in\operatorname{Vert}(P_i)}F_i(v).
\end{equation}
For a segment its vertices are its endpoints; for a point the sole
vertex is the point.  All boundaries and ties are retained.
If the inputs are rational, every clipping intersection and every
price in \eqref{eq:new-exact-rank-envelope} is computed exactly by
rational arithmetic.
\end{proposition}
\begin{proof}
At every point of $P$ at least one affine term
$A_i u+B_i r+C_i$ is minimal.  Thus the closed $P_i$ cover $P$;
ties simply belong to more than one polygon.  On $P_i$ the common
quadratic cancels from all comparisons, so $\min_jF_j=F_i$.
Each $F_i$ is convex because its Hessian is
$\operatorname{diag}(2,2b)\succ0$.  Expressing any point of $P_i$
as a convex combination of its vertices proves
$F_i\leq\max_{v\in\operatorname{Vert}(P_i)}F_i(v)$.
Taking maxima over the covering polygons proves equality.
The same argument handles lower-dimensional polygons, and closed
half-plane clipping preserves their boundary points.
\end{proof}

\paragraph{Use in a global all-sign cover.}
Each $C_i$ must be a proved uniform upper price over the same fixed
cumulant--shape cell, and every $F_i$ must bound the same operator
at one fixed coefficient $b$.  Then the right side of
\eqref{eq:new-exact-rank-envelope} is an upper bound for every physical
sign pair in that source polygon.  An arbitrary nonempty subset of
the valid criteria is sufficient; it can only increase the pointwise
minimum.  A floating calculation may nominate this subset, but not
replace exact clipping and pricing.  An inherited uniform bound may
be retained by taking its minimum with this new envelope bound.
For rank-one R3 criteria, exact source transport gives the common
coefficient $b$ of $r^2$.  More generally a profile is eligible
whenever its computed rational $r^2$ coefficient equals the same
$b$; a profile label is not a substitute for that equality.


\section{A smaller truncation grid and physical tensor constraints}

% A smaller implicit grid at each fixed dimension, not merely reindexing.
\begin{proposition}[A logarithmic-cutoff reduced sequence]
\label{prop:new-log-cutoff-sequence}
Fix $0<b<1$ and a separately certified complete norm cap
$N(P_1-bP_3)\leq B$.  For each integer $d\geq4$, let
\[
 R_d=\min\{R\in\mathbb Z_{\geq1}:\exp(R^2/10)\geq d\}
           =\lceil\sqrt{10\log d}\rceil.
\]
Use mesh $h_d=1/d$ on $[-R_d,R_d]$ in every coordinate, the two
exterior intervals, odd half-cell compression, and total entrywise
dyadic approximation error $\eta_d=d^{-2}$.  Apply the same exact
reduced-class maps.  Write the resulting rational matrix as $H_d^*(b)$.
Put
\[
 p_d=\min\left(1,\frac{2d\varphi(R_d)}{R_d}\right),\quad
 q_d(b)=a_d\left(1-\frac b{2(d+2)}\right),
\]
\begin{align*}
 \mathcal L_d^*(b)={}&\frac1d
 \left[\frac{1-b/2}{d-1}+\frac b{2(d+2)}+\frac b{d-2}\right]\\
 &+4\left[1-\frac b2+\frac b2\sqrt{\frac{d+4}{d+2}}\right]
                                      \sqrt{p_d}+4bp_d.
\end{align*}
Then
\[
 G(H_d^*(b))\geq
 \frac{\max(0,q_d(b)-\mathcal L_d^*(b)-d^{-2})}{B+d^{-2}}.
\]
The actual ratios converge to $1/N(P_1-bP_3)$.
The displayed certified lower bounds converge to $1/B$ at rate
$O(d^{-1})$, with $\mathcal L_d^*(b)=O(d^{-2})$.
Each coordinate has only $2dR_d+2=O(d\sqrt{\log d})$ intervals,
instead of $O(d^{3/2})$ for cutoff $\sqrt d$ at the same mesh.
The implicit half-cell count is $(2dR_d+2)^d/2$, before further
rational cloning; no practical enumeration guarantee is asserted.
\end{proposition}
\begin{proof}
Mills' inequality and the coordinate union bound give clipping
probability at most $p_d$.  From the definition of $R_d$,
\[
 \varphi(R_d)\leq(2\pi)^{-1/2}d^{-5},\qquad
 p_d\leq\sqrt{2/\pi}\,d^{-4}/R_d\leq d^{-4}.
\]
The exact radial/unit-field calculation in the powered-dimension
proposition gives $\mathcal L_d^*$ after substituting $h=1/d$ and
this clipping probability.  Its bulk term is $O(d^{-2})$, and
the displayed probability bound gives the same order for its tail.
Entrywise approximation and exact cloning/balancing give the finite
ratio bound.  Both $R_d\to\infty$ and $h_d\to0$, so conditional
averaging of every fixed cylindrical odd function still converges
in $L_2$, proving the qualitative scalar norm limit.  The radial
unit field proves vector norm convergence to one.
Finally $q_d(b)=1-(1+b)/(2d)+O(d^{-2})$, proving the stated
rate only for the certified lower bounds.

The integer cutoff can be computed without assuming a floating
logarithm is exact.  Test $\exp(R^2/10)\geq d$ with increasingly
precise directed intervals.  Equality is impossible for positive
$R$ and integer $d$, because the exponential of a nonzero algebraic
number is transcendental.  Therefore every comparison terminates
and determines the smallest admissible integer.
\end{proof}

\paragraph{Compatibility with a powered index and an odd tail.}
Taking $d=n^4$ retains certified lower-bound rate $O(n^{-4})$, now
with $2n^4R_{n^4}+2$ intervals in each coordinate and rational error
$n^{-8}$.  For the normalized tail family with
$\alpha=1-\tau\geq b/2$, replace $q_d$ by
$\tau+\alpha a_d-be_3$ and replace each $1-b/2$ in the loss formula
by $\alpha-b/2$.  Its uniform scalar norm needs a separate complete
certificate.  There is still no quantitative rate claimed for the
actual denominator convergence.


% New necessary physical constraints and complete hyperbolic-boundary pricing.
% They require independent numeric cache/geometry acceptance when implemented.
\begin{proposition}[A necessary cubic cumulant--shape constraint]
\label{prop:new-cubic-shape-product}
Let $f$ be a real symmetric order-three Hilbert tensor of norm one,
$Q=I_3(f)/\sqrt6$, and
\[
 A=\|f\otimes_2f\|^2,\quad
 B=\|\operatorname{sym}(f\otimes_1f)\|^2,\quad
 \kappa=\mathbb EQ^4-3,\quad
 V=\operatorname{Var}(\|DQ\|^2/3).
\]
Then
\[
 \kappa=36A+54B,\qquad V=8A+6B,
 \qquad 0\leq B\leq A\leq1.
\]
For $\kappa>0$, define the physical shape $z=9V/\kappa-1$.
It satisfies
\begin{equation}\label{eq:new-physical-shape-product}
 \frac25\leq z\leq1,\qquad \kappa z=36A\leq36.
\end{equation}
\end{proposition}
\begin{proof}
The product and derivative formulas give
\[
 Q^2=1+3I_2(f\otimes_2f)
 +\tfrac32 I_4(\operatorname{sym}(f\otimes_1f))
 +\tfrac16 I_6(\operatorname{sym}(f\otimes f)),
\]
\[
 \|DQ\|^2/3=1+2I_2(f\otimes_2f)
                    +\tfrac12 I_4(\operatorname{sym}(f\otimes_1f)).
\]
The norm-square of the last order-six tensor is $(2+18A)/20$:
in the twenty possible groupings of its six indices, two preserve
the original grouping and eighteen have contraction inner product $A$.
Chaos orthogonality now proves the two moment identities.
Flatten $f$ into a matrix $F$ with one index on its row.  Its two
Gram matrices have the same nonzero eigenvalues, so
$\|f\otimes_1f\|^2=\|f\otimes_2f\|^2=A$.
Orthogonal symmetrization gives $B\leq A$.  Also $FF^T$ is positive
semidefinite with trace one; hence $A=\operatorname{tr}(FF^T)^2\leq1$.
Finally $9V-\kappa=36A$, and
$z=36A/(36A+54B)$ gives its stated range.
\end{proof}

\begin{proposition}[Complete pricing on the hyperbolic physical boundary]
\label{prop:new-curved-shape-price}
On a box $\kappa\in[k_-,k_+]\subset[0,90]$,
$z\in[z_-,z_+]\subset[2/5,1]$, consider
\[
 C(\kappa,z)=C_0+(L_0+L_1z)\kappa+D_0\kappa^2
 +\kappa^{3/2}\left[A_0(K-t(z-2/5))+sD_0z^{3/2}\right],
\]
where $A_0,D_0\geq0$, $K=194/\sqrt{10}$, $t=633/9$, $s=9/4$,
and the other coefficients are arbitrary reals.  Its maximum on the
physical domain $\kappa z\leq36$ is obtained by the following finite
boundary-endpoint check.
Put $k_* = \min(k_+,36/z_-)$ and $k_t=36/z_+$.
If $k_->k_*$ the domain is empty.  Otherwise test
\begin{itemize}
 \item $(k_-,z_-)$ and $(k_*,z_-)$;
 \item if $k_-\leq\min(k_*,k_t)$, test
       $(k_-,z_+)$ and $(\min(k_*,k_t),z_+)$;
 \item if $\max(k_-,k_t)\leq k_*$, test
       $(\max(k_-,k_t),36/\max(k_-,k_t))$
       and $(k_*,36/k_*)$.
\end{itemize}
Repeated points may be removed.  The resulting maximum is exact for
this scalar price, including boxes touching the physical boundary.
Its point evaluations can be certified by directed arithmetic.
\end{proposition}
\begin{proof}
For fixed $\kappa$, the price is convex in $z$, so its maximum on
the feasible vertical interval is at $z_-$ or
$\min(z_+,36/\kappa)$.  At fixed $z$, it is convex in $\kappa$:
the coefficient of $\kappa^{3/2}$ is nonnegative because
$K-t(z-2/5)\geq194/\sqrt{10}-211/5>0$.
The lower horizontal edge and the upper horizontal portion therefore
need only their endpoints.  On the remaining hyperbolic edge,
\begin{align*}
 C(\kappa,36/\kappa)={}&C_0+36L_1+sD_0\,36^{3/2}
 +L_0\kappa+D_0\kappa^2\\
 &+A_0(K+2t/5)\kappa^{3/2}-36A_0t\sqrt\kappa.
\end{align*}
Each non-affine term here is convex: $D_0,A_0,t\geq0$ and
$-\sqrt\kappa$ is convex.  Its interval maximum is therefore at
the two endpoints as well.  The displayed candidates enumerate all
these boundary endpoints, proving the claim.  The hyperbolic
portion starts at $k_t\geq36$, so no division by zero occurs.
\end{proof}

\paragraph{Application and scope.}
The frozen R3 moment-price formula has this exact form with
$A_0=q_6+\zeta_{\rm profile}^2$ and
$D_0=G-d_0-\zeta_{\rm profile}^2$, provided both are nonnegative.
The accepted scalar B0 plane prices have the same form with
$A_0=q_6\geq0$ and $D_0=0$.  Thus each can be priced on the
necessary physical constraint instead of an unrestricted shape box.
Simply deleting infeasible rectangle corners would be invalid;
the newly exposed curved-boundary endpoints must also be priced.
This lemma does not assert that any resulting numerical cache or
full operator cover has already been certified.

\begin{proposition}[An optional stronger high-contraction constraint]
\label{prop:new-high-contraction-shape}
In the tensor notation above,
\[
 B\geq(3A-2)_+^2.
\]
More strongly, for $A\geq5/9$,
\[
 B\geq\left(\frac{3\sqrt{2A-1}-1}{2}\right)^2.
\]
Writing $w=36A=\kappa z$, define
\[
 F(w)=\begin{cases}
 w,&0\leq w\leq20,\\
 \dfrac{31w}{4}-108-81\sqrt{w/18-1},&20\leq w\leq36.
 \end{cases}
\]
Every physical source satisfies $\kappa\geq F(w)$.  The function
$F$ is increasing and convex; $F(w)/w$ is nondecreasing for $w>0$.
Consequently an explicit stronger shape cap is
\[
 z\leq1\quad(\kappa\leq20),\qquad
 z\leq\frac{18}{\kappa}
 \left[1+\left(\frac{27+\sqrt{62\kappa-1224}}{93}\right)^2\right]
                               \quad(20\leq\kappa\leq90).
\]
The latter cap decreases from one to $2/5$.
\end{proposition}
\begin{proof}
Let $u$ be a unit top eigenvector of $FF^T$, with eigenvalue
$\lambda$.  Since its trace is one, $\lambda\geq A$.
Symmetry of $f$ means that projection onto $u$ in each tensor leg
has squared norm $\lambda$.  These three projections commute, so
the union bound for their complements gives
$|f(u,u,u)|^2\geq3\lambda-2$.
For $T=\operatorname{sym}(f\otimes_1f)$,
\[
 T(u,u,u,u)=\sum_a f(u,u,e_a)^2\geq|f(u,u,u)|^2,
 \qquad \|T\|^2\geq T(u,u,u,u)^2.
\]
This proves the first inequality.  If $A>1/2$, the remaining Gram
eigenvalues sum to $1-\lambda$, so
$A\leq\lambda^2+(1-\lambda)^2$.  Because $\lambda\geq A>1/2$,
this implies $\lambda\geq(1+\sqrt{2A-1})/2$ and proves the stronger
inequality when its lower bound is nonnegative.

Substitute it into $\kappa=36A+54B$ to obtain $\kappa\geq F(w)$.
On its curved part put $s=\sqrt{w/18-1}\in[1/3,1]$.
Then $F'(w)=31/4-9/(4s)\geq1$ and $F''(w)=1/(16s^3)>0$;
both match the affine part at $w=20$.  Also
$wF'(w)-F(w)=108+(81/2)(s-1/s)\geq0$.
The inverse relation is $s=(27+\sqrt{62\kappa-1224})/93$,
which gives the stated cap.  Its monotonicity follows from that
of $F(w)/w$, and its endpoint values follow from $F(20)=20$,
$F(36)=90$.
\end{proof}


% Strong physical cubic constraint.  Analytic lemma only: a numerical
% norm improvement requires freshly accepted directed prices and cover.
\begin{proposition}[A stronger physical cubic shape constraint]
\label{prop:new-strong-cubic-shape}
Let $f$ be a real symmetric order-three Hilbert tensor with $\|f\|=1$,
and put
\[
 Q=I_3(f)/\sqrt6,\qquad
 A=\|f\otimes_2f\|^2,\qquad
 B=\|\operatorname{sym}(f\otimes_1f)\|^2.
\]
Then
\begin{equation}\label{eq:new-strong-cubic-contraction}
 B\geq(2A-1)_+.
\end{equation}
In particular, if $\kappa=\mathbb E Q^4-3$ and
$z=9\operatorname{Var}(\|DQ\|^2/3)/\kappa-1$ for $\kappa>0$, then
\begin{equation}\label{eq:new-strong-cubic-shape}
 \frac25\leq z\leq1,\qquad 0\leq\kappa\leq90,
 \qquad 4\kappa z-\kappa\leq54.
\end{equation}
Thus $z\leq\tfrac14+\tfrac{27}{2\kappa}$ whenever $\kappa>0$.
This last constraint also implies $\kappa z\leq36$ on
$0\leq\kappa\leq90$.
\end{proposition}
\begin{proof}
Flatten $f$ to an operator $F$ with one index on its row, and set
$G=FF^*$.  It is positive with trace one.  Its largest eigenvalue
$\lambda$ therefore satisfies
$\lambda\geq\operatorname{tr}G^2=A$.  Let $u$ be a corresponding
unit eigenvector.  If $A>1/2$, then $\lambda>1/2$ and
\[
 A\leq\lambda^2+(1-\lambda)^2.
\]
Indeed, the remaining nonnegative eigenvalues sum to $1-\lambda$,
so the sum of their squares is at most $(1-\lambda)^2$.
The orthogonal projections onto $u$ in the first and second tensor
legs commute.  Each projects $f$ to squared norm $\lambda$ by
symmetry, and hence
\[
 p:=\|f(u,u,\cdot)\|^2\geq2\lambda-1>0.
\]
For the symmetric fourth-order tensor
$T=\operatorname{sym}(f\otimes_1f)$ we have
$T(u,u,u,u)=p$.  Cauchy--Schwarz and $\|u^{\otimes4}\|=1$ give
\[
 B=\|T\|^2\geq p^2\geq(2\lambda-1)^2\geq2A-1.
\]
If $A\leq1/2$, nonnegativity of $B$ proves the same positive-part
bound.

The Gaussian product formulas give
$\kappa=36A+54B$ and
$\operatorname{Var}(\|DQ\|^2/3)=8A+6B$, so $\kappa z=36A$.
Also $0\leq B\leq A\leq1$, since orthogonal symmetrization is a
contraction and the two flattening Gram matrices have the same
nonzero eigenvalues.  These facts give the first two bounds in
\eqref{eq:new-strong-cubic-shape}.  If $A>1/2$, the new contraction
bound gives
$\kappa\geq36A+54(2A-1)=144A-54=4\kappa z-54$.
If $A\leq1/2$, then $4\kappa z-54\leq\kappa z\leq\kappa$.
Finally $(\kappa+54)/4\leq36$ when $\kappa\leq90$.
\end{proof}

\begin{proposition}[Complete pricing on the stronger physical boundary]
\label{prop:new-strong-curved-shape-price}
Consider a box $\kappa\in[k_-,k_+]\subset[0,90]$ and
$z\in[z_-,z_+]\subset[2/5,1]$.  Let
\[
 C(\kappa,z)=C_0+(L_0+L_1z)\kappa+D_0\kappa^2
 +\kappa^{3/2}\left[A_0(K-t(z-2/5))+sD_0z^{3/2}\right],
\]
where $A_0,D_0\geq0$, $K=194/\sqrt{10}$, $t=633/9$,
$s=9/4$, and $C_0,L_0,L_1$ are arbitrary reals.  Its exact maximum
on the necessary physical domain $4\kappa z-\kappa\leq54$ is
obtained by the following finite endpoint check.  Set
\[
 c=\tfrac14,\qquad d=\tfrac{27}{2},\qquad
 k_* =\min(k_+,d/(z_--c)),\qquad k_t=d/(z_+-c).
\]
The domain is empty if $k_->k_*$.  Otherwise test
\begin{itemize}
 \item $(k_-,z_-)$ and $(k_*,z_-)$;
 \item if $k_-\leq\min(k_*,k_t)$, test
       $(k_-,z_+)$ and $(\min(k_*,k_t),z_+)$;
 \item if $\max(k_-,k_t)\leq k_*$, test
       $(\max(k_-,k_t),c+d/\max(k_-,k_t))$
       and $(k_*,c+d/k_*)$.
\end{itemize}
Repeated points may be removed.  Every candidate is rational for
rational box endpoints, and its price may be certified by directed
arithmetic.
\end{proposition}
\begin{proof}
For each fixed $\kappa$, the price is convex in $z$, so a maximum
occurs at $z_-$ or $\min(z_+,c+d/\kappa)$, with the natural
interpretation at $\kappa=0$.  For fixed $z$ it is convex in
$\kappa$: the coefficient of $\kappa^{3/2}$ is nonnegative because
\[
 K-t(z-2/5)\geq194/\sqrt{10}-211/5>0.
\]
The horizontal boundary portions consequently need only their
endpoints.  On the curved portion, direct substitution gives
\begin{align*}
 C(\kappa,c+d/\kappa)={}&C_0+dL_1+(L_0+cL_1)\kappa+D_0\kappa^2\\
 &+A_0[K+t(2/5-c)]\kappa^{3/2}
       -A_0td\sqrt\kappa+sD_0(c\kappa+d)^{3/2}.
\end{align*}
Each non-affine term is convex in $\kappa$: all its displayed
coefficients are nonnegative except that of $\sqrt\kappa$, whose
negative is convex.  This portion also needs only its endpoints.
The listed candidates are exactly the endpoints of these portions.
Since $k_t\geq d/(1-c)=18$, the curved portion never divides by
zero.  The argument retains points and degenerate intervals on
the boundary.
\end{proof}

\paragraph{Application and scope.}
The frozen R3 price has this form when
$A_0=q_6+\zeta_{\rm profile}^2\geq0$ and
$D_0=G-d_0-\zeta_{\rm profile}^2\geq0$.
Accepted B0 source-plane prices have $A_0=q_6\geq0$ and $D_0=0$.
Combined native scalar/L1 endpoint prices also have $D_0=0$ and
$A_0=q_6\geq0$; their additional shape terms are affine in
$\kappa$ and $\kappa z$.
The same endpoint rule therefore applies to all these guarded
families.  The physical restriction is an analytic theorem;
any numerical improvement still requires fresh directed pricing
of every exposed curved endpoint and a complete source cover.


% Complete independently accepted checkpoint; subsequent improvements are separate.
\section{The first complete refinement}
\label{sec:new-complete-refined-certificate}
\begin{theorem}[The October 2 refined checkpoint]
\label{thm:new-complete-refined-certificate}
For the fixed coefficient $b=711/1000$,
\[
 \mathcal N(b)\leq\frac{140904959}{250000000},\qquad
 G_{\mathbb R}\geq\frac{250000000}{140904959}
             =1.7742455749907283\ldots.
\]
The logarithmic-cutoff rational reduced matrices $H_d^*(b)$ satisfy
\[
 \lim_{d\to\infty}G(H_d^*(b))=\frac1{\mathcal N(b)}
       \geq\frac{250000000}{140904959}.
\]
Their explicitly evaluated lower bounds converge at rate $O(d^{-1})$.
The actual denominator limit is not replaced by its certificate cap.
\end{theorem}
\begin{proof}
The protected all-dimension source cover and its complete same-$b$
lower branch remain unchanged. The latter covers
$0\leq u\leq1381/2000$ with price at most
$563307766601/10^{12}$.
For larger amplitudes every one of the $272579$ protected physical
source parents is retained, as are the $8737$ proved empty cells.

Forty-three complete whole-line local inequalities supplement the
old criteria at the single coefficient $b$. Their exact six spectral
gates and all frozen profile/reserve/proof bindings have independent
checks. The stronger tensor constraint
$4\kappa z-\kappa\leq54$ is priced using the complete boundary
rule above. Independent directed MPFR computations at $768$ bits
replay $3010$ R3 tables with $12040$ exposed endpoint prices and
$700$ selected source planes from $70$ source boxes.
All underlying $60$ source-plane rows are also checked.

The producer-free exact rational geometry replay retains all source
polygons and ties, checks $3176$ rank envelopes with $7421$ closed
dominance regions, and checks every knot of the $953$ selected native
offers. Further subdivision of $4564$ amplitude parents yields a
closed cover by $19693$ children, including $5$ proved empty children;
all other parents keep their preceding bounds. The child replay checks
$994$ rank envelopes and $3275$ native offers using exactly the same
accepted moment prices. Its exact global maximum is
\[
 \frac{563619834426673}{10^{15}}
           <\frac{140904959}{250000000}.
\]
The smaller-amplitude branch is also below the displayed cap. Thus
every physical odd sign pair in every finite Gaussian dimension is
covered by the same fixed operator inequality.

The combined independently accepted certificate is
\texttt{audit/\allowbreak BANK43\_\allowbreak STRONG\_\allowbreak SHAPEU\_\allowbreak REFINED\_\allowbreak GLOBAL\_\allowbreak ACCEPTANCE.json}
in the new-round source tree, with SHA-256
\begin{center}\small\ttfamily
e7af38399f6e31535645eb9cac7c9ab16b1103b29092bcabc75e294fe16e1663
\end{center}
It binds the complete geometry, numeric prices, analytic tensor
restriction, same-coefficient lower branch, and $593$ bound files.
The logarithmic-cutoff limit theorem and the exact ratio-preserving
reductions prove the matrix-sequence conclusions.
\end{proof}

% October 2 analytic upper and hybrid source improvements.
\section{Combining source estimates point by point}
% Analytic criterion only; every numerical use needs its own bound/partition
% replay. Native radial pieces must be split before applying this lemma.
\begin{proposition}[A common affine envelope for source prices]
\label{prop:new-hybrid-affine-source-envelope}
Let $P$ be a nonempty closed polygon, segment, or point contained in
$[u_-,u_+]\times[r_-,r_+]$.  Suppose a finite family of independently
valid upper bounds for a sign-pair objective has the form
\[
 F_i(u,r)=u^2+\gamma_i r^2+a_i u+b_i r+c_i.
\]
Write $r_0=(r_-+r_+)/2$ and define the affine functions
\[
 U_i(u,r)=(u_-+u_++a_i)u-u_-u_++b_i r+c_i+
 \begin{cases}
  \gamma_i((r_-+r_+)r-r_-r_+),&\gamma_i\geq0,\\
  \gamma_i(2r_0r-r_0^2),&\gamma_i<0.
 \end{cases}
\]
For each $i$ form the closed dominance region
$P_i=P\cap\bigcap_j\{U_i\leq U_j\}$, retaining all its vertices even
when it is a segment or point.  Then the objective is bounded by
\[
 \max_{i:P_i\ne\varnothing}\ \max_{v\in\operatorname{vert}(P_i)}U_i(v)
       =\sup_{(u,r)\in P}\min_i U_i(u,r).
\]
Moreover, with
$\varepsilon=((u_+-u_-)^2+\max_i|\gamma_i|(r_+-r_-)^2)/4$,
\[
 \sup_P\min_iF_i\ \leq\ \sup_P\min_iU_i
                          \ \leq\ \sup_P\min_iF_i+\varepsilon.
\]
All intersections and values are exactly rational when the input is
rational.  Any nonempty subset of the valid bounds gives a valid
upper bound by the same construction.
\end{proposition}
\begin{proof}
For $u\in[u_-,u_+]$, the secant error is
$(u_-+u_+)u-u_-u_+-u^2=(u-u_-)(u_+-u)$, between zero and
$(u_+-u_-)^2/4$.  For $\gamma_i\geq0$ the corresponding error in
$r$ is $\gamma_i(r-r_-)(r_+-r)$; for $\gamma_i<0$ the tangent error
is $-\gamma_i(r-r_0)^2$.  Thus $F_i\leq U_i\leq F_i+\varepsilon$
on the rectangle.  Taking pointwise minima and then suprema proves
both inequalities.

The closed regions $P_i$ cover $P$, including ties, and on $P_i$ the
minimum of the affine functions is $U_i$.  A linear function reaches
its maximum on a nonempty compact polygon at a vertex; the same
statement holds for a segment or point.  This proves the displayed
identity.  Removing bounds increases their pointwise minimum and
therefore preserves validity.
\end{proof}

For a native price that is given by different quadratic formulas
between radial knots, first partition the polygon at every such
knot and apply the proposition separately to each closed piece.
Moment and source constants used on a piece must bound the entire
same moment cell; the lemma does not license mixing coefficients
from different valid profiles.

\section{A self-contained upper bound and numerical precision}
% Self-contained universal upper bound; the final rounding implication is
% explicitly conditional until the goal global scalar certificate is accepted.
\begin{proposition}[An elementary universal upper bound]
\label{prop:new-self-contained-upper}
With $c=\log(1+\sqrt2)$,
\[
 G_{\mathbb R}\leq\frac{\pi}{2c}
                <\frac{1760}{987}<\frac{357}{200}=1.785.
\]
\end{proposition}
\begin{proof}
For a real unit vector $u$, in the Hilbert direct sum of its odd
real tensor powers define
\[
 \Phi(u)=\bigoplus_{k\geq0}
       \sqrt{\frac{c^{2k+1}}{(2k+1)!}}\,u^{\otimes(2k+1)},\qquad
 \Psi(u)=\bigoplus_{k\geq0}(-1)^k
       \sqrt{\frac{c^{2k+1}}{(2k+1)!}}\,u^{\otimes(2k+1)}.
\]
Since $\sinh c=1$, both vectors have norm one and
\[
 \langle\Phi(u),\Psi(v)\rangle
                    =\sin\bigl(c\langle u,v\rangle\bigr).
\]
Also $c<1$, because $\sinh1>1$, whereas
$\pi=4\int_0^1(1+x^2)^{-1}\,dx>2$; hence $c<\pi/2$.
For any finite families $u_i,v_j$, take a standard Gaussian vector
in the finite-dimensional span of their transformed vectors.  Put
$\epsilon_i=\operatorname{sgn}\langle G,\Phi(u_i)\rangle$ and
$\delta_j=\operatorname{sgn}\langle G,\Psi(v_j)\rangle$.
Two Gaussian linear forms with correlation $\rho$ have sign-product
expectation $2\arcsin(\rho)/\pi$: in their two-dimensional span,
rotational invariance makes the angle uniform and their signs
 differ with probability $\arccos(\rho)/\pi$.  Consequently
\[
 \mathbb E[\epsilon_i\delta_j]
 =\frac2\pi\arcsin\!\left(\sin(c\langle u_i,v_j\rangle)\right)
 =\frac{2c}{\pi}\langle u_i,v_j\rangle.
\]
For every real matrix $A$, its scalar sign norm bounds each rounded
bilinear form, so
\[
 \frac{2c}{\pi}\sum_{i,j}a_{ij}\langle u_i,v_j\rangle
 =\mathbb E\sum_{i,j}a_{ij}\epsilon_i\delta_j\leq C(A).
\]
Taking the vector supremum and then the supremum over matrices
proves $G_{\mathbb R}\leq\pi/(2c)$.

It remains to certify the numerical comparisons using rational
arithmetic.  The exact identity
\[
 \frac{x^4(1-x)^4}{1+x^2}
 =x^6-4x^5+5x^4-4x^2+4-\frac4{1+x^2}
\]
gives
\[
 0<\int_0^1\frac{x^4(1-x)^4}{1+x^2}\,dx
                         =\frac{22}{7}-\pi.
\]
Set $x_0=141/160$.  The positive exponential series and the
successive-tail ratio bound $x_0/(k+1)\leq x_0/14$ for $k\geq13$
give
\begin{align*}
 e^{x_0}
 &<\sum_{k=0}^{12}\frac{x_0^k}{k!}
       +\frac{x_0^{13}}{13!}\frac1{1-x_0/14}\\
 &=\frac{36546709859637513399380063330480378767}
          {15140013877549634840166400000000000000}
  <\frac{239}{99}<1+\sqrt2.
\end{align*}
The middle comparison follows by clearing positive denominators;
the final one follows from
$2-(140/99)^2=2/9801>0$.  Thus $c>141/160$, and
\[
 \frac\pi{2c}<\frac{22/7}{2(141/160)}
                =\frac{1760}{987},\qquad
 \frac{357}{200}-\frac{1760}{987}=\frac{359}{197400}>0.
\]
All polynomial and rational computations are exactly replayed by
\texttt{certify\_elementary\_upper.py}; they require no decimal
approximation to $\pi$, a logarithm, or a square root.
\end{proof}

\begin{proposition}[The three-significant-digit implication]
\label{prop:new-three-significant-digit-implication}
If a separately accepted complete matrix-sequence certificate proves
$G_{\mathbb R}\geq71/40$, then
\[
 \frac{71}{40}=1.775\leq G_{\mathbb R}
                         <\frac{1760}{987}<1.785.
\]
Thus $G_{\mathbb R}$ rounds to $1.78$ to three significant digits.
In particular, a complete scalar cap $N(T)\leq40/71$ for a seed
whose vector numerator tends to one proves this implication.
\end{proposition}
\begin{proof}
The upper endpoint is Proposition~\ref{prop:new-self-contained-upper}.
The interval $[1.775,1.785)$ is precisely the rounding interval for
$1.78$ at the hundredths place.  For the stated seed, its actual
matrix ratio limit is $1/N(T)\geq71/40$, so its matrices give the
required lower bound for the unrestricted constant.
\end{proof}


% Compact analytic matrix entries, independently checked October 2.
\begin{proposition}[Compact evaluation of cubic cell entries]
\label{prop:new-compact-cubic-cell-entry}
Let $h_\alpha(x)=\prod_{j=1}^d
\operatorname{He}_{\alpha_j}(x_j)/\sqrt{\alpha_j!}$ denote the
orthonormal probabilists' Hermite polynomials for standard Gaussian
measure. The degree-three projection kernel has the exact formula
\[
 \mathcal K_3(x,y):=\sum_{|\alpha|=3}h_\alpha(x)h_\alpha(y)
 =\frac{(x\cdot y)^3-3(x\cdot y)(\|x\|^2+\|y\|^2)
                  +3(d+2)(x\cdot y)}6.
\]
For two product cells $C,D$, let
\[
 a_{j,k}=\mathbb E[X_j^k\mid X\in C],\qquad
 b_{j,k}=\mathbb E[Y_j^k\mid Y\in D],\qquad k=1,2,3,
\]
where $X,Y$ are independent standard Gaussian vectors. Put
\[
 u_j=a_{j,1}b_{j,1},\quad
 e_j=(a_{j,2}-1)(b_{j,2}-1),\quad
 t_j=(a_{j,3}-3a_{j,1})(b_{j,3}-3b_{j,1}),\quad S=\sum_j u_j.
\]
Then
\[
 \mathbb E[\mathcal K_3(X,Y)\mid X\in C,Y\in D]
 =\frac16\left\{
 \sum_j t_j+3S\sum_j e_j-3\sum_j e_ju_j
          +S^3-3S\sum_j u_j^2+2\sum_j u_j^3\right\}.
\]
Consequently, one compressed cubic matrix entry requires $O(d)$
arithmetic operations after its one-dimensional conditional moments are
known. No enumeration of the $\binom{d+2}{3}$ degree-three features is
required.
\end{proposition}
\begin{proof}
The degree-three multi-indices have patterns $(3)$, $(2,1)$ and
$(1,1,1)$. Since $\operatorname{He}_1(x)=x$,
$\operatorname{He}_2(x)=x^2-1$ and
$\operatorname{He}_3(x)=x^3-3x$, their conditional contributions are
\[
 \frac16\sum_j t_j
 +\frac12\sum_{i\ne j}e_i u_j
 +\sum_{i<j<k}u_i u_j u_k.
\]
Use $\sum_{i\ne j}e_i u_j=S\sum_i e_i-\sum_i e_i u_i$ and
$6\sum_{i<j<k}u_i u_j u_k
 =S^3-3S\sum_i u_i^2+2\sum_i u_i^3$.
Product conditioning preserves the independence of the coordinates,
so these identities prove the stated expectation. Substituting the
point moments $a_{j,k}=x_j^k$, $b_{j,k}=y_j^k$ and collecting terms
gives the displayed projection kernel.

For an interval $I=(\ell,r)$ with positive Gaussian mass
$p_I=\Phi(r)-\Phi(\ell)$, its required conditional moments are
\[
 m_1(I)=\frac{\varphi(\ell)-\varphi(r)}{p_I},\qquad
 m_2(I)=1+\frac{\ell\varphi(\ell)-r\varphi(r)}{p_I},
\]
\[
 m_3(I)=
 \frac{(\ell^2+2)\varphi(\ell)-(r^2+2)\varphi(r)}{p_I}.
\]
The expressions at infinite endpoints are interpreted by their zero
limits. These formulas follow by integrating
$\varphi'(x)=-x\varphi(x)$.

For a selected half-cell $C$ in the antipodal fold, let
$w_C=2\gamma_d(C)$. The analytic cubic seed entry is exactly
\[
 A_{C,D}=w_Cw_D\left(S-
 b\,\mathbb E[\mathcal K_3(X,Y)\mid C,D]\right).
\]
Indeed, every degree-one or degree-three Hermite cell feature equals
$w_C\mathbb E[h_\alpha(X)\mid C]$, so its outer product is exactly
the folded cell compression. Rational approximation and the reduced
matrix maps are applied subsequently as before; this entry formula
changes neither the sequence nor its certified bounds.
\end{proof}

% Fully rational finite lower bound, independently checked October 2.
\begin{corollary}[A fully rational finite bound]
\label{cor:new-rational-log-cutoff-bound}
In Proposition~\ref{prop:new-log-cutoff-sequence}, assume
$0<b<1$, $d\geq4$ and $N(P_1-bP_3)\leq B$. Define the rational numbers
\[
 \underline q_d=
 \left(1-\frac1{2d}\right)
 \left(1-\frac b{2(d+2)}\right),
\]
\[
 \overline L_d=
 \frac1d\left(\frac{1-b/2}{d-1}
                   +\frac b{2(d+2)}+\frac b{d-2}\right)
 +\frac4{d^2}\left(1+\frac b{2(d+2)}\right)
 +\frac{4b}{d^4}.
\]
Then the same rational reduced matrices satisfy
\[
 G(H_d^*(b))\geq
 \frac{\max(0,\underline q_d-\overline L_d-d^{-2})}{B+d^{-2}}.
\]
If $B$ is rational, this bound uses only rational arithmetic. It does
not require a numerical gamma-function evaluation or a Gaussian tail
evaluation. The right-hand side converges to $1/B$ at rate $O(d^{-1})$.
\end{corollary}
\begin{proof}
For $R=\|X\|$, its centered Hermite expansion has only even degrees,
starting at degree two. The Hermite derivative identity therefore gives
\[
 \operatorname{Var}(R)\leq
 \frac12\sum_{k\geq1}k\|P_kR\|_2^2
 =\frac12\mathbb E\|\nabla R\|^2=\frac12.
\]
Consequently $a_d=(\mathbb ER)^2/d\geq1-1/(2d)$, giving
$q_d(b)\geq\underline q_d$. The logarithmic cutoff gives
$p_d\leq d^{-4}$. Finally,
\[
 \sqrt{\frac{d+4}{d+2}}
 =\sqrt{1+\frac2{d+2}}\leq1+\frac1{d+2}.
\]
Substitute these inequalities into the stated midpoint-error bound,
using $\sqrt{p_d}\leq d^{-2}$ and the positive coefficients of its
tail terms. Thus $\mathcal L_d^*(b)\leq\overline L_d$.
The finite ratio bound of the logarithmic-cutoff proposition now proves
the assertion. Since $\underline q_d=1-(1+b)/(2d)+O(d^{-2})$ and
$\overline L_d=O(d^{-2})$, the stated rate follows.
\end{proof}

% Final exact lower endpoint, emitted after independent complete acceptance.
\section{Final constructive enclosure and finite matrices}\label{sec:new-final-goal}
% Generated only after the bound complete aggregate passes the goal.
% Aggregate SHA-256: f1e6bf6641fa20dd23e1f60cd1c2deb743ac0be98035210ab6d58b24718c6f08
\begin{theorem}[Certified constructive enclosure]
\label{thm:new-final-goal-global-norm}
Let $b=711/1000$ and $B_*=\frac{40}{71}$. Then
\[
 \mathcal N(b):=N(P_1-bP_3)\leq B_*,\qquad
 G_{\mathbb R}\geq \frac{71}{40}=1.775.
\]
For the logarithmic-cutoff rational reduced matrices of
Proposition~\ref{prop:new-log-cutoff-sequence},
\[
 \lim_{d\to\infty}G(H_d^*(b))=\frac1{\mathcal N(b)}\geq\frac1{B_*}.
\]
In particular the matrices are symmetric, have zero diagonal and zero row
sums, and their entrywise absolute values are doubly stochastic.
For every stated finite index their certified lower bound is
\[
 G(H_d^*(b))\geq
 \frac{\max(0,q_d(b)-\mathcal L_d^*(b)-d^{-2})}{B_*+d^{-2}}.
\]
These lower formulas converge to $1/B_*$ at rate $O(d^{-1})$, or
$O(n^{-4})$ on $d=n^4$. This rate does not assert a quantitative rate
for the actual scalar denominator.
Together with Proposition~\ref{prop:new-self-contained-upper},
\[
 1.775\leq G_{\mathbb R}<\frac{1760}{987}<1.785.
\]
Thus $G_{\mathbb R}$ rounds to $1.78$ to three significant digits.
\end{theorem}
\begin{proof}
The complete same-coefficient lower-amplitude proof covers
$0\leq u\leq1381/2000$ with cap
$563307766601/10^{12}$. The accepted upper cover retains all $272579$
physical source parents, $8737$ inherited proved-empty cells, and $19693$
closed amplitude children. The shared cumulant--shape refinement has
$18480$ inherited closed product children. The final step replaces exactly
two of them by their eight closed Cartesian descendants, giving
$18486$ current closed moment children. Every inherited bound remains
available; any empty child is
removed only by an accepted necessary physical constraint or a uniformly
valid source halfspace.

There are 89 independently accepted complete whole-line criteria.
The accepted affine parent has $5280$ independently repriced uniform
moment tables in addition to the complete shared-moment ancestor.
The final four moment boxes have $356$ uniform tables: $344$ are checked
equal to previously accepted tables and the remaining $12$ are freshly
repriced with directed arithmetic. The four corresponding source-plane
boxes are also independently repriced. These constants cover the actual shared physical
moment domain, and every spectral gate and literal remainder allowance
is bound to the corresponding complete local proof. Such criteria remain
valid outside the box in which they were discovered. The full even Gram
criteria retain their actual radial quadratic coefficient, including
negative coefficients.

At each source cell all radial knots are split before a native chord is
used, and each chord retains one full native kernel. Uniform quadratic
prices are replaced by valid affine secants or tangents. The closed
dominance polygons of Proposition~\ref{prop:new-hybrid-affine-source-envelope}
cover all ties and degenerate source sets. Their exact vertex maxima,
with every old cap retained, prove the displayed complete upper bound.
The final eight descendants are all nonempty and pass the exact threshold
$40/71$; there is no uncovered or failing source cell.
Any further closed source subdivision is checked against its parent and
retains the entire parent domain. The lower and upper amplitude proofs
use the same coefficient $b$ and cover $[0,\sqrt{2/\pi}]$.

The complete independent aggregate is
\texttt{FINAL\_GOAL\_AFFINE\_MOMENT\_GLOBAL\_ACCEPTANCE.json}. It binds every local, numerical, geometric,
ancestor and lower-amplitude input. The final envelope has SHA-256
\begin{center}\small\ttfamily
bdcde16a1efcc9a1e414236394ea0edb1f63722e3711bcd1a01c2530e005c8ff
\end{center}
The Gaussian-to-matrix limit and exact reduced-class maps prove the
constructive matrix claims. The unrestricted upper corollary gives the
rounding interval. Neither the scalar cap nor this matrix-family limit
is identified with the exact unrestricted Grothendieck constant.
\end{proof}

\paragraph{Finite matrices evaluated with rational arithmetic.}
Taking $B=40/71$ in Corollary~\ref{cor:new-rational-log-cutoff-bound}
gives the following certified lower bounds. Every displayed decimal is
rounded down to twelve places from an exact rational value:
\begin{center}
\begin{tabular}{rr}
Gaussian dimension $d$ & Lower bound for $G(H_d^*(711/1000))$\\\hline
$10^3$ & $1.773467997012$ \\
$10^4$ & $1.774848013893$ \\
$10^6$ & $1.774998481474$ \\
$10^8$ & $1.774999984814$ \\
\hline
\end{tabular}
\end{center}
The index $d$ is the Gaussian dimension. The matrix order includes the
product-cell compression and the subsequent exact cloning maps; these
large matrices are specified implicitly. All total entrywise rational
approximation errors $d^{-2}$ are charged in the displayed bounds.
The exact arithmetic inputs are recorded in
\texttt{FINAL\_GOAL\_RATIONAL\_FINITE\_BOUND.json}.
These finite lower certificates converge to $71/40$. The actual matrix
ratios converge to $1/\N(711/1000)$, which is at least $71/40$.
The limiting inequality and the unrestricted upper bound together prove
the three-significant-digit conclusion.

The source package also supplies independent Arb768 and MPFR768
evaluations of the exact logarithmic-cutoff formula on $d=n^4$, joined to
the same accepted norm proof by
\texttt{FINAL\_GOAL\_FINITE\_FORMULA\_ACCEPTANCE.json}.
No numerical extrapolation is used to prove the limit.

\paragraph{Replaying the completed proof.}
The source package retains every explicitly bound input and the complete
historical source inventory. The frozen current input whitelist is
\texttt{FINAL\_GOAL\_INPUT\_MANIFEST.json}; its replay guide is
\texttt{README\_GOAL\_REPLAY.md}. The \texttt{--all-local} mode reconstructs
the protected baseline, every current whole-line criterion, the directed
moment and source prices, the exact closed partition and dominance
geometry, and the elementary upper-bound arithmetic. Every criterion
uses its own literal remainder allowances and its certified tail cutoff.
All input hashes are checked before and after replay; accepted certificate
bytes are preserved and new outputs are written separately.

\end{document}
